% !TeX program = xelatex
% !TeX root = main.tex
\documentclass[12pt,a4paper,oneside,openany]{book}

% !TeX root = ../main.tex
% ----------------------------------------------------------------------
% Paquetes, colores, geometría y macros matemáticas
% ----------------------------------------------------------------------

% --- Fuentes y codificación (válido para XeLaTeX, LuaLaTeX y pdfLaTeX) ---
\usepackage{iftex}
\ifPDFTeX
  \usepackage[T1]{fontenc}
  \usepackage[utf8]{inputenc}
  \usepackage{lmodern}
\else
  \usepackage{fontspec}
\fi
\usepackage[spanish,es-nodecimaldot]{babel}
\spanishdeactivate{<>}
\AtBeginDocument{\shorthandoff{<>}}
\usepackage{microtype}

% --- Matemática ---
\usepackage{amsmath, amssymb, amsthm, mathtools}
\usepackage{mathrsfs}
\usepackage{cancel}
\usepackage{interval}

% --- Utilidades ---
\usepackage{graphicx}
\usepackage{enumitem}
\usepackage[table]{xcolor}
\usepackage{tikz}
\usepackage{pgfplots}
\pgfplotsset{compat=1.18}

% --- Colores (los mismos que en la portada del final de Temas) ---
\definecolor{darkblue}{RGB}{0, 51, 102}
\definecolor{accentblue}{RGB}{0, 102, 204}
\definecolor{mediumgray}{RGB}{100, 100, 100}
\definecolor{lightgray}{RGB}{240, 240, 240}

% --- Geometría y párrafos ---
\usepackage{geometry}
\geometry{
  a4paper,
  left=2.8cm, right=2.8cm,
  top=3.2cm, bottom=3.2cm,
  headheight=15pt
}
\setlength{\parindent}{0pt}
\setlength{\parskip}{0.6em plus 0.3em minus 0.2em}
\setlist{itemsep=0.2em, topsep=0.3em, parsep=0pt}

% --- Conjuntos numéricos y atajos ---
\newcommand{\N}{\mathbb{N}}
\newcommand{\Z}{\mathbb{Z}}
\newcommand{\Q}{\mathbb{Q}}
\newcommand{\R}{\mathbb{R}}
\newcommand{\C}{\mathbb{C}}
\newcommand{\K}{\mathbb{K}}
\newcommand{\bb}[1]{\mathbb{#1}}
\newcommand{\e}{\varepsilon}
\renewcommand{\O}{\varnothing}

% --- Operadores ---
\DeclareMathOperator{\card}{card}
\DeclareMathOperator{\id}{id}
\DeclareMathOperator{\Img}{Im}
\DeclareMathOperator{\diam}{diam}
\DeclareMathOperator{\dist}{dist}
\DeclareMathOperator{\interior}{int}
\newcommand{\cl}[1]{\overline{#1}}
\newcommand{\abs}[1]{\left\lvert #1 \right\rvert}
\newcommand{\norm}[1]{\left\lVert #1 \right\rVert}
\newcommand{\dd}{\mathop{}\!\mathrm{d}}

% --- Integrales superior e inferior ---
\def\upint{\mathchoice%
  {\mkern13mu\overline{\vphantom{\intop}\mkern7mu}\mkern-20mu}%
  {\mkern7mu\overline{\vphantom{\intop}\mkern7mu}\mkern-14mu}%
  {\mkern7mu\overline{\vphantom{\intop}\mkern7mu}\mkern-14mu}%
  {\mkern7mu\overline{\vphantom{\intop}\mkern7mu}\mkern-14mu}%
  \int}
\def\lowint{\mkern3mu\underline{\vphantom{\intop}\mkern7mu}\mkern-10mu\int}

% !TeX root = ../main.tex
% ----------------------------------------------------------------------
% Entornos: definiciones, teoremas, ejemplos, observaciones y demostraciones
% ----------------------------------------------------------------------
\usepackage{thmtools}
\usepackage[framemethod=default]{mdframed}

\mdfsetup{skipabove=0.9em, skipbelow=0.6em}

% Caja con barra lateral de color
\mdfdefinestyle{barra}{
  topline=false, rightline=false, bottomline=false,
  linewidth=2.5pt,
  innerleftmargin=10pt, innerrightmargin=8pt,
  innertopmargin=10pt, innerbottommargin=10pt,
  leftmargin=0pt, rightmargin=0pt,
  splittopskip=0.5em, splitbottomskip=0.5em
}

% Teoremas, proposiciones, lemas y corolarios
\declaretheoremstyle[
  headfont=\bfseries\color{darkblue},
  notefont=\bfseries\color{accentblue},
  notebraces={(}{)},
  bodyfont=\normalfont,
  headpunct={.},
  postheadspace=0.5em,
  mdframed={style=barra, linecolor=darkblue, backgroundcolor=darkblue!5}
]{caja-teorema}

% Definiciones
\declaretheoremstyle[
  headfont=\bfseries\color{accentblue},
  notefont=\bfseries\color{accentblue},
  notebraces={(}{)},
  bodyfont=\normalfont,
  headpunct={.},
  postheadspace=0.5em,
  mdframed={style=barra, linecolor=accentblue, backgroundcolor=accentblue!5}
]{caja-definicion}

% Ejemplos
\declaretheoremstyle[
  headfont=\bfseries\color{mediumgray},
  notefont=\bfseries\color{mediumgray},
  notebraces={(}{)},
  bodyfont=\normalfont,
  headpunct={.},
  postheadspace=0.5em,
  mdframed={style=barra, linecolor=mediumgray!60, backgroundcolor=white}
]{caja-ejemplo}

% Observaciones
\declaretheoremstyle[
  headfont=\bfseries\color{mediumgray},
  notefont=\bfseries\color{mediumgray},
  notebraces={(}{)},
  bodyfont=\normalfont,
  headpunct={.},
  postheadspace=0.5em,
  mdframed={style=barra, linecolor=mediumgray!60, backgroundcolor=lightgray}
]{caja-observacion}

\declaretheorem[numberwithin=chapter, style=caja-definicion, name=Definición]{definition}
\declaretheorem[sibling=definition, style=caja-teorema, name=Proposición]{prop}
\declaretheorem[sibling=definition, style=caja-teorema, name=Teorema]{theorem}
\declaretheorem[sibling=definition, style=caja-teorema, name=Lema]{lemma}
\declaretheorem[sibling=definition, style=caja-teorema, name=Corolario]{corollary}
\declaretheorem[sibling=definition, style=caja-ejemplo, numbered=no, name=Ejemplo]{eg}
\declaretheorem[style=caja-observacion, numbered=no, name=Observación]{note}
\declaretheorem[style=caja-observacion, numbered=no, name=Ejercicio]{ejercicio}

% Demostraciones (fuera de las cajas)
\renewcommand{\proofname}{Demostración}
\renewcommand{\qedsymbol}{\(\blacksquare\)}
\makeatletter
\renewenvironment{proof}[1][\proofname]{%
  \par\pushQED{\qed}\normalfont\topsep=0.3em\relax
  \trivlist\item[\hskip\labelsep\itshape\textcolor{darkblue}{#1}\@addpunct{.}]\ignorespaces
}{%
  \popQED\endtrivlist\@endpefalse
}
\makeatother

% !TeX root = ../main.tex
% ----------------------------------------------------------------------
% Encabezados, títulos, índice y enlaces
% ----------------------------------------------------------------------
\usepackage{epigraph}
\usepackage{fancyhdr}
\usepackage{titlesec}
\usepackage{titletoc}
\usepackage{emptypage}

\addto\captionsspanish{%
  \renewcommand{\contentsname}{Contenido}%
  \renewcommand{\bibname}{Bibliografía}%
}

% --- Epígrafe del prefacio ---
\renewcommand{\epigraphflush}{flushright}
\renewcommand{\epigraphsize}{\normalsize\itshape}
\setlength{\epigraphwidth}{0.7\textwidth}
\renewcommand{\textflush}{flushepinormal}
\renewcommand{\sourceflush}{flushepinormal}

% --- Encabezado y pie de página ---
\newcommand{\fechaclase}{}
\newcommand{\clase}[2]{\renewcommand{\fechaclase}{#1}\chapter{#2}}
\newcommand{\capitulosinnumero}[1]{%
  \chapter*{#1}\addcontentsline{toc}{chapter}{#1}\markboth{#1}{}%
}

\pagestyle{fancy}
\fancyhf{}
\fancyhead[L]{\small\textcolor{darkblue}{\textbf{Notas del teórico}}}
\fancyhead[R]{\small\textcolor{mediumgray}{\nouppercase{\leftmark}}}
\fancyfoot[C]{\thepage}
\renewcommand{\headrulewidth}{0.5pt}
\renewcommand{\headrule}{\color{darkblue}\hrule width\headwidth height\headrulewidth \kern-\headrulewidth}
\renewcommand{\footrulewidth}{0pt}
\renewcommand{\chaptermark}[1]{\markboth{Clase \thechapter\ \textbullet\ #1}{}}

\fancypagestyle{plain}{%
  \fancyhf{}%
  \fancyfoot[C]{\thepage}%
  \renewcommand{\headrulewidth}{0pt}%
}

% --- Títulos ---
\titleformat{\chapter}[display]
  {\normalfont}
  {\large\bfseries\color{accentblue}Clase \thechapter%
    \normalfont\normalsize\color{mediumgray}\quad\textbullet\quad\fechaclase}
  {4pt}
  {\Huge\bfseries\color{darkblue}}
  [\vspace{0.4em}{\color{darkblue}\titlerule[1.5pt]}]
\titleformat{name=\chapter, numberless}[display]
  {\normalfont}
  {}
  {0pt}
  {\Huge\bfseries\color{darkblue}}
  [\vspace{0.4em}{\color{darkblue}\titlerule[1.5pt]}]
\titlespacing*{\chapter}{0pt}{-1.5em}{2em}

\titleformat{\section}
  {\Large\bfseries\color{darkblue}}
  {\thesection.}{0.5em}{}
  [\vspace{0.15em}{\color{darkblue}\titlerule[0.6pt]}]
\titleformat{\subsection}
  {\large\bfseries\color{accentblue}}
  {\thesubsection.}{0.5em}{}
\titleformat{\subsubsection}
  {\normalsize\bfseries\color{accentblue}}
  {\thesubsubsection.}{0.5em}{}
\titlespacing*{\section}{0pt}{2em plus 0.5em minus 0.2em}{0.8em}
\titlespacing*{\subsection}{0pt}{1.5em plus 0.4em minus 0.2em}{0.4em}

% --- Índice general ---
\setcounter{tocdepth}{1}
\titlecontents{chapter}[5.2em]
  {\addvspace{0.9em}\bfseries\color{darkblue}}
  {\contentslabel[Clase~\thecontentslabel]{5.2em}}
  {\hspace*{-5.2em}}
  {\hfill\contentspage}
\titlecontents{section}[7.7em]
  {\small}
  {\contentslabel{2.5em}}
  {\hspace*{-2.5em}}
  {\titlerule*[0.6pc]{.}\contentspage}

% --- Enlaces ---
\usepackage{hyperref}
\hypersetup{
  colorlinks=true,
  linkcolor=darkblue,
  urlcolor=darkblue,
  citecolor=darkblue,
  bookmarksnumbered=true,
  breaklinks=true,
  pdftitle={Notas del teórico - Complementos de Análisis Matemático},
  pdfauthor={Jordi Bustos}
}

% --- Anclas únicas para los entornos que comparten contador ---
\def\theHdefinition{\thechapter.\arabic{definition}}
\def\theHprop{\thechapter.\arabic{definition}}
\def\theHtheorem{\thechapter.\arabic{definition}}
\def\theHlemma{\thechapter.\arabic{definition}}
\def\theHcorollary{\thechapter.\arabic{definition}}



\begin{document}

% ----------------------------------------------------------------------
% Portada (mismo diseño que la del final de Temas)
% ----------------------------------------------------------------------
\begin{titlepage}
  \thispagestyle{empty}
  \begin{center}
    \includegraphics[width=4cm]{img/logo_unlp.png}

    \vspace{0.5cm}
    \textcolor{darkblue}{\rule{\linewidth}{2pt}}

    \vspace{0.5cm}
    {\large \textbf{UNIVERSIDAD NACIONAL DE LA PLATA}}\\
    {\normalsize Facultad de Ciencias Exactas}\\
    {\normalsize Departamento de Matemática}

    \vspace{0.5cm}
    \textcolor{darkblue}{\rule{\linewidth}{2pt}}

    \vspace{2.5cm}
    {\Huge \textbf{\textcolor{darkblue}{NOTAS DEL TEÓRICO}}}

    \vspace{3cm}
    {\large \textbf{Complementos de Análisis Matemático}}\\
    {\normalsize Profesora: Irene Drelichman}

    \vspace{2.5cm}
    {\normalsize
      \renewcommand{\arraystretch}{1.3}
      \begin{tabular}{>{\bfseries\color{darkblue}}l l}
        Autor:   & Jordi Bustos                  \\
        Fecha:   & \today                        \\
        Carrera: & Licenciatura en Matemática    \\
      \end{tabular}}

    \vfill
    {\footnotesize \textcolor{mediumgray}{%
        Notas del curso dictado en el segundo cuatrimestre de 2024.}}

    \vspace{1cm}
  \end{center}
\end{titlepage}

\frontmatter
\tableofcontents

\capitulosinnumero{Prefacio}
\epigraph{``Considero a cada hombre como un deudor \\ de su profesión, \\ y ya que de ella recibe sustento y provecho, \\ así debe procurar, \\ mediante el estudio, \\ servirle de ayuda y ornato.''}{Francis Bacon}

Este libro recoge las notas tomadas durante el curso de Complementos de Análisis Matemático, dictado por Irene Drelichman en el segundo cuatrimestre de 2024.

Las clases 1 a 8, junto con la primera parte de la clase 9, abarcan la primera mitad del curso. Las clases 9 a 22 cubren la segunda mitad, que incluye temas de topología, continuidad, diferenciación, integración y sucesiones de funciones. Se tomaron dos exámenes parciales, uno por cada mitad del curso; sus enunciados se reúnen al final de este libro.

Las clases 9 a 11 (hasta la mitad) cubren los conceptos básicos de topología. Desde la segunda mitad de la clase 11 hasta la 14 se aborda la continuidad. Los teoremas de la función inversa e implícita se tratan en las clases 15 y 16. La integración se desarrolla en las clases 17 a 19. Las clases 20 a 22 se dedican a las sucesiones de funciones y, finalmente, se introducen las ecuaciones diferenciales.

Estas notas se basan principalmente en los libros \textit{Principles of Mathematical Analysis} de Walter Rudin, \textit{Curso de análise, vol.~1} de Elon Lages Lima, \textit{Cálculo y Análisis} de Larotonda (para la demostración de los teoremas de la función inversa e implícita) y \textit{Calculus} de Tom M. Apostol.


\mainmatter
\clase{23/08/2024}{Funciones, naturales y cuerpos}
\section{Repaso de funciones}

Una función $f\colon A \to B$ consta de tres partes: un conjunto $A$ (dominio), un conjunto $B$ (codominio) y una regla que asocia a cada elemento $x \in A$ un único elemento $f(x) \in B$.

El \emph{gráfico} de $f$ es el subconjunto de $A \times B$ dado por
\[
  G(f) = \{ (x, y) \in A \times B : y = f(x) \}.
\]

\begin{definition}[Inyectividad]
  $f\colon A \to B$ es \emph{inyectiva} si para todo $x, y \in A$ vale $f(x) = f(y) \Rightarrow x = y$.
\end{definition}

\begin{definition}[Suryectividad]
  $f\colon A \to B$ es \emph{suryectiva} si para todo $y \in B$ existe $x \in A$ tal que $f(x) = y$.
\end{definition}

\begin{definition}[Biyectividad]
  $f\colon A \to B$ es \emph{biyectiva} si es inyectiva y suryectiva.
\end{definition}

\begin{definition}[Imagen]
  Dados $f\colon A \to B$ y $X \subset A$, la \emph{imagen} de $X$ por $f$ es $f(X) = \{ f(x) : x \in X \}$.
\end{definition}

\subsection{Propiedades de la imagen}

\begin{prop}
  Sean $X, Y \subset A$. Entonces:
  \begin{enumerate}[label=(\arabic*)]
    \item $f(X \cup Y) = f(X) \cup f(Y)$.
    \item $X \subset Y \Rightarrow f(X) \subset f(Y)$.
    \item $f(\varnothing) = \varnothing$.
  \end{enumerate}
\end{prop}

\begin{prop}
  Si $X, Y \subset A$, entonces $f(X \cap Y) \subset f(X) \cap f(Y)$. La igualdad vale para todo par $X, Y$ si y solo si $f$ es inyectiva.
\end{prop}

\begin{proof}
  Si $b \in f(X \cap Y)$, existe $x \in X \cap Y$ con $f(x) = b$. Como $x \in X$ e $x \in Y$, resulta $b \in f(X) \cap f(Y)$.

  Si $f$ es inyectiva y $b \in f(X) \cap f(Y)$, existen $a \in X$ y $a' \in Y$ con $f(a) = f(a') = b$. Por inyectividad $a = a'$, luego $a \in X \cap Y$ y $b \in f(X \cap Y)$.

  Si $f$ no es inyectiva, existen $x \neq y$ con $f(x) = f(y)$. Tomando $X = \{x\}$ e $Y = \{y\}$ se tiene $X \cap Y = \varnothing$, de modo que $f(X \cap Y) = \varnothing$, pero $f(X) \cap f(Y) = \{f(x)\} \neq \varnothing$.
\end{proof}

\subsection{Preimagen}

\begin{definition}[Preimagen]
  Dados $f\colon A \to B$ e $Y \subset B$, la \emph{preimagen} de $Y$ por $f$ es
  \[
    f^{-1}(Y) = \{ x \in A : f(x) \in Y \}.
  \]
\end{definition}

\begin{prop}
  Sean $X, Y \subset B$. Entonces:
  \begin{enumerate}[label=(\arabic*)]
    \item $f^{-1}(X \cup Y) = f^{-1}(X) \cup f^{-1}(Y)$.
    \item $f^{-1}(X \cap Y) = f^{-1}(X) \cap f^{-1}(Y)$.
    \item $f^{-1}(Y^c) = {\bigl(f^{-1}(Y)\bigr)}^c$.
    \item $X \subset Y \Rightarrow f^{-1}(X) \subset f^{-1}(Y)$.
    \item $f^{-1}(B) = A$ y $f^{-1}(\varnothing) = \varnothing$.
  \end{enumerate}
\end{prop}

\begin{proof}
  Todas se deducen directamente de la definición. Por ejemplo, para (3):
  \[
    x \in f^{-1}(Y^c) \iff f(x) \notin Y \iff x \notin f^{-1}(Y) \iff x \in {\bigl(f^{-1}(Y)\bigr)}^c. \qedhere
  \]
\end{proof}

\subsection{Composición e inversas}

Sean $f\colon A \to B$ y $g\colon B \to C$. La \emph{composición} $g \circ f\colon A \to C$ se define como $(g \circ f)(x) = g(f(x))$ para todo $x \in A$. Basta con que $f(A)$ esté contenido en el dominio de $g$.

\begin{prop}
  La composición de funciones inyectivas es inyectiva y la composición de funciones suryectivas es suryectiva. Además, ${(g \circ f)}^{-1}(Z) = f^{-1}\bigl(g^{-1}(Z)\bigr)$ para todo $Z \subset C$.
\end{prop}

\begin{definition}[Restricción]
  La \emph{restricción} de $f\colon A \to B$ a un subconjunto $X \subset A$ es la función $f|_X\colon X \to B$ dada por $f|_X(x) = f(x)$.
\end{definition}

\begin{definition}[Inversas laterales]
  Sea $f\colon A \to B$. Una función $g\colon B \to A$ es:
  \begin{itemize}
    \item una \emph{inversa a izquierda} de $f$ si $g \circ f = \id_A$;
    \item una \emph{inversa a derecha} de $f$ si $f \circ g = \id_B$;
    \item una \emph{inversa} de $f$ si es inversa a izquierda y a derecha.
  \end{itemize}
\end{definition}

\begin{prop}
  Sea $f\colon A \to B$ con $A \neq \varnothing$. Entonces:
  \begin{enumerate}[label=(\arabic*)]
    \item $f$ tiene inversa a izquierda si y solo si $f$ es inyectiva;
    \item $f$ tiene inversa a derecha si y solo si $f$ es suryectiva (esta implicación usa el axioma de elección);
    \item $f$ es biyectiva si y solo si tiene una inversa, que en tal caso es única.
  \end{enumerate}
\end{prop}

\subsection{Familias}

Sea $L$ un conjunto de \emph{índices}, cuyos elementos se representan genéricamente con $\lambda$. Dado un conjunto $X$, una \emph{familia} de elementos de $X$ con índices en $L$ es una función $x\colon L \to X$. El valor de $x$ en $\lambda \in L$ se nota $x_\lambda$ y la familia, ${(x_\lambda)}_{\lambda \in L}$.

\begin{eg}
  Si $L = \{1, 2\}$, una familia $x\colon \{1, 2\} \to X$ se representa por los valores $x_1$ y $x_2$, es decir, se identifica con un par ordenado $(x_1, x_2)$ de elementos de $X$.
\end{eg}

Una familia indexada por $\N$ se llama \emph{sucesión}: ${(x_n)}_{n \in \N}$ es una función $x\colon \N \to X$ con $x_n = x(n)$.

Las propiedades anteriores se extienden a familias arbitrarias de conjuntos; por ejemplo, $f\bigl(\bigcup_\lambda X_\lambda\bigr) = \bigcup_\lambda f(X_\lambda)$ y $f^{-1}\bigl(\bigcap_\lambda Y_\lambda\bigr) = \bigcap_\lambda f^{-1}(Y_\lambda)$.

\section{Números naturales}

Partimos de un conjunto $\N$ y una función $S\colon \N \to \N$ (el \emph{sucesor}) que cumple los \emph{axiomas de Peano}:
\begin{enumerate}[label=(\arabic*)]
  \item $S$ es inyectiva.
  \item $\N \setminus S(\N)$ tiene un único elemento, que llamamos $1$.
  \item (\emph{Principio de inducción}) Si $X \subset \N$ cumple $1 \in X$ y $S(m) \in X$ para todo $m \in X$, entonces $X = \N$.
\end{enumerate}

El principio de inducción permite definir operaciones recursivamente. La \emph{suma} se define por
\[
  m + 1 = S(m), \qquad m + S(n) = S(m + n).
\]

\begin{prop}[Asociatividad]
  Para todos $m, n, p \in \N$, $m + (n + p) = (m + n) + p$.
\end{prop}

\begin{proof}
  Sea $X = \{ p \in \N : m + (n + p) = (m + n) + p \text{ para todos } m, n \in \N \}$.

  Para $p = 1$: $m + (n + 1) = m + S(n) = S(m + n) = (m + n) + 1$, así que $1 \in X$.

  Si $p \in X$, entonces
  \[
    m + (n + S(p)) = m + S(n + p) = S\bigl(m + (n + p)\bigr) = S\bigl((m + n) + p\bigr) = (m + n) + S(p),
  \]
  de modo que $S(p) \in X$. Por inducción, $X = \N$.
\end{proof}

\begin{prop}
  La suma es conmutativa: $m + n = n + m$. Además vale la ley de cancelación: $m + n = m + p \Rightarrow n = p$.
\end{prop}

\begin{prop}[Tricotomía]
  Dados $m, n \in \N$, se cumple exactamente una de las siguientes condiciones:
  \[
    m = n, \qquad \exists\, p \in \N \colon m = n + p, \qquad \exists\, p \in \N \colon n = m + p.
  \]
  Se escribe $n < m$ en el segundo caso y $m < n$ en el tercero.
\end{prop}

\begin{definition}[Producto]
  La \emph{multiplicación} se define recursivamente por $m \cdot 1 = m$ y $m \cdot S(n) = m \cdot n + m$. Es asociativa, conmutativa, distributiva respecto de la suma, cumple la ley de cancelación y es compatible con el orden.
\end{definition}

\begin{theorem}[Principio de buena ordenación]
  Todo subconjunto no vacío $A \subset \N$ tiene un elemento mínimo.
\end{theorem}

\begin{proof}
  Para cada $m \in \N$ sea $I_m = \{ p \in \N : p \leq m \}$ y sea
  \[
    X = \{ m \in \N : I_m \subset \N \setminus A \}.
  \]
  Si $1 \in A$, entonces $1$ es el mínimo de $A$. Supongamos $1 \notin A$; entonces $1 \in X$. Además $X \neq \N$, porque si fuera $X = \N$ tendríamos $m \in I_m \subset \N \setminus A$ para todo $m$, es decir $A = \varnothing$.

  Por el principio de inducción, existe $m \in X$ tal que $m + 1 \notin X$. Entonces todos los elementos menores o iguales que $m$ están en $\N \setminus A$, pero $I_{m+1} \not\subset \N \setminus A$, de donde $m + 1 \in A$. Luego $m + 1$ es el mínimo de $A$.
\end{proof}

\section{Cuerpos}

Un \emph{cuerpo} es un conjunto $\K$ dotado de dos operaciones, suma y producto, que cumplen los siguientes axiomas para todos $x, y, z \in \K$.

\subsection{Axiomas de la suma}

\begin{enumerate}[label=(S\arabic*)]
  \item Asociatividad: $(x + y) + z = x + (y + z)$.
  \item Conmutatividad: $x + y = y + x$.
  \item Neutro: existe $0 \in \K$ tal que $x + 0 = x$.
  \item Opuesto: existe $-x \in \K$ tal que $x + (-x) = 0$.
\end{enumerate}

\subsection{Axiomas del producto}

\begin{enumerate}[label=(P\arabic*)]
  \item Asociatividad: $(xy)z = x(yz)$.
  \item Conmutatividad: $xy = yx$.
  \item Neutro: existe $1 \in \K$, $1 \neq 0$, tal que $x \cdot 1 = x$.
  \item Inverso: si $x \neq 0$, existe $x^{-1} \in \K$ tal que $x \cdot x^{-1} = 1$.
\end{enumerate}

Además se pide la \emph{distributividad}: $x(y + z) = xy + xz$.

\begin{eg}
  $\Q$, $\R$, $\C$ y $\Z_2 = \Z/2\Z$ son cuerpos.
\end{eg}

\subsection{Cuerpos ordenados}

Un \emph{cuerpo ordenado} es un cuerpo $\K$ con un subconjunto $P \subset \K$, llamado conjunto de elementos \emph{positivos}, tal que:
\begin{enumerate}[label=(O\arabic*)]
  \item si $x, y \in P$, entonces $x + y \in P$ y $xy \in P$;
  \item para todo $x \in \K$ se cumple exactamente una de las siguientes: $x \in P$, $-x \in P$ o $x = 0$.
\end{enumerate}

\begin{eg}
  $\Q$ con $P = \{ p/q : p, q \in \N \}$ es un cuerpo ordenado. En cambio, $\Z_2$ no admite un orden: si $1 \in P$, entonces $1 + 1 = 0 \in P$, lo que contradice (O2).
\end{eg}

\begin{prop}
  En un cuerpo ordenado, si $a \neq 0$ entonces $a^2 \in P$.
\end{prop}

\begin{proof}
  Si $a \in P$, entonces $a^2 = a \cdot a \in P$ por (O1). Si $a \notin P$, por (O2) es $-a \in P$ y entonces $a^2 = (-a)(-a) \in P$.
\end{proof}

En un cuerpo ordenado se define $x < y$ cuando $y - x \in P$. Esta relación cumple:

\begin{prop}
  Sean $x, y, z \in \K$. Entonces:
  \begin{enumerate}[label=(\arabic*)]
    \item (Transitividad) $x < y$ e $y < z$ implican $x < z$.
    \item (Tricotomía) vale exactamente una de las relaciones $x = y$, $x < y$, $y < x$.
    \item (Monotonía de la suma) $x < y \Rightarrow x + z < y + z$.
    \item (Monotonía del producto) $x < y$ y $0 < z \Rightarrow xz < yz$.
  \end{enumerate}
\end{prop}

Se escribe $x \leq y$ si $x < y$ o $x = y$, es decir, si $y - x \in P \cup \{0\}$. Con esta relación se cumplen todas las propiedades anteriores y además la antisimetría: $x \leq y$ e $y \leq x$ implican $x = y$.

También se tiene la noción de intervalo. Dados $a, b \in \K$ se definen
\[
  [a, b] = \{ x \in \K : a \leq x \leq b \}, \quad (a, b) = \{ x \in \K : a < x < b \},
\]
y análogamente $[a, b)$, $(a, b]$, $[a, +\infty)$, $(a, +\infty)$, $(-\infty, b]$ y $(-\infty, b)$. Notar que $[a, a] = \{a\}$.

Un subconjunto $X \subset \K$ es \emph{acotado superiormente} si existe $b \in \K$ tal que $x \leq b$ para todo $x \in X$ (se dice que $b$ es una \emph{cota superior}). Análogamente se define \emph{acotado inferiormente} y \emph{cota inferior}. Un conjunto es \emph{acotado} si lo es superior e inferiormente.


\clase{27/08/2024}{Cuerpos completos y cardinalidad}
\section{Cuerpos arquimedianos}

Si $\K$ es un cuerpo ordenado, se tiene $\N \subset \Z \subset \Q \subset \K$, pero esto no implica que $\N$ sea no acotado en $\K$.

\begin{eg}
  Sea $\Q(t)$ el cuerpo de funciones racionales $r(t) = p(t)/q(t)$, con $p, q \in \Q[t]$ y $q \neq 0$. Se lo ordena declarando que $r(t)$ es positivo si y solo si el coeficiente principal del polinomio $p \cdot q$ es positivo.

  Para cada $n \in \N$, el polinomio $t - n$ tiene coeficiente principal $1 > 0$, de modo que $t - n \in P$, es decir, $t > n$. Por lo tanto, en $\Q(t)$ el conjunto $\N$ está acotado superiormente, por ejemplo por $t$.
\end{eg}

\begin{definition}[Cuerpo arquimediano]
  Un cuerpo ordenado $\K$ es \emph{arquimediano} si $\N \subset \K$ no es acotado superiormente.
\end{definition}

\begin{theorem}
  Sea $\K$ un cuerpo ordenado. Son equivalentes:
  \begin{enumerate}[label=(\arabic*)]
    \item $\K$ es arquimediano, es decir, $\N$ no es acotado superiormente.
    \item Dados $a, b \in \K$ con $a > 0$, existe $n \in \N$ tal que $n a > b$.
    \item Dado $a \in \K$ con $a > 0$, existe $n \in \N$ tal que $0 < \frac{1}{n} < a$.
  \end{enumerate}
\end{theorem}

\begin{proof}
  \begin{description}[font=\normalfont, style=nextline, leftmargin=2em, itemsep=0.3em]
    \item[(1) $\Rightarrow$ (2)] Como $\N$ no es acotado, $b/a$ no es cota superior de $\N$, así que existe $n \in \N$ con $n > b/a$. Luego $na > b$.
    \item[(2) $\Rightarrow$ (3)] Dado $a > 0$, aplicamos (2) con $b = 1$: existe $n$ con $na > 1$, es decir, $0 < \frac{1}{n} < a$.
    \item[(3) $\Rightarrow$ (1)] Sea $b \in \K$. Si $b \leq 0$, claramente $b$ no es cota superior de $\N$. Si $b > 0$, aplicamos (3) con $a = 1/b$: existe $n$ con $0 < \frac{1}{n} < \frac{1}{b}$, o sea $n > b$. Luego ningún elemento de $\K$ es cota superior de $\N$.\qedhere
  \end{description}
\end{proof}

\section{Supremo e ínfimo}

\begin{definition}[Supremo]
  Sean $\K$ un cuerpo ordenado y $X \subset \K$ acotado superiormente. Un elemento $b \in \K$ es el \emph{supremo} de $X$ si es la menor de las cotas superiores de $X$ en $\K$. Es decir:
  \begin{enumerate}[label=(\arabic*)]
    \item $x \leq b$ para todo $x \in X$;
    \item si $c \in \K$ y $x \leq c$ para todo $x \in X$, entonces $b \leq c$.
  \end{enumerate}
  Equivalentemente, (2) puede reemplazarse por: para todo $c < b$ existe $x \in X$ tal que $c < x$.
\end{definition}

\begin{note}
  \begin{enumerate}[label=(\arabic*)]
    \item El supremo de un conjunto, si existe, es único.
    \item Si un conjunto tiene máximo, este es su supremo.
    \item Si $X = \varnothing$, todo elemento de $\K$ es cota superior. Como $\K$ no tiene mínimo, $\varnothing$ no tiene supremo en $\K$.
  \end{enumerate}
\end{note}

\begin{definition}[Ínfimo]
  Sean $\K$ un cuerpo ordenado y $X \subset \K$ acotado inferiormente. Un elemento $b \in \K$ es el \emph{ínfimo} de $X$ si es la mayor de las cotas inferiores de $X$ en $\K$. Es decir:
  \begin{enumerate}[label=(\arabic*)]
    \item $x \geq b$ para todo $x \in X$;
    \item si $c \in \K$ y $x \geq c$ para todo $x \in X$, entonces $b \geq c$.
  \end{enumerate}
  Equivalentemente, (2) puede reemplazarse por: para todo $c > b$ existe $x \in X$ tal que $x < c$.
\end{definition}

\begin{eg}
  Dados $a < b$ en $\K$, si $X = (a, b)$ entonces $\inf X = a$ y $\sup X = b$.

  Por definición, $a$ es cota inferior y $b$ es cota superior. Veamos que $a$ es la mayor cota inferior. Sea $c > a$. Si $c \geq b$, entonces $c$ no es cota inferior pues cualquier $x \in X$ cumple $x < b \leq c$. Si $c < b$, entonces $\frac{a + c}{2} \in X$ y $\frac{a + c}{2} < c$, de modo que $c$ tampoco es cota inferior. Luego $a = \inf X$. El argumento para el supremo es análogo.
\end{eg}

\begin{eg}
  Sea $Y = \left\{ \frac{1}{2^n} : n \in \N \right\} \subset \Q$. Entonces $\sup Y = \frac{1}{2}$ y $\inf Y = 0$.

  Como $\frac{1}{2} \in Y$ y $\frac{1}{2^n} \leq \frac{1}{2}$ para todo $n$, se tiene $\sup Y = \frac{1}{2}$.

  Además $0 < \frac{1}{2^n}$ para todo $n$, así que $0$ es cota inferior. Sea $c > 0$. Por la desigualdad de Bernoulli, $2^n = {(1 + 1)}^n \geq 1 + n > n$. Como $\Q$ es arquimediano, existe $n$ con $n > \frac{1}{c}$, y entonces $2^n > \frac{1}{c}$, es decir, $\frac{1}{2^n} < c$. Luego $c$ no es cota inferior y $\inf Y = 0$.
\end{eg}

El problema más serio de los racionales desde el punto de vista del análisis es que algunos conjuntos acotados de números racionales no tienen supremo (o ínfimo) en $\Q$.

\begin{eg}
  Sean
  \[
    X = \{ x \in \Q : x \geq 0,\ x^2 < 2 \}, \qquad Y = \{ y \in \Q : y > 0,\ y^2 > 2 \}.
  \]
  Si $z > 2$, entonces $z^2 > 4$ y $z \notin X$; luego $X \subset [0, 2]$ está acotado. Además $Y \subset (0, +\infty)$ está acotado inferiormente. Veamos que $X$ no tiene supremo en $\Q$.

  \begin{enumerate}[label=(\arabic*)]
    \item \emph{$X$ no tiene máximo.} Dado $x \in X$ se elige $r \in \Q$ con $0 < r < \min\left\{1, \frac{2 - x^2}{2x + 1}\right\}$. Como $r < 1$, resulta $r^2 < r$ y
          \[
            {(x + r)}^2 = x^2 + 2xr + r^2 < x^2 + r(2x + 1) < 2,
          \]
          de modo que $x + r \in X$ y $x + r > x$.
    \item \emph{$Y$ no tiene mínimo.} Dado $y \in Y$ se elige $r \in \Q$ con $0 < r < \frac{y^2 - 2}{2y}$. Notar que $\frac{y^2-2}{2y} < \frac{y}{2}$, de modo que $y - r > 0$. Además
          \[
            {(y - r)}^2 = y^2 - 2yr + r^2 > y^2 - 2yr > y^2 - (y^2 - 2) = 2,
          \]
          así que $y - r \in Y$ e $y - r < y$.
    \item \emph{Si $x \in X$ e $y \in Y$, entonces $x < y$.} En efecto, si fuese $y \leq x$, sería $y^2 \leq x^2 < 2$, lo que contradice $y^2 > 2$.
  \end{enumerate}

  Supongamos que existe $\alpha = \sup X \in \Q$. Como $1 \in X$, es $\alpha \geq 1 > 0$.
  \begin{itemize}
    \item Si $\alpha^2 < 2$, entonces $\alpha \in X$ sería el máximo de $X$, lo que contradice (1).
    \item Si $\alpha^2 > 2$, entonces $\alpha \in Y$. Por (2) existe $\beta \in Y$ con $\beta < \alpha$ y por (3) $\beta$ es cota superior de $X$, lo que contradice que $\alpha$ sea la menor cota superior.
  \end{itemize}
  Por lo tanto $\alpha^2 = 2$. Pero esto es imposible en $\Q$: si ${(p/q)}^2 = 2$ con $p, q$ coprimos, entonces $p^2 = 2q^2$ es par, luego $p$ es par, y entonces $4 \mid 2q^2$ implica que $q$ también es par, absurdo. Así $X$ no tiene supremo en $\Q$. El caso de $\inf Y$ es análogo.
\end{eg}

\section{Cuerpos completos}

Si $\K$ es un cuerpo ordenado no arquimediano, $\N$ está acotado superiormente. Pero además no tiene supremo: si $b$ es una cota superior de $\N$, entonces $n + 1 \leq b$ para todo $n \in \N$, es decir, $n \leq b - 1$ para todo $n$, de modo que $b - 1 < b$ también es cota superior.

\begin{definition}[Cuerpo completo]
  Un cuerpo ordenado $\K$ es \emph{completo} si todo subconjunto no vacío y acotado superiormente de $\K$ tiene supremo en $\K$.
\end{definition}

\begin{prop}
  Sea $\K$ un cuerpo ordenado completo. Entonces:
  \begin{enumerate}[label=(\arabic*)]
    \item todo subconjunto no vacío y acotado inferiormente tiene ínfimo en $\K$;
    \item $\K$ es arquimediano.
  \end{enumerate}
\end{prop}

\begin{proof}
  (1) Sea $Y \subset \K$ no vacío y acotado inferiormente. Entonces $X = -Y = \{ -y : y \in Y \}$ es no vacío y acotado superiormente, de modo que existe $a = \sup X$, y se verifica que $-a = \inf Y$.

  (2) Si $\N$ estuviese acotado superiormente, tendría un supremo $s$. Como $s - 1 < s$ no es cota superior de $\N$, existe $n \in \N$ con $n > s - 1$, y entonces $n + 1 > s$ con $n + 1 \in \N$, lo que contradice que $s$ sea cota superior.
\end{proof}

\begin{note}
  Se asume como axioma que existe un cuerpo ordenado completo, que llamamos $\R$ (y que es único salvo isomorfismo).
\end{note}

\begin{ejercicio}
  Dados $a \in \R$ con $a > 0$ y $n \in \N$, existe un único $b > 0$ tal que $b^n = a$.

  \emph{Sugerencia.} Definir $X = \{ x \in \R : x \geq 0,\ x^n < a \}$ e $Y = \{ y \in \R : y > 0,\ y^n > a \}$ e imitar el ejemplo anterior. Para ello, probar y usar que, dado $x > 0$, existe una constante $A_n(x) > 0$ tal que ${(x + d)}^n \leq x^n + A_n(x)\, d$ para todo $0 < d < 1$.
\end{ejercicio}

\begin{definition}[Conjunto denso]
  Un conjunto $X \subset \R$ es \emph{denso} en $\R$ si todo intervalo abierto $(a, b)$ con $a < b$ contiene algún punto de $X$.
\end{definition}

\begin{prop}
  $\Q$ es denso en $\R$.
\end{prop}

\begin{proof}
  Sean $a < b$. Como $b - a > 0$ y $\R$ es arquimediano, existe $p \in \N$ con $0 < \frac{1}{p} < b - a$. Sea $A = \{ m \in \Z : \frac{m}{p} \geq b \}$. Por la propiedad arquimediana $A \neq \varnothing$, y está acotado inferiormente por $bp$. Entonces $A$ tiene un mínimo $m_0$, que cumple $\frac{m_0}{p} \geq b$ y $\frac{m_0 - 1}{p} < b$.

  Veamos que $\frac{m_0 - 1}{p} > a$. Si no, tendríamos $\frac{m_0 - 1}{p} \leq a < b \leq \frac{m_0}{p}$, y entonces $b - a \leq \frac{m_0}{p} - \frac{m_0 - 1}{p} = \frac{1}{p}$, lo que es absurdo. Luego $\frac{m_0 - 1}{p} \in (a, b) \cap \Q$.
\end{proof}

\begin{prop}
  $\R \setminus \Q$ es denso en $\R$.
\end{prop}

\begin{proof}
  Sean $a < b$. Como $\Q$ es denso, existe $r \in \Q$ con $a - \sqrt{2} < r < b - \sqrt{2}$. Entonces $r + \sqrt{2} \in (a, b)$ y $r + \sqrt{2} \notin \Q$, pues de lo contrario $\sqrt{2} = (r + \sqrt{2}) - r$ sería racional.
\end{proof}

\section{Introducción a la cardinalidad}

\begin{definition}[Equipotencia]
  Dos conjuntos $X$ e $Y$ tienen el \emph{mismo cardinal} (o son \emph{coordinables} o \emph{equipotentes}) si existe una función biyectiva $f\colon X \to Y$. Se escribe $X \sim Y$ o $\card(X) = \card(Y)$.
\end{definition}

La relación $\sim$ es de equivalencia.

Dado $n \in \N$, sea $I_n = \{1, 2, \dots, n\}$.

\begin{lemma}
  Para todo $n \in \N$ no existe una función suryectiva $h\colon I_n \to I_{n+1}$.
\end{lemma}

\begin{proof}
  Por inducción en $n$. Para $n = 1$, $I_1$ tiene un solo elemento, así que su imagen no puede ser $I_2$.

  Supongamos que el resultado vale para $n - 1$ y que existe $h\colon I_n \to I_{n+1}$ suryectiva. Sea $r$ la permutación de $I_{n+1}$ que intercambia $h(n)$ con $n + 1$. Entonces $r \circ h$ es suryectiva y $(r \circ h)(n) = n + 1$. Cada $j \in I_n$ es imagen de algún $k \in I_n$ y necesariamente $k \neq n$ (pues $(r \circ h)(n) = n + 1 \notin I_n$). Por lo tanto, la restricción $(r \circ h)|_{I_{n-1}}\colon I_{n-1} \to I_n$ es suryectiva, lo que contradice la hipótesis inductiva.
\end{proof}

\begin{theorem}
  Sean $n, m \in \N$. Entonces $I_n \sim I_m$ si y solo si $n = m$.
\end{theorem}

\begin{proof}
  Si $n = m$ es trivial. Recíprocamente, sea $f\colon I_n \to I_m$ biyectiva y supongamos, sin pérdida de generalidad, que $n < m$ (si no, se usa $f^{-1}$). Definimos $g\colon I_m \to I_{n+1}$ por
  \[
    g(k) =
    \begin{cases}
      k & \text{si } 1 \leq k \leq n + 1, \\
      1 & \text{si } k > n + 1.
    \end{cases}
  \]
  Esta función es suryectiva y, como $f$ es biyectiva, $g \circ f\colon I_n \to I_{n+1}$ también lo es, lo que contradice el lema.
\end{proof}

\begin{definition}[Conjunto finito]
  Un conjunto $X$ es \emph{finito} si es vacío o si existe $n \in \N$ tal que $X \sim I_n$; en ese caso $\card(X) = n$ (y $\card(\varnothing) = 0$). Por el teorema anterior, $n$ está bien definido. Un conjunto que no es finito se dice \emph{infinito}.
\end{definition}


\clase{30/08/2024}{Conjuntos numerables}
\section{Conjuntos numerables}

\begin{definition}[Conjunto numerable]
  Un conjunto \(X\) es \emph{numerable} si \(X \sim \N\). Cada biyección \(\N \to X\) se llama una \emph{enumeración} de \(X\). Un conjunto es \emph{a lo sumo numerable} si es finito o numerable.
\end{definition}

\begin{eg}
  El conjunto de los números pares \(P = \{ 2n : n \in \N \}\) es numerable, pues \(f\colon \N \to P\), \(f(n) = 2n\), es biyectiva.
\end{eg}

\begin{eg}
  \(\Z\) es numerable. La función \(f\colon \Z \to \N\) dada por
  \[
    f(n) =
    \begin{cases}
      2n      & \text{si } n > 0,    \\
      -2n + 1 & \text{si } n \leq 0,
    \end{cases}
  \]
  es biyectiva (los positivos van a los pares y los no positivos a los impares), y su inversa es una enumeración de \(\Z\).
\end{eg}

\begin{theorem}[Cantor]
  Sea \(X\) un conjunto y \(\mathcal{P}(X) = \{ A : A \subset X \}\). Entonces \(\card(X) \neq \card(\mathcal{P}(X))\).
\end{theorem}

\begin{proof}
  Supongamos que existe \(f\colon X \to \mathcal{P}(X)\) biyectiva. Para cada \(x \in X\) puede ocurrir \(x \in f(x)\) o \(x \notin f(x)\). Sea
  \[
    B = \{ x \in X : x \notin f(x) \} \subset X.
  \]
  Como \(f\) es suryectiva, existe \(y \in X\) tal que \(f(y) = B\).
  \begin{itemize}
    \item Si \(y \in B = f(y)\), entonces \(y \in f(y)\), lo que contradice la definición de \(B\).
    \item Si \(y \notin B = f(y)\), entonces \(y \notin f(y)\), y por definición \(y \in B\).
  \end{itemize}
  En ambos casos llegamos a un absurdo, de modo que no existe tal \(f\).
\end{proof}

\begin{definition}[Orden entre cardinales]
  Se escribe \(\card(X) \leq \card(Y)\) si existe una función inyectiva \(X \to Y\). Se escribe \(\card(X) < \card(Y)\) si \(\card(X) \leq \card(Y)\) pero \(X \not\sim Y\).
\end{definition}

\begin{theorem}[Cantor--Schröder--Bernstein]
  Si existen funciones inyectivas \(f\colon X \to Y\) y \(g\colon Y \to X\), entonces existe una función biyectiva \(h\colon X \to Y\).
\end{theorem}

\begin{proof}
  Basta encontrar particiones \(X = X_1 \cup X_2\) e \(Y = Y_1 \cup Y_2\) tales que \(f\colon X_1 \to Y_1\) y \(g\colon Y_2 \to X_2\) sean biyectivas: entonces
  \[
    h(x) =
    \begin{cases}
      f(x)      & \text{si } x \in X_1, \\
      g^{-1}(x) & \text{si } x \in X_2,
    \end{cases}
  \]
  es biyectiva.

  Definimos \(\phi\colon \mathcal{P}(X) \to \mathcal{P}(X)\) por \(\phi(A) = X \setminus g\bigl(Y \setminus f(A)\bigr)\).

  \emph{\(\phi\) es creciente}, es decir, \(A \subset B \Rightarrow \phi(A) \subset \phi(B)\). En efecto,
  \begin{multline*}
    A \subset B \Rightarrow f(A) \subset f(B) \Rightarrow Y \setminus f(B) \subset Y \setminus f(A) \\
    \Rightarrow g\bigl(Y \setminus f(B)\bigr) \subset g\bigl(Y \setminus f(A)\bigr) \Rightarrow \phi(A) \subset \phi(B).
  \end{multline*}

  Sea \(\mathscr{C} = \{ C \subset X : \phi(C) \subset C \}\). Es no vacío pues \(X \in \mathscr{C}\). Sea \(A = \bigcap_{C \in \mathscr{C}} C\). Para cada \(C \in \mathscr{C}\) se tiene \(A \subset C\), luego \(\phi(A) \subset \phi(C) \subset C\); por lo tanto \(\phi(A) \subset A\). Aplicando \(\phi\) nuevamente, \(\phi(\phi(A)) \subset \phi(A)\), es decir, \(\phi(A) \in \mathscr{C}\) y entonces \(A \subset \phi(A)\). Concluimos que \(A = \phi(A)\).

  Sean \(X_1 = A\), \(X_2 = X \setminus A\), \(Y_1 = f(A)\) e \(Y_2 = Y \setminus f(A)\). De \(A = \phi(A)\) se obtiene \(X \setminus A = g\bigl(Y \setminus f(A)\bigr)\), es decir, \(X_2 = g(Y_2)\). Como \(f\) y \(g\) son inyectivas, \(f\colon X_1 \to Y_1\) y \(g\colon Y_2 \to X_2\) son biyectivas.
\end{proof}

\begin{prop}
  La relación \(\leq\) entre cardinales es reflexiva, transitiva y antisimétrica.
\end{prop}

\begin{proof}
  Reflexiva: la identidad es inyectiva. Transitiva: la composición de funciones inyectivas es inyectiva. Antisimétrica: si \(\card(X) \leq \card(Y)\) y \(\card(Y) \leq \card(X)\), por el teorema de Cantor--Schröder--Bernstein \(X \sim Y\).
\end{proof}

\begin{eg}
  \(\N \sim \Z \sim \Q\).
\end{eg}

\begin{proof}
  Ya vimos que \(\N \sim \Z\). La inclusión \(\Z \to \Q\) es inyectiva. Recíprocamente, todo \(q \in \Q\) se escribe de manera única como \(q = a/b\) con \(a \in \Z\), \(b \in \N\) coprimos. Definimos \(F\colon \Q \to \Z\) por \(F(q) = \operatorname{sgn}(a) \cdot 2^{|a|} \cdot 3^{b}\). Por el teorema fundamental de la aritmética, \(F\) es inyectiva. Por el teorema de Cantor--Schröder--Bernstein, \(\Z \sim \Q\).
\end{proof}

\begin{eg}
  \(\N \times \N\) es numerable.
\end{eg}

\begin{proof}
  La función \(n \mapsto (1, n)\) es inyectiva de \(\N\) en \(\N \times \N\), y la función \((n, m) \mapsto 2^n 3^m\) es inyectiva de \(\N \times \N\) en \(\N\) por el teorema fundamental de la aritmética. Por el teorema de Cantor--Schröder--Bernstein, \(\N \times \N \sim \N\).
\end{proof}

\begin{eg}
  El conjunto de sucesiones de ceros y unos, \(A = {\{0, 1\}}^{\N} = \{ {(a_n)}_{n \geq 1} : a_n \in \{0, 1\} \}\), no es numerable.
\end{eg}

\begin{proof}
  Supongamos que es numerable y escribamos \(A = \{ a^1, a^2, \dots, a^j, \dots \}\) con \(a^j = {(a^j_n)}_{n \geq 1}\). Consideremos la sucesión \(b = {(b_n)}_{n \geq 1}\) dada por \(b_n = 1 - a^n_n\), que en el lugar \(n\) tiene lo contrario de lo que tiene la \(n\)-ésima sucesión de la lista. Entonces \(b \in A\) pero \(b \neq a^n\) para todo \(n\) (difieren en el lugar \(n\)), lo que es absurdo.
\end{proof}

Este argumento (\emph{diagonal de Cantor}) puede adaptarse para probar que el intervalo \((0, 1)\) no es numerable. A continuación damos otra demostración, basada en el siguiente teorema.

\begin{theorem}[Principio de encaje de intervalos]
  Sea \(I_1 \supset I_2 \supset I_3 \supset \cdots\) una sucesión decreciente de intervalos cerrados y acotados, \(I_n = [a_n, b_n]\). Entonces
  \[
    \bigcap_{n=1}^{\infty} I_n = [\alpha, \beta], \qquad \alpha = \sup_n a_n, \quad \beta = \inf_n b_n,
  \]
  y en particular la intersección es no vacía.
\end{theorem}

\begin{proof}
  Como \(I_{n+1} \subset I_n\), se tiene \(a_n \leq a_{n+1} \leq b_{n+1} \leq b_n\). Luego, para todos \(m, k \in \N\), tomando \(j = \max\{m, k\}\) resulta \(a_m \leq a_j \leq b_j \leq b_k\). Entonces el conjunto \(A = \{ a_n : n \in \N \}\) está acotado superiormente por cada \(b_k\), y \(B = \{ b_n : n \in \N \}\) está acotado inferiormente por cada \(a_m\). Así existen \(\alpha = \sup A\) y \(\beta = \inf B\), y \(\alpha \leq b_k\) para todo \(k\), de donde \(\alpha \leq \beta\). Se cumple
  \[
    a_n \leq \alpha \leq \beta \leq b_n \quad \text{para todo } n \in \N,
  \]
  y por lo tanto \([\alpha, \beta] \subset I_n\) para todo \(n\), es decir, \([\alpha, \beta] \subset \bigcap_n I_n\).

  Recíprocamente, si \(x < \alpha\), existe \(a_i\) con \(x < a_i\), de modo que \(x \notin I_i\). Análogamente, si \(x > \beta\) existe \(b_i < x\) y \(x \notin I_i\). Luego \(\bigcap_n I_n \subset [\alpha, \beta]\).
\end{proof}

\begin{theorem}
  \(\R\) no es numerable.
\end{theorem}

\begin{proof}
  Supongamos que existe una enumeración \(\N \to \R\), \(n \mapsto x_n\). Construimos recursivamente intervalos cerrados \([0, 1] = I_0 \supset I_1 \supset I_2 \supset \cdots\) tales que \(I_n\) tiene longitud \(3^{-n}\) y \(x_n \notin I_n\). Dado \(I_{n-1}\), lo dividimos en tres intervalos cerrados de igual longitud; como \(x_n\) pertenece a lo sumo a dos de ellos, alguno no lo contiene, y lo llamamos \(I_n\).

  Por el principio de encaje, existe \(x \in \bigcap_{n \geq 1} I_n \subset [0, 1]\). Pero \(x \neq x_n\) para todo \(n\) (pues \(x \in I_n\) y \(x_n \notin I_n\)), en contradicción con que la enumeración recorre todo \(\R\).
\end{proof}

\begin{note}
  Otra idea para ver que \(\R\) no es numerable es la siguiente. Si \(\R = \{x_n : n \in \N\}\), fijado \(0 < \e < 1\) se asigna a cada \(x_n\) un intervalo abierto \(I_n\) centrado en \(x_n\) de longitud \(\e\, 2^{-n}\). Estos intervalos cubrirían \([0, 1]\) con longitud total \(\sum_n \e\, 2^{-n} = \e < 1\), lo cual es imposible, pues un intervalo cerrado no puede cubrirse con intervalos abiertos de longitud total menor que la suya. Este hecho se probará más adelante.
\end{note}

\section{Propiedades de los conjuntos numerables}

\begin{theorem}
  Si \(X\) es numerable e \(Y \subset X\), entonces \(Y\) es a lo sumo numerable.
\end{theorem}

\begin{proof}
  Supongamos que \(Y\) es infinito y escribamos \(X = \{ x_n : n \in \N \}\) (enumeración). Definimos recursivamente
  \[
    n_1 = \min\{ n \in \N : x_n \in Y \}, \qquad n_{k+1} = \min\{ n > n_k : x_n \in Y \},
  \]
  que están bien definidos porque \(Y\) es infinito. La sucesión \((n_k)\) es estrictamente creciente. Sea \(g\colon \N \to Y\) dada por \(g(k) = x_{n_k}\).

  \(g\) es inyectiva pues \(k \neq l \Rightarrow n_k \neq n_l\). Veamos que es suryectiva: sea \(y = x_j \in Y\). Como \((n_k)\) es estrictamente creciente, existe un menor \(k\) con \(n_k \geq j\). Si \(k = 1\), entonces \(n_1 \geq j\) y, como \(x_j \in Y\), \(n_1 \leq j\); luego \(n_1 = j\). Si \(k > 1\), entonces \(n_{k-1} < j \leq n_k\) y como \(x_j \in Y\), la definición de \(n_k\) como mínimo da \(n_k \leq j\); luego \(n_k = j\). En ambos casos \(y = g(k)\).
\end{proof}

\begin{corollary}
  Si \(f\colon X \to Y\) es inyectiva e \(Y\) es a lo sumo numerable, entonces \(X\) es a lo sumo numerable.
\end{corollary}

\begin{proof}
  Como \(f\) es inyectiva, \(X \sim f(X)\), y \(f(X) \subset Y\) es a lo sumo numerable por el teorema anterior.
\end{proof}

\begin{theorem}
  Si \(f\colon X \to Y\) es suryectiva y \(X\) es a lo sumo numerable, entonces \(Y\) es a lo sumo numerable.
\end{theorem}

\begin{proof}
  Como \(f\) es suryectiva, existe \(g\colon Y \to X\) con \(f \circ g = \id_Y\) (inversa a derecha). Entonces \(g\) es inyectiva y, por el corolario, \(Y\) es a lo sumo numerable.
\end{proof}

\begin{theorem}
  Si \((X_n)_{n \in \N}\) es una familia de conjuntos numerables, entonces \(X = \bigcup_{n \in \N} X_n\) es numerable.
\end{theorem}

\begin{proof}
  Para cada \(n\) elegimos una biyección \(f_n\colon \N \to X_n\). Definimos \(F\colon \N \times \N \to X\) por \(F(n, m) = f_n(m)\). \(F\) es suryectiva: si \(x \in X\), existe \(n\) con \(x \in X_n\) y existe \(m\) con \(x = f_n(m) = F(n, m)\). Como \(\N \times \N\) es numerable, \(X\) es a lo sumo numerable por el teorema anterior; y como contiene al conjunto infinito \(X_1\), es numerable.
\end{proof}

\begin{eg}
  El conjunto \(\Q[x]\) de polinomios con coeficientes racionales es numerable.
\end{eg}

\begin{proof}
  Para cada \(n\), sea \(\Q_n[x] = \{ p \in \Q[x] : \operatorname{gr}(p) \leq n \}\). La función \(\Q^{n+1} \to \Q_n[x]\), \((a_0, \dots, a_n) \mapsto a_0 + a_1 x + \cdots + a_n x^n\), es biyectiva. Como \(\Q^{n+1}\) es numerable (producto finito de numerables) y \(\Q[x] = \bigcup_{n \geq 0} \Q_n[x]\) es unión numerable de conjuntos numerables, \(\Q[x]\) es numerable.
\end{proof}

\begin{definition}[Números algebraicos y trascendentes]
  Un número real es \emph{algebraico} si es raíz de algún polinomio no nulo con coeficientes enteros; por ejemplo, \(\sqrt{2}\) es raíz de \(x^2 - 2\). Un número real que no es algebraico se llama \emph{trascendente}.
\end{definition}

\begin{ejercicio}
  Probar que el conjunto de números algebraicos es numerable. Concluir que existen números trascendentes (de hecho, el conjunto de números trascendentes no es numerable).
\end{ejercicio}


\clase{03/09/2024}{Cardinales y números reales}
\section{Operaciones con cardinales}

Dados cardinales $n, m$ (no necesariamente finitos) y conjuntos \emph{disjuntos} $X, Y$ con $\card(X) = n$ y $\card(Y) = m$, se definen:
\begin{enumerate}[label=(\arabic*)]
  \item \textbf{Suma:} $n + m = \card(X \cup Y)$.
  \item \textbf{Producto:} $n \cdot m = \card(X \times Y)$.
  \item \textbf{Potencia:} $n^m = \card\bigl(X^Y\bigr)$, donde $X^Y = \{ f\colon Y \to X \}$.
\end{enumerate}

\begin{note}
  Suponer que $X \cap Y = \varnothing$ no es restrictivo: siempre se pueden reemplazar $X$ e $Y$ por $X \times \{1\} \sim X$ e $Y \times \{2\} \sim Y$, que son disjuntos.
\end{note}

\begin{note}
  Las definiciones no dependen de los conjuntos elegidos. Si $X \sim \tilde{X}$ e $Y \sim \tilde{Y}$, con $X \cap Y = \varnothing = \tilde{X} \cap \tilde{Y}$, y $f\colon X \to \tilde{X}$, $g\colon Y \to \tilde{Y}$ son biyecciones, entonces $h\colon X \cup Y \to \tilde{X} \cup \tilde{Y}$ dada por
  \[
    h(z) =
    \begin{cases}
      f(z) & \text{si } z \in X, \\
      g(z) & \text{si } z \in Y,
    \end{cases}
  \]
  es biyectiva, así que $\card(X \cup Y) = \card(\tilde{X} \cup \tilde{Y})$. Para el producto y la potencia se procede de manera similar.
\end{note}

En lo que sigue, sean $n = \card(X)$, $m = \card(Y)$ y $p = \card(Z)$, con $X, Y, Z$ disjuntos dos a dos.

\begin{prop}
  La suma cumple:
  \begin{enumerate}[label=(\arabic*)]
    \item Conmutatividad: $n + m = m + n$, pues $X \cup Y = Y \cup X$.
    \item Asociatividad: $(n + m) + p = n + (m + p)$.
    \item Neutro: $0 + n = n$, pues $\varnothing \cup X = X$.
  \end{enumerate}
\end{prop}

\begin{prop}
  El producto cumple:
  \begin{enumerate}[label=(\arabic*)]
    \item Conmutatividad: $n \cdot m = m \cdot n$, pues $X \times Y \sim Y \times X$.
    \item $0 \cdot n = 0$, pues $\varnothing \times X = \varnothing$.
    \item $1 \cdot n = n$, pues $\{1\} \times X \sim X$ (la función $(1, x) \mapsto x$ es biyectiva).
    \item Distributividad: $n \cdot (m + p) = n \cdot m + n \cdot p$, pues $X \times (Y \cup Z) = (X \times Y) \cup (X \times Z)$.
  \end{enumerate}
\end{prop}

\begin{note}
  No vale la ley de cancelación: $\aleph_0 + 1 = \aleph_0 + 2$ y $\aleph_0 \cdot 1 = \aleph_0 \cdot 2$, pero $1 \neq 2$.
\end{note}

\begin{eg}
  Sean $\aleph_0 = \card(\N)$ y $c = \card(\R)$. Entonces:
  \begin{enumerate}[label=(\arabic*)]
    \item $\aleph_0 + \aleph_0 = \aleph_0$ y $\aleph_0 \cdot \aleph_0 = \aleph_0$.
    \item $c + c = c$ y $c \cdot c = c$.
  \end{enumerate}
\end{eg}

\begin{proof}
  (1) Los pares $\{2n : n \in \N\}$ y los impares $\{2n - 1 : n \in \N\}$ son numerables, disjuntos y su unión es $\N$; luego $\aleph_0 + \aleph_0 = \aleph_0$. Además $\N \times \N \sim \N$, de modo que $\aleph_0 \cdot \aleph_0 = \aleph_0$.

  (2) Veamos primero que todo intervalo no degenerado tiene cardinal $c$.
  \begin{itemize}
    \item $\arctan\colon \R \to (-\frac{\pi}{2}, \frac{\pi}{2})$ es una biyección, y la función afín $x \mapsto a + (b - a)x$ es una biyección entre $(0, 1)$ y $(a, b)$. Luego $(a, b) \sim \R$ para todo $a < b$.
    \item Si $I$ es un intervalo con extremos $a < b$, entonces $(a, b) \subset I \subset \R$, y por el teorema de Cantor--Schröder--Bernstein, $\card(I) = c$. En particular $(0, 1) \sim [0, 1] \sim \R$.
  \end{itemize}
  Como $(0, 1) = (0, \tfrac{1}{2}] \cup (\tfrac{1}{2}, 1)$ es una unión disjunta de dos intervalos, $c + c = \card(0, 1) = c$.

  Para $c \cdot c = c$ basta ver que $(0, 1) \times (0, 1) \sim (0, 1)$. La función $x \mapsto (x, \tfrac{1}{2})$ es inyectiva de $(0, 1)$ en $(0, 1) \times (0, 1)$. Recíprocamente, escribimos cada $x, y \in (0, 1)$ con su desarrollo decimal infinito (el que no termina en infinitos $9$) $x = 0, x_1 x_2 x_3 \dots$ e $y = 0, y_1 y_2 y_3 \dots$ y definimos $(x, y) \mapsto 0, x_1 y_1 x_2 y_2 \dots$. Esta función es inyectiva, pues el desarrollo obtenido tampoco termina en infinitos $9$, y de él se recuperan $x$ e $y$. Por el teorema de Cantor--Schröder--Bernstein, $(0, 1) \times (0, 1) \sim (0, 1)$.
\end{proof}

\begin{prop}
  $n^m \cdot n^p = n^{m + p}$.
\end{prop}

\begin{proof}
  Hay que ver que $X^Y \times X^Z \sim X^{Y \cup Z}$. La función que a un par $(f, g)$ con $f\colon Y \to X$ y $g\colon Z \to X$ le asigna la función $h\colon Y \cup Z \to X$ dada por $h|_Y = f$ y $h|_Z = g$ (bien definida porque $Y \cap Z = \varnothing$) es biyectiva, con inversa $h \mapsto (h|_Y, h|_Z)$.
\end{proof}

\begin{prop}
  ${(n^m)}^p = n^{m p}$.
\end{prop}

\begin{proof}
  Hay que ver que ${(X^Y)}^Z \sim X^{Y \times Z}$. A cada $F\colon Z \to X^Y$ le asignamos la función $G\colon Y \times Z \to X$ dada por $G(y, z) = F(z)(y)$. Esta correspondencia es biyectiva: su inversa asigna a $G$ la función $F$ dada por $F(z) = G(\cdot, z)$.
\end{proof}

\begin{theorem}
  Sea $n = \aleph_0$ o $n = c$ y sea $m$ un cardinal con $2 \leq m \leq 2^n$. Entonces $m^n = 2^n$.
\end{theorem}

\begin{proof}
  Si $\card(X) \leq \card(Y)$ y $i\colon X \to Y$ es inyectiva, entonces $\varphi \mapsto i \circ \varphi$ es una función inyectiva de $X^Z$ en $Y^Z$, de modo que $\card(X)^{\card(Z)} \leq \card(Y)^{\card(Z)}$. Usando esto,
  \[
    2^n \leq m^n \leq {(2^n)}^n = 2^{n \cdot n} = 2^n,
  \]
  pues $n \cdot n = n$ tanto para $n = \aleph_0$ como para $n = c$. Por el teorema de Cantor--Schröder--Bernstein, $m^n = 2^n$.
\end{proof}

\begin{note}
  ${\{0, 1\}}^{\N} \sim \mathcal{P}(\N)$, pues a cada $f\colon \N \to \{0, 1\}$ se le asigna el subconjunto $A_f = \{ n \in \N : f(n) = 1 \}$ (su función característica es $f$).
\end{note}

\begin{theorem}
  $\R \sim {\{0, 1\}}^{\N}$, es decir, $2^{\aleph_0} = c$.
\end{theorem}

\begin{proof}
  Usaremos el teorema de Cantor--Schröder--Bernstein con $(0, 1) \sim \R$.

  Todo $x \in (0, 1)$ tiene un único desarrollo binario $x = \sum_{n \geq 1} \frac{x_n}{2^n}$ con $x_n \in \{0, 1\}$ que no termina en infinitos $1$. La función $x \mapsto {(x_n)}_{n \geq 1}$ es inyectiva de $(0, 1)$ en ${\{0, 1\}}^{\N}$.

  Recíprocamente, definimos $g\colon {\{0, 1\}}^{\N} \to [0, 1]$ por $g\bigl({(x_n)}_{n}\bigr) = \sum_{n \geq 1} \frac{x_n}{3^n}$. Veamos que es inyectiva. Si ${(x_n)} \neq {(y_n)}$, sea $k$ el primer índice con $x_k \neq y_k$, digamos $x_k = 1$ e $y_k = 0$. Entonces
  \[
    g(x) - g(y) = \frac{1}{3^k} + \sum_{n > k} \frac{x_n - y_n}{3^n} \geq \frac{1}{3^k} - \sum_{n > k} \frac{1}{3^n} = \frac{1}{3^k} - \frac{1}{2 \cdot 3^k} > 0.
  \]
  Luego $g$ es inyectiva (la imagen de $g$ es el conjunto de Cantor, que se verá más adelante). Por el teorema de Cantor--Schröder--Bernstein, $(0, 1) \sim {\{0, 1\}}^{\N}$.
\end{proof}

\section{Hipótesis del continuo}

La \emph{hipótesis del continuo} afirma que no existe ningún cardinal $\kappa$ con $\aleph_0 < \kappa < c = 2^{\aleph_0}$.

Es independiente de los axiomas usuales de la teoría de conjuntos (ZFC, es decir, Zermelo--Fraenkel más el axioma de elección): Gödel probó en 1940 que no puede refutarse y Cohen probó en 1963 que no puede demostrarse.

\section{Construcción de los números reales}

Consideremos sucesiones de números racionales ${(a_n)}_{n \in \N} \subset \Q$. Decimos que ${(a_n)}$ \emph{tiende a $0$}, y escribimos $a_n \to 0$, si para todo $\e > 0$ existe $n_0 \in \N$ tal que $|a_n| < \e$ para todo $n > n_0$.

\begin{definition}[Sucesión de Cauchy]
  Una sucesión ${(a_n)}_{n \in \N} \subset \Q$ es de \emph{Cauchy} si para todo $\e > 0$ existe $n_0 \in \N$ tal que $|a_n - a_m| < \e$ para todos $n, m > n_0$.
\end{definition}

\begin{theorem}
  Si ${(a_n)} \subset \Q$ es convergente, es decir, existe $q \in \Q$ con $a_n - q \to 0$, entonces es de Cauchy.
\end{theorem}

\begin{proof}
  Sea $\e > 0$. Existe $n_0$ tal que $|a_n - q| < \frac{\e}{2}$ para todo $n > n_0$. Entonces, para $n, m > n_0$,
  \[
    |a_n - a_m| \leq |a_n - q| + |q - a_m| < \frac{\e}{2} + \frac{\e}{2} = \e. \qedhere
  \]
\end{proof}

\begin{theorem}
  Toda sucesión de Cauchy ${(a_n)} \subset \Q$ es acotada.
\end{theorem}

\begin{proof}
  Tomando $\e = 1$, existe $n_0$ tal que $|a_n - a_m| < 1$ para todos $n, m > n_0$. En particular, $|a_n - a_{n_0 + 1}| < 1$ para $n > n_0$, de modo que $|a_n| < |a_{n_0 + 1}| + 1$. Luego, con
  \[
    M = \max\{ |a_1|, \dots, |a_{n_0}|, |a_{n_0 + 1}| + 1 \},
  \]
  se tiene $|a_n| \leq M$ para todo $n \in \N$.
\end{proof}

\begin{definition}
  Sea $\mathscr{C}_{\Q}$ el conjunto de todas las sucesiones de Cauchy de números racionales. Dos sucesiones ${(a_n)}, {(b_n)} \in \mathscr{C}_{\Q}$ son \emph{equivalentes} si $a_n - b_n \to 0$.
\end{definition}

\begin{prop}
  La relación anterior es de equivalencia.
\end{prop}

\begin{proof}
  \begin{itemize}
    \item Reflexiva: $a_n - a_n = 0 \to 0$.
    \item Simétrica: si $a_n - b_n \to 0$, entonces $b_n - a_n = -(a_n - b_n) \to 0$.
    \item Transitiva: si $a_n - b_n \to 0$ y $b_n - c_n \to 0$, dado $\e > 0$ existen $n_0$ y $n_1$ tales que $|a_n - b_n| < \frac{\e}{2}$ para $n > n_0$ y $|b_n - c_n| < \frac{\e}{2}$ para $n > n_1$. Para $n > \max\{n_0, n_1\}$ resulta
          \[
            |a_n - c_n| \leq |a_n - b_n| + |b_n - c_n| < \e. \qedhere
          \]
  \end{itemize}
\end{proof}

\begin{definition}[Números reales]
  El conjunto $\R$ de los números reales es el conjunto de clases de equivalencia $[(a_n)]$ de sucesiones de $\mathscr{C}_{\Q}$ respecto de la relación anterior.
\end{definition}


\clase{06/09/2024}{Los números reales}
\section{Operaciones en los reales}

\begin{note}
  A cada \(q \in \Q\) se le asocia la clase \(\tilde{q}\) de la sucesión constante \((q, q, q, \dots)\). La función \(q \mapsto \tilde{q}\) es inyectiva, y de esta manera consideramos \(\Q \subset \R\).
\end{note}

\begin{definition}
  Sean \(s = [(a_n)]\) y \(t = [(b_n)]\) en \(\R\). Definimos
  \begin{itemize}
    \item \(s + t = [(a_n + b_n)]\),
    \item \(s \cdot t = [(a_n b_n)]\).
  \end{itemize}
\end{definition}

Hay que verificar que las definiciones son correctas.

\begin{itemize}
  \item \emph{Las sucesiones \((a_n + b_n)\) y \((a_n b_n)\) son de Cauchy.} Para la suma es inmediato. Para el producto, si \(|a_n|, |b_n| \leq M\) para todo \(n\) (las sucesiones de Cauchy son acotadas), entonces
        \[
          a_n b_n - a_m b_m = a_n (b_n - b_m) + b_m (a_n - a_m)
        \]
        y por lo tanto \(|a_n b_n - a_m b_m| \leq M \bigl( |b_n - b_m| + |a_n - a_m| \bigr)\), que es menor que \(\e\) para \(n, m\) grandes.
  \item \emph{No dependen de los representantes.} Sean \([(a_n)] = [(c_n)]\) y \([(b_n)] = [(d_n)]\), es decir, \(a_n - c_n \to 0\) y \(b_n - d_n \to 0\). Para la suma,
        \[
          (a_n + b_n) - (c_n + d_n) = (a_n - c_n) + (b_n - d_n) \to 0.
        \]
        Para el producto,
        \[
          a_n b_n - c_n d_n = b_n (a_n - c_n) + c_n (b_n - d_n).
        \]
        Como \((b_n)\) y \((c_n)\) son acotadas, existe \(M\) con \(|b_n|, |c_n| \leq M\) para todo \(n\), de donde \(|a_n b_n - c_n d_n| \leq M \bigl( |a_n - c_n| + |b_n - d_n| \bigr) \to 0\).
\end{itemize}

\begin{prop}
  \(\R\) con estas operaciones es un cuerpo.
\end{prop}

\begin{proof}
  Los axiomas de cuerpo (asociatividad, conmutatividad, distributividad, neutros \(\tilde{0}\) y \(\tilde{1}\), opuesto \(-[(a_n)] = [(-a_n)]\)) se heredan de \(\Q\) operando término a término. Queda como ejercicio. La existencia del inverso multiplicativo se prueba en el siguiente teorema.
\end{proof}

\begin{lemma}
  Sea \((a_n)\) una sucesión de Cauchy con \(a_n \not\to 0\). Entonces existen \(\delta \in \Q\), \(\delta > 0\), y \(n_0 \in \N\) tales que \(|a_n| > \delta\) para todo \(n > n_0\). Además, o bien \(a_n > \delta\) para todo \(n > n_0\), o bien \(a_n < -\delta\) para todo \(n > n_0\).
\end{lemma}

\begin{proof}
  Como \(a_n \not\to 0\), existe \(\e_1 > 0\) tal que \(|a_n| \geq \e_1\) para infinitos valores de \(n\). Como \((a_n)\) es de Cauchy, existe \(n_0\) tal que \(|a_n - a_m| < \frac{\e_1}{2}\) para todos \(n, m > n_0\). Elegimos \(M > n_0\) con \(|a_M| \geq \e_1\). Entonces, para \(n > n_0\),
  \[
    |a_n| \geq |a_M| - |a_n - a_M| > \e_1 - \frac{\e_1}{2} = \frac{\e_1}{2}.
  \]
  Sea \(\delta = \frac{\e_1}{2}\). Si \(a_M > 0\), entonces \(a_n > a_M - \frac{\e_1}{2} \geq \delta\) para \(n > n_0\); si \(a_M < 0\), análogamente \(a_n < a_M + \frac{\e_1}{2} \leq -\delta\).
\end{proof}

\begin{theorem}
  Si \(s \in \R\), \(s \neq 0\), entonces existe \(t \in \R\) tal que \(s \cdot t = 1\).
\end{theorem}

\begin{proof}
  Sea \(s = [(a_n)]\). Como \(s \neq \tilde{0}\), se tiene \(a_n \not\to 0\). Por el lema, existen \(\delta > 0\) y \(n_0\) tales que \(|a_n| > \delta\) para \(n > n_0\); en particular \(a_n \neq 0\) para \(n > n_0\), aunque puede haber finitos términos nulos al principio. Definimos
  \[
    b_n =
    \begin{cases}
      0              & \text{si } n \leq n_0, \\
      \dfrac{1}{a_n} & \text{si } n > n_0.
    \end{cases}
  \]
  La sucesión \((b_n)\) es de Cauchy: para \(n, m > n_0\),
  \[
    |b_n - b_m| = \frac{|a_m - a_n|}{|a_n| |a_m|} < \frac{|a_n - a_m|}{\delta^2},
  \]
  que es arbitrariamente pequeño para \(n, m\) grandes. Además \(a_n b_n = 1\) para \(n > n_0\), de modo que \((1, 1, \dots) - (a_n b_n) \to 0\) y \([(a_n b_n)] = \tilde{1}\). Es decir, \(t = [(b_n)]\) cumple \(s \cdot t = 1\).
\end{proof}

\section{Cuerpo ordenado}

Para ver que \(\R\) es un cuerpo ordenado hay que decir quiénes son los positivos.

\begin{definition}
  Un real \(s \in \R\) es \emph{positivo} si \(s \neq 0\) y \(s = [(a_n)]\) con \(a_n > 0\) para todo \(n\) mayor que cierto \(n_0\); es decir, si todos los términos son positivos a partir de cierto punto. Escribimos \(s > t\) si \(s - t\) es positivo.
\end{definition}

\begin{ejercicio}
  Probar que la definición no depende del representante elegido, que la suma y el producto de positivos son positivos, y que para todo \(s \in \R\) vale exactamente una de las opciones \(s > 0\), \(s = 0\), \(-s > 0\). (Sugerencia: usar el lema anterior.)
\end{ejercicio}

Veamos, a modo de ejemplo, cómo se prueban las propiedades de orden.

\begin{theorem}
  Si \(s, t, r \in \R\) y \(s > t\), entonces \(s + r > t + r\).
\end{theorem}

\begin{proof}
  Sean \(s = [(a_n)]\), \(t = [(b_n)]\) y \(r = [(c_n)]\). Como \((a_n + c_n) - (b_n + c_n) = a_n - b_n\) para todo \(n\),
  \[
    (s + r) - (t + r) = [(a_n - b_n)] = s - t > 0. \qedhere
  \]
\end{proof}

\begin{theorem}
  \(\R\) es arquimediano: si \(s, t > 0\), existe \(m \in \N\) tal que \(m \cdot s > t\).
\end{theorem}

\begin{proof}
  Sean \(s = [(a_n)]\) y \(t = [(b_n)]\). Como \(s > 0\), por el lema existen \(\delta \in \Q\), \(\delta > 0\), y \(n_0\) tales que \(a_n > \delta\) para \(n > n_0\). Como \((b_n)\) es de Cauchy, es acotada: existe \(B \in \Q\) con \(b_n \leq B\) para todo \(n\). Como \(\Q\) es arquimediano, existe \(m \in \N\) con \(m \delta > B + 1\). Entonces, para \(n > n_0\),
  \[
    m a_n - b_n > m \delta - B > 1.
  \]
  Luego \(m a_n - b_n \not\to 0\) y es positivo a partir de \(n_0\), es decir, \(m s - t > 0\).
\end{proof}

\begin{theorem}
  \(\Q\) es denso en \(\R\): dados \(r \in \R\) y \(\e \in \Q\), \(\e > 0\), existe \(q \in \Q\) tal que \(|r - q| < \e\).
\end{theorem}

\begin{proof}
  Sea \(r = [(a_n)]\) con \((a_n) \subset \Q\) de Cauchy. Existe \(n_0\) tal que \(|a_n - a_m| < \frac{\e}{2}\) para todos \(n, m > n_0\). Fijamos \(l > n_0\) y sea \(q = \tilde{a_l} = [(a_l, a_l, \dots)]\). Entonces
  \[
    r - q = [(a_n - a_l)_n], \qquad \e - (r - q) = [(\e - (a_n - a_l))_n],
  \]
  y para \(n > n_0\) se cumple \(|a_n - a_l| < \frac{\e}{2}\), de modo que \(\e - (a_n - a_l) > \frac{\e}{2} > 0\) y \(\e + (a_n - a_l) > \frac{\e}{2}\) (estas sucesiones no tienden a \(0\)). Luego \(r - q < \e\) y \(q - r < \e\), es decir, \(|r - q| < \e\).
\end{proof}

\section{Propiedad del supremo}

Sea \(S \subset \R\) un conjunto no vacío y acotado superiormente por \(M \in \R\). Fijamos \(s_0 \in S\) y construimos recursivamente dos sucesiones \((u_n)_{n \geq 0}\) y \((l_n)_{n \geq 0}\) con \(u_0 = M\) y \(l_0 = s_0\). Dados \(u_n\) y \(l_n\), sea \(m_n = \frac{l_n + u_n}{2}\) el punto medio.
\begin{enumerate}[label=(\roman*)]
  \item Si \(m_n\) es cota superior de \(S\), definimos \(u_{n+1} = m_n\) y \(l_{n+1} = l_n\).
  \item Si \(m_n\) no es cota superior de \(S\), definimos \(u_{n+1} = u_n\) y \(l_{n+1} = m_n\).
\end{enumerate}

Por inducción se verifica que para todo \(n\):
\begin{enumerate}[label=(\alph*)]
  \item \(u_n\) es cota superior de \(S\);
  \item existe \(s_n \in S\) con \(l_n \leq s_n\) (en el caso (ii), \(m_n\) no es cota superior, y existe \(s \in S\) con \(s > m_n\));
  \item \(l_n \leq s_n \leq u_n\), y además \(l_n \leq l_{n+1} \leq u_{n+1} \leq u_n\);
  \item \(u_n - l_n = \dfrac{M - s_0}{2^n}\), pues en ambos casos \(u_{n+1} - l_{n+1} = \frac{u_n - l_n}{2}\).
\end{enumerate}

\begin{lemma}
  Las sucesiones \((u_n)\) y \((l_n)\) son de Cauchy en \(\R\).
\end{lemma}

\begin{proof}
  Sea \(\e > 0\). Como \(\R\) es arquimediano y \(2^n \geq n\), existe \(n_0\) tal que \(\frac{M - s_0}{2^n} < \e\) para todo \(n > n_0\). Si \(m \geq n > n_0\), por (c) se tiene \(l_n \leq l_m \leq u_m \leq u_n\), de donde
  \[
    0 \leq l_m - l_n \leq u_n - l_n < \e, \qquad 0 \leq u_n - u_m \leq u_n - l_n < \e. \qedhere
  \]
\end{proof}

\begin{lemma}
  Toda sucesión de Cauchy \((u_n)\) en \(\R\) converge a algún \(u \in \R\).
\end{lemma}

\begin{proof}
  Como \(\Q\) es denso en \(\R\), para cada \(n\) existe \(q_n \in \Q\) con \(|u_n - q_n| < \frac{1}{n}\). Veamos que \((q_n)\) es de Cauchy en \(\Q\). Dado \(\e > 0\) racional, existe \(n_0\) tal que \(|u_n - u_m| < \frac{\e}{3}\) para \(n, m > n_0\), y por la arquimedianidad existe \(n_1\) con \(\frac{1}{n} < \frac{\e}{3}\) para \(n > n_1\). Si \(n, m > \max\{n_0, n_1\}\),
  \[
    |q_n - q_m| \leq |q_n - u_n| + |u_n - u_m| + |u_m - q_m| < \e.
  \]
  Sea \(u = [(q_n)] \in \R\). Para cada \(n\), \(|u_n - u| \leq |u_n - \tilde{q_n}| + |\tilde{q_n} - u|\). El primer término es menor que \(\frac{1}{n}\). Para el segundo, \(\tilde{q_n} - u = [(q_n - q_k)_k]\) y, dado \(\e > 0\), existe \(n_2\) tal que \(|q_n - q_k| < \e\) para \(n, k > n_2\), de modo que \(|\tilde{q_n} - u| \leq \e\) para \(n > n_2\). Por lo tanto \(u_n \to u\).
\end{proof}

\begin{lemma}
  Existe \(u \in \R\) tal que \(u_n \to u\) y \(l_n \to u\).
\end{lemma}

\begin{proof}
  Por los lemas anteriores existe \(u\) con \(u_n \to u\). Como \(u_n - l_n = \frac{M - s_0}{2^n} \to 0\), también \(l_n \to u\).
\end{proof}

\begin{theorem}
  \(\R\) tiene la propiedad del supremo: todo subconjunto no vacío y acotado superiormente tiene supremo.
\end{theorem}

\begin{proof}
  Probamos que el \(u\) del lema anterior es \(\sup S\).

  \emph{\(u\) es cota superior.} Si existiera \(s \in S\) con \(u < s\), sea \(\e = s - u > 0\). Como \(u_n \to u\), existe \(n\) con \(u_n - u < \e\), es decir, \(u_n < u + \e = s\), en contradicción con que \(u_n\) es cota superior de \(S\).

  \emph{\(u\) es la menor cota superior.} Sea \(c < u\). Como \(l_n \to u\), existe \(n\) con \(l_n > c\). Por (b), \(s_n \geq l_n > c\) con \(s_n \in S\), de modo que \(c\) no es cota superior de \(S\).
\end{proof}


\clase{10/09/2024}{Sucesiones}
\section{Sucesiones}

Una \emph{sucesión} de números reales es una función \(x\colon \N \to \R\). Se escribe \(x(n) = x_n\) y se lo llama el \(n\)-ésimo término de la sucesión. La sucesión se indica \({(x_n)}_{n \in \N}\) o \((x_1, x_2, \dots)\).

Una \emph{subsucesión} de \(x\) es la restricción de \(x\) a un subconjunto infinito \(A = \{ n_1 < n_2 < \cdots \} \subset \N\). Se escribe \({(x_n)}_{n \in A}\) o \({(x_{n_i})}_{i \in \N}\).

\begin{note}
  Estrictamente, la subsucesión no está definida en \(\N\), pero se la puede considerar como una función definida en \(\N\) componiendo con la biyección \(i \mapsto n_i\). Por eso se usa la notación \({(x_{n_i})}_{i \in \N}\).
\end{note}

\begin{definition}[Límite]
  Decimos que \(a = \lim_{n \to \infty} x_n\) si
  \[
    (\forall \e > 0)(\exists n_0 \in \N) \colon |x_n - a| < \e \quad \forall n > n_0.
  \]
  Equivalentemente, si para todo \(\e > 0\) el intervalo \((a - \e, a + \e)\) contiene a todos los términos de la sucesión salvo, quizás, una cantidad finita.
\end{definition}

\begin{theorem}[Unicidad del límite]
  Si \(\lim_{n \to \infty} x_n = a\) y \(\lim_{n \to \infty} x_n = b\), entonces \(a = b\).
\end{theorem}

\begin{proof}
  Sea \(\e > 0\). Existen \(n_0, n_1 \in \N\) tales que \(|x_n - a| < \e\) para \(n > n_0\) y \(|x_n - b| < \e\) para \(n > n_1\). Luego, para \(n > \max\{n_0, n_1\}\),
  \[
    |a - b| \leq |a - x_n| + |x_n - b| < 2\e.
  \]
  Como \(\e > 0\) es arbitrario, \(|a - b| = 0\), es decir, \(a = b\).
\end{proof}

\begin{theorem}
  Si \(\lim_{n \to \infty} x_n = a\), toda subsucesión de \({(x_n)}_{n \in \N}\) converge a \(a\).
\end{theorem}

\begin{proof}
  Sea \({(x_{n_i})}_{i \in \N}\) una subsucesión. Dado \(\e > 0\), existe \(n_0\) tal que \(|x_n - a| < \e\) para \(n > n_0\). Como \(n_1 < n_2 < \cdots\) es estrictamente creciente, existe \(i_0\) con \(n_{i_0} > n_0\), y entonces para \(i > i_0\) se tiene \(n_i > n_0\) y por lo tanto \(|x_{n_i} - a| < \e\).
\end{proof}

\begin{theorem}
  Toda sucesión convergente es acotada.
\end{theorem}

\begin{proof}
  Sea \(a = \lim_{n \to \infty} x_n\). Tomando \(\e = 1\), existe \(n_1\) tal que \(x_n \in (a - 1, a + 1)\) para todo \(n > n_1\). Sean \(A = \{ x_1, \dots, x_{n_1}, a - 1, a + 1 \}\), \(c = \min A\) y \(d = \max A\). Entonces \(x_n \in [c, d]\) para todo \(n \in \N\).
\end{proof}

\begin{theorem}
  Toda sucesión monótona y acotada es convergente.
\end{theorem}

\begin{proof}
  Supongamos que \({(x_n)}\) es creciente y acotada, y sea \(a = \sup \{ x_n : n \in \N \}\). Dado \(\e > 0\), \(a - \e\) no es cota superior de \(\{ x_n \}\), luego existe \(n_0\) con \(x_{n_0} > a - \e\). Como la sucesión es creciente, para \(n > n_0\) vale
  \[
    a - \e < x_{n_0} \leq x_n \leq a < a + \e,
  \]
  es decir, \(x_n \to a\). El caso decreciente es análogo, con el ínfimo.
\end{proof}

\begin{corollary}
  Si una sucesión monótona tiene una subsucesión convergente, entonces es convergente.
\end{corollary}

\begin{proof}
  La subsucesión es acotada, y como la sucesión es monótona, está acotada por los mismos valores (el conjunto de sus términos está entre \(x_1\) y una cota de la subsucesión). Por el teorema anterior, es convergente.
\end{proof}

\begin{eg}
  Sea \(x_n = a^n\) con \(a \in \R\).
  \begin{itemize}
    \item Si \(a = 0\) o \(a = 1\), la sucesión es constante, luego convergente.
    \item Si \(a = -1\), la sucesión diverge, pues \(x_{2n} = 1 \to 1\) y \(x_{2n+1} = -1 \to -1\) son subsucesiones con límites distintos.
    \item Si \(|a| > 1\), escribimos \(|a| = 1 + h\) con \(h > 0\); por la desigualdad de Bernoulli, \(|a|^n \geq 1 + nh\), y la sucesión no es acotada, luego diverge.
    \item Si \(0 < a < 1\), la sucesión es decreciente y acotada inferiormente por \(0\), luego converge a algún \(l\). Como \(x_{n+1} = a \, x_n\), pasando al límite \(l = a l\), de donde \(l = 0\).
    \item Si \(-1 < a < 0\), se tiene \(|a^n - 0| = |a|^n \to 0\) por el caso anterior, de modo que \(a^n \to 0\).
  \end{itemize}
\end{eg}

\section{Propiedades de los límites}

\begin{theorem}
  Si \(\lim_{n \to \infty} x_n = 0\) e \({(y_n)}_{n \in \N}\) es una sucesión acotada, entonces \(\lim_{n \to \infty} x_n y_n = 0\).
\end{theorem}

\begin{proof}
  Sea \(c > 0\) con \(|y_n| < c\) para todo \(n\). Dado \(\e > 0\), existe \(n_0\) tal que \(|x_n| < \frac{\e}{c}\) para \(n > n_0\). Entonces \(|x_n y_n| < c \cdot \frac{\e}{c} = \e\) para todo \(n > n_0\).
\end{proof}

Sean \(\lim_{n \to \infty} x_n = a\) y \(\lim_{n \to \infty} y_n = b\).

\begin{prop}
  \(x_n + y_n \to a + b\).
\end{prop}

\begin{proof}
  Dado \(\e > 0\), existen \(n_0\) y \(n_1\) tales que \(|x_n - a| < \frac{\e}{2}\) para \(n > n_0\) y \(|y_n - b| < \frac{\e}{2}\) para \(n > n_1\). Para \(n > \max\{n_0, n_1\}\),
  \[
    |(x_n + y_n) - (a + b)| \leq |x_n - a| + |y_n - b| < \e. \qedhere
  \]
\end{proof}

\begin{prop}
  \(x_n - y_n \to a - b\).
\end{prop}

\begin{proof}
  Análoga a la anterior.
\end{proof}

\begin{prop}
  \(x_n y_n \to a b\).
\end{prop}

\begin{proof}
  Se tiene
  \[
    x_n y_n - a b = x_n (y_n - b) + b (x_n - a).
  \]
  Como \(y_n - b \to 0\) y \((x_n)\) es acotada por ser convergente, \(x_n (y_n - b) \to 0\). Además \(b (x_n - a) \to 0\). Luego \(x_n y_n - ab \to 0\).
\end{proof}

\begin{prop}
  Si \(b \neq 0\), entonces \(y_n \neq 0\) salvo finitos términos y \(\dfrac{x_n}{y_n} \to \dfrac{a}{b}\).
\end{prop}

\begin{proof}
  Tomando \(\e = \frac{|b|}{2}\), existe \(n_0\) tal que \(|y_n - b| < \frac{|b|}{2}\) para \(n > n_0\), de modo que \(|y_n| > \frac{|b|}{2} > 0\) para \(n > n_0\); en particular \(y_n \neq 0\). Para \(n > n_0\),
  \[
    \frac{x_n}{y_n} - \frac{a}{b} = \frac{b x_n - a y_n}{b\, y_n} = \frac{1}{b\, y_n} \bigl( b (x_n - a) - a (y_n - b) \bigr).
  \]
  La sucesión \(\bigl( \frac{1}{b y_n} \bigr)_{n > n_0}\) es acotada, pues \(\left| \frac{1}{b y_n} \right| < \frac{2}{b^2}\), y el factor entre paréntesis tiende a \(0\). Por el teorema anterior, \(\frac{x_n}{y_n} - \frac{a}{b} \to 0\).
\end{proof}

\section{Ejemplos con subsucesiones}

\begin{eg}
  Si \(a > 0\), entonces \(\sqrt[n]{a} \to 1\).
\end{eg}

\begin{proof}
  Si \(a = 1\) es trivial. Si \(a > 1\), la sucesión \(x_n = a^{1/n}\) es decreciente: \(a^{1/n} > a^{1/(n+1)} \iff a^{n+1} > a^n\), lo cual es cierto pues \(a > 1\) (se elevó a la \(n(n+1)\)). Además \(x_n > 1\), de modo que converge a un límite \(l \geq 1\). Si \(0 < a < 1\), la sucesión es creciente y acotada por \(1\), y converge a un \(l \geq a > 0\).

  En ambos casos \(l > 0\). La subsucesión \(y_n = a^{\frac{1}{n(n+1)}}\) (índices \(n(n+1)\)) converge también a \(l\), y
  \[
    y_n = a^{\frac{1}{n} - \frac{1}{n+1}} = \frac{a^{1/n}}{a^{1/(n+1)}} \longrightarrow \frac{l}{l} = 1,
  \]
  usando que \(l \neq 0\). Por unicidad del límite, \(l = 1\).
\end{proof}

\begin{eg}
  \(\sqrt[n]{n} \to 1\).
\end{eg}

\begin{proof}
  Veamos que la sucesión es decreciente a partir del tercer término. Elevando a la \(n(n+1)\),
  \[
    \sqrt[n]{n} > \sqrt[n+1]{n+1} \iff n^{n+1} > {(n+1)}^n \iff n > {\left( \frac{n+1}{n} \right)}^n = {\left( 1 + \frac{1}{n} \right)}^n.
  \]
  Esto vale para todo \(n \geq 3\) porque \({(1 + \frac{1}{n})}^n < 3\) para todo \(n\). En efecto, por el binomio de Newton,
  \[
    {\left( 1 + \frac{1}{n} \right)}^n = \sum_{k=0}^{n} \binom{n}{k} \frac{1}{n^k} \leq \sum_{k=0}^{n} \frac{1}{k!} \leq 1 + \sum_{k=1}^{n} \frac{1}{2^{k-1}} < 1 + 2 = 3,
  \]
  pues \(\binom{n}{k} \frac{1}{n^k} = \frac{n (n-1) \cdots (n-k+1)}{n^k} \cdot \frac{1}{k!} \leq \frac{1}{k!}\) y \(k! \geq 2^{k-1}\).

  Además \(\sqrt[n]{n} \geq 1\) para todo \(n\), de modo que la sucesión converge a un \(l \geq 1\); en particular \(l \neq 0\). Considerando la subsucesión \({(2n)}^{\frac{1}{2n}}\),
  \[
    l^2 = {\left( \lim_{n \to \infty} {(2n)}^{\frac{1}{2n}} \right)}^2 = \lim_{n \to \infty} {(2n)}^{\frac{1}{n}} = \lim_{n \to \infty} 2^{\frac{1}{n}} \cdot \lim_{n \to \infty} n^{\frac{1}{n}} = 1 \cdot l,
  \]
  donde usamos el ejemplo anterior. Luego \(l^2 = l\) y como \(l \neq 0\), \(l = 1\).
\end{proof}

\begin{note}
  La definición de límite puede reformularse así: \(a\) es el límite de \({(x_n)}_{n \in \N}\) si y sólo si para todo \(\e > 0\) el conjunto \(\{ n \in \N : x_n \in (a - \e, a + \e) \}\) tiene complemento finito en \(\N\).
\end{note}

\begin{theorem}
  Un número \(a \in \R\) es el límite de alguna subsucesión de \({(x_n)}_{n \in \N}\) si y sólo si para todo \(\e > 0\) existen infinitos índices \(n\) tales que \(x_n \in (a - \e, a + \e)\).
\end{theorem}

\begin{proof}
  (\(\Rightarrow\)) Sea \(a = \lim_{i \to \infty} x_{n_i}\). Dado \(\e > 0\), existe \(i_0\) tal que \(x_{n_i} \in (a - \e, a + \e)\) para todo \(i > i_0\). Como los \(n_i\) son todos distintos, hay infinitos índices \(n\) con \(x_n \in (a - \e, a + \e)\).

  (\(\Leftarrow\)) Construimos recursivamente índices \(n_1 < n_2 < \cdots\) con \(|x_{n_i} - a| < \frac{1}{i}\). Como \(\{ n : x_n \in (a - 1, a + 1) \}\) es infinito, existe \(n_1\) en él. Dados \(n_1 < \cdots < n_i\), como \(\{ n : x_n \in (a - \frac{1}{i+1}, a + \frac{1}{i+1}) \}\) es infinito, contiene algún \(n_{i+1} > n_i\). Entonces \(|x_{n_i} - a| < \frac{1}{i}\) para todo \(i\), y por lo tanto \(x_{n_i} \to a\).
\end{proof}

\section{Punto de acumulación}

\begin{definition}[Punto de acumulación de una sucesión]
  Un número \(a \in \R\) es un \emph{punto de acumulación} de \({(x_n)}_{n \in \N}\) si es el límite de alguna subsucesión de \({(x_n)}_{n \in \N}\).
\end{definition}

\begin{definition}[Límite inferior y superior]
  Sea \({(x_n)}_{n \in \N}\) una sucesión acotada, digamos \(\alpha \leq x_n \leq \beta\) para todo \(n \in \N\). Sea \(X_n = \{ x_n, x_{n+1}, \dots \}\), de modo que \([\alpha, \beta] \supset X_1 \supset X_2 \supset \cdots\). Llamando \(a_n = \inf X_n\) y \(b_n = \sup X_n\), se tiene
  \[
    \alpha \leq a_1 \leq a_2 \leq \cdots \leq a_n \leq b_n \leq \cdots \leq b_2 \leq b_1 \leq \beta.
  \]
  Por lo tanto existen
  \begin{align*}
    a & = \lim_{n \to \infty} a_n = \sup_{n} \, \inf_{k \geq n} x_k, \\
    b & = \lim_{n \to \infty} b_n = \inf_{n} \, \sup_{k \geq n} x_k.
  \end{align*}
  Se llama \emph{límite inferior} de \({(x_n)}\) a \(a = \liminf_{n \to \infty} x_n\) y \emph{límite superior} a \(b = \limsup_{n \to \infty} x_n\).
\end{definition}

\begin{eg}
  Sea \({(x_n)} = (-1, 2, -\frac{1}{2}, \frac{3}{2}, -\frac{1}{3}, \frac{4}{3}, \dots)\), es decir, \(x_{2n-1} = -\frac{1}{n}\) y \(x_{2n} = 1 + \frac{1}{n}\). Entonces
  \begin{align*}
    \inf X_{2n-2} = \inf X_{2n-1} & = -\frac{1}{n},    \\
    \sup X_{2n-1} = \sup X_{2n}   & = 1 + \frac{1}{n},
  \end{align*}
  de donde \(\liminf_{n \to \infty} x_n = 0\) y \(\limsup_{n \to \infty} x_n = 1\).
\end{eg}


\clase{13/09/2024}{Límites y series}
\section{Límite superior e inferior}

\begin{theorem}
  Sea ${(x_n)}_{n \in \N}$ una sucesión acotada. Entonces $\liminf x_n$ es el menor punto de acumulación de la sucesión y $\limsup x_n$ es el mayor.
\end{theorem}

\begin{proof}
  Sea $a = \liminf_{n \to \infty} x_n = \lim_{n \to \infty} a_n$, con $a_n = \inf X_n$.

  \emph{$a$ es punto de acumulación.} Sean $\e > 0$ y $n_0 \in \N$. Como $a_n \to a$, existe $n_1 > n_0$ con $a - \e < a_{n_1} \leq a$. Como $a_{n_1} = \inf X_{n_1}$, el número $a + \e > a_{n_1}$ no es cota inferior de $X_{n_1}$, luego existe $n \geq n_1$ con $x_n < a + \e$; además $x_n \geq a_{n_1} > a - \e$. Entonces $x_n \in (a - \e, a + \e)$ con $n > n_0$. Como $n_0$ es arbitrario, hay infinitos índices con $x_n \in (a - \e, a + \e)$, y por el teorema de la sección anterior, $a$ es punto de acumulación.

  \emph{Ningún $c < a$ es punto de acumulación.} Como $a_n \to a$ y $c < a$, existe $n_0$ con $c < a_{n_0} \leq a$. Sea $\e = a_{n_0} - c > 0$. Para $n \geq n_0$ se cumple $x_n \geq a_{n_0} = c + \e$, de modo que $x_n \notin (c - \e, c + \e)$ salvo quizás para $n < n_0$ (finitos índices). Luego $c$ no es punto de acumulación.

  La afirmación sobre $\limsup$ se prueba de manera análoga.
\end{proof}

\begin{theorem}[Bolzano--Weierstrass]
  Toda sucesión acotada de $\R$ tiene una subsucesión convergente.
\end{theorem}

\begin{proof}
  $\limsup x_n$ es punto de acumulación de ${(x_n)}$, es decir, el límite de alguna subsucesión.
\end{proof}

\begin{corollary}
  Una sucesión acotada de $\R$ es convergente si y sólo si $\limsup x_n = \liminf x_n$, es decir, si posee un único punto de acumulación.
\end{corollary}

\begin{proof}
  ($\Rightarrow$) Si ${(x_n)}$ converge a $L$, toda subsucesión converge a $L$. Como $\liminf x_n$ y $\limsup x_n$ son límites de subsucesiones, ambos son iguales a $L$.

  ($\Leftarrow$) Sea $a = \liminf x_n = \limsup x_n = \lim a_n = \lim b_n$. Dado $\e > 0$, existe $n_0$ tal que $a - \e < a_n$ y $b_n < a + \e$ para $n \geq n_0$. Como $x_n \in X_n$, se tiene $a_n \leq x_n \leq b_n$, y entonces $a - \e < x_n < a + \e$ para todo $n \geq n_0$. Es decir, $x_n \to a$.
\end{proof}

\section{Sucesiones de Cauchy}

Una sucesión ${(x_n)}_{n \in \N} \subset \R$ es \emph{de Cauchy} si para todo $\e > 0$ existe $n_0 \in \N$ tal que $|x_n - x_m| < \e$ para todos $n, m > n_0$, tal como se hizo en $\Q$. Valen, con las mismas demostraciones que en $\Q$:
\begin{enumerate}[label=(\arabic*)]
  \item Toda sucesión convergente es de Cauchy.
  \item Toda sucesión de Cauchy es acotada.
\end{enumerate}

\begin{lemma}
  Si una sucesión de Cauchy tiene una subsucesión que converge a $a \in \R$, entonces $\lim_{n \to \infty} x_n = a$.
\end{lemma}

\begin{proof}
  Dado $\e > 0$, existe $n_0$ tal que $|x_n - x_m| < \frac{\e}{2}$ para $n, m > n_0$. Como $a$ es el límite de una subsucesión, existe un índice $n_1 > n_0$ de la subsucesión con $|x_{n_1} - a| < \frac{\e}{2}$. Entonces, para $n > n_0$,
  \[
    |x_n - a| \leq |x_n - x_{n_1}| + |x_{n_1} - a| < \frac{\e}{2} + \frac{\e}{2} = \e. \qedhere
  \]
\end{proof}

\begin{theorem}
  Toda sucesión de Cauchy en $\R$ es convergente.
\end{theorem}

\begin{proof}
  Por ser de Cauchy es acotada; por el teorema de Bolzano--Weierstrass, alguna subsucesión converge, y por el lema anterior toda la sucesión converge al mismo límite.
\end{proof}

\section{Límites infinitos}

\begin{definition}
  Decimos que $\lim_{n \to \infty} x_n = +\infty$ si para todo $A > 0$ existe $n_0 \in \N$ tal que $x_n > A$ para todo $n > n_0$. Decimos que $\lim_{n \to \infty} x_n = -\infty$ si para todo $A > 0$ existe $n_0 \in \N$ tal que $x_n < -A$ para todo $n > n_0$.
\end{definition}

\begin{prop}
  Si $\lim_{n \to \infty} x_n = +\infty$ e ${(y_n)}$ es acotada inferiormente, entonces $\lim_{n \to \infty} (x_n + y_n) = +\infty$.
\end{prop}

\begin{proof}
  Sea $c \in \R$ con $c < y_n$ para todo $n$. Dado $A > 0$, existe $n_0$ tal que $x_n > A - c$ para $n > n_0$. Entonces $x_n + y_n > (A - c) + c = A$ para $n > n_0$.
\end{proof}

\begin{prop}
  Si $\lim_{n \to \infty} x_n = +\infty$ y existe $c > 0$ tal que $y_n > c$ para todo $n$, entonces $\lim_{n \to \infty} x_n y_n = +\infty$.
\end{prop}

\begin{proof}
  Dado $A > 0$, existe $n_0$ tal que $x_n > \frac{A}{c}$ para $n > n_0$. Entonces $x_n y_n > \frac{A}{c} \cdot c = A$ para $n > n_0$.
\end{proof}

\begin{prop}
  Si $x_n > 0$ para todo $n$, entonces $\lim_{n \to \infty} x_n = 0$ si y sólo si $\lim_{n \to \infty} \frac{1}{x_n} = +\infty$.
\end{prop}

\begin{proof}
  ($\Rightarrow$) Dado $A > 0$, existe $n_0$ tal que $0 < x_n < \frac{1}{A}$ para $n > n_0$, es decir, $\frac{1}{x_n} > A$.

  ($\Leftarrow$) Dado $\e > 0$, existe $n_0$ tal que $\frac{1}{x_n} > \frac{1}{\e}$ para $n > n_0$, es decir, $0 < x_n < \e$.
\end{proof}

\begin{prop}
  Sean ${(x_n)}$ e ${(y_n)}$ sucesiones de términos positivos.
  \begin{enumerate}[label=(\arabic*)]
    \item Si existe $c > 0$ con $x_n > c$ para todo $n$ y $\lim_{n \to \infty} y_n = 0$, entonces $\lim_{n \to \infty} \frac{x_n}{y_n} = +\infty$.
    \item Si ${(x_n)}$ es acotada y $\lim_{n \to \infty} y_n = +\infty$, entonces $\lim_{n \to \infty} \frac{x_n}{y_n} = 0$.
  \end{enumerate}
\end{prop}

\begin{proof}
  (1) Dado $A > 0$, existe $n_0$ tal que $0 < y_n < \frac{c}{A}$ para $n > n_0$, de modo que $\frac{x_n}{y_n} > c \cdot \frac{A}{c} = A$.

  (2) Sea $k > 0$ una cota superior de ${(x_n)}$. Dado $\e > 0$, existe $n_0$ tal que $y_n > \frac{k}{\e}$ para $n > n_0$, de modo que $0 < \frac{x_n}{y_n} < k \cdot \frac{\e}{k} = \e$.
\end{proof}

\section{Series numéricas}

Dada una sucesión ${(a_n)}_{n \in \N} \subset \R$ formamos la sucesión ${(S_n)}_{n \in \N}$ de \emph{sumas parciales} dada por $S_n = a_1 + \cdots + a_n$. La \emph{serie} $\sum_{n \geq 1} a_n$ \emph{converge} si ${(S_n)}$ converge, y su suma es $\lim_{n \to \infty} S_n$.

\begin{theorem}
  Si $\sum_{n \geq 1} a_n$ converge, entonces $\lim_{n \to \infty} a_n = 0$.
\end{theorem}

\begin{proof}
  Sea $s = \lim_{n \to \infty} S_n$. Como $S_{n-1} \to s$ también,
  \[
    a_n = S_n - S_{n-1} \longrightarrow s - s = 0. \qedhere
  \]
\end{proof}

\begin{eg}
  El recíproco del teorema anterior es falso: $\sum_{n \geq 1} \frac{1}{n}$ diverge.
\end{eg}

\begin{proof}
  Agrupando los términos en bloques de $2^{k-1}$ sumandos, cada uno de los cuales es al menos $2^{k-1} \cdot \frac{1}{2^k} = \frac{1}{2}$,
  \begin{align*}
    S_{2^n} & = 1 + \frac{1}{2} + \left( \frac{1}{3} + \frac{1}{4} \right) + \cdots + \left( \frac{1}{2^{n-1} + 1} + \cdots + \frac{1}{2^n} \right) \\
            & \geq 1 + \frac{1}{2} + \frac{2}{4} + \cdots + \frac{2^{n-1}}{2^n} = 1 + \frac{n}{2}.
  \end{align*}
  Luego $S_{2^n} \to +\infty$ y la serie diverge.
\end{proof}

\begin{theorem}
  Si $a_n \geq 0$ para todo $n$, entonces $\sum_{n \geq 1} a_n$ converge si y sólo si la sucesión ${(S_n)}$ es acotada.
\end{theorem}

\begin{proof}
  Por ser $a_n \geq 0$, ${(S_n)}$ es creciente, y por lo tanto converge si y sólo si es acotada.
\end{proof}

\begin{theorem}[Criterio de Cauchy]
  $\sum a_n$ converge si y sólo si para todo $\e > 0$ existe $n_0 \in \N$ tal que
  \[
    |a_{n+1} + \cdots + a_{n+p}| < \e \quad \text{para todo } n > n_0 \text{ y todo } p \in \N.
  \]
\end{theorem}

\begin{proof}
  Como $|S_{n+p} - S_n| = |a_{n+1} + \cdots + a_{n+p}|$, la condición dice exactamente que ${(S_n)}$ es de Cauchy, lo cual equivale a que sea convergente.
\end{proof}

\begin{theorem}[Criterio de comparación]
  Sean $a_n, b_n \geq 0$ para todo $n$ y supongamos que existen $c > 0$ y $n_0 \in \N$ con $a_n \leq c\, b_n$ para todo $n > n_0$.
  \begin{enumerate}[label=(\alph*)]
    \item Si $\sum b_n$ converge, entonces $\sum a_n$ converge.
    \item Si $\sum a_n$ diverge, entonces $\sum b_n$ diverge.
  \end{enumerate}
\end{theorem}

\begin{proof}
  Notemos que la convergencia de una serie sólo depende de su cola: $\sum a_n$ converge si y sólo si $\sum_{n > n_1} a_n$ converge para algún (todo) $n_1$.

  (a) Si $\sum b_n$ converge, sus sumas parciales $T_n$ están acotadas, digamos por $K$. Para $n > n_0$ las sumas parciales de $\sum a_n$ cumplen
  \[
    S_n = S_{n_0} + \sum_{m = n_0 + 1}^{n} a_m \leq S_{n_0} + c \sum_{m = n_0 + 1}^{n} b_m \leq S_{n_0} + c K.
  \]
  Luego ${(S_n)}$ es acotada y $\sum a_n$ converge.

  (b) Es la contrarrecíproca de (a).
\end{proof}

\begin{definition}
  La serie $\sum a_n$ es \emph{absolutamente convergente} si $\sum |a_n|$ converge. Si $\sum a_n$ converge pero $\sum |a_n|$ diverge, se dice que es \emph{condicionalmente convergente}.
\end{definition}

\begin{theorem}
  Toda serie absolutamente convergente es convergente.
\end{theorem}

\begin{proof}
  Dado $\e > 0$, como $\sum |a_n|$ converge, por el criterio de Cauchy existe $n_0$ tal que $|a_{n+1}| + \cdots + |a_{n+p}| < \e$ para $n > n_0$ y $p \in \N$. Entonces
  \[
    |a_{n+1} + \cdots + a_{n+p}| \leq |a_{n+1}| + \cdots + |a_{n+p}| < \e
  \]
  para $n > n_0$ y $p \in \N$, y por el criterio de Cauchy, $\sum a_n$ converge.
\end{proof}

\begin{theorem}[Criterio de condensación de Cauchy]
  Sea ${(a_n)}$ una sucesión con $a_1 \geq a_2 \geq \cdots \geq 0$. Entonces $\sum_{n \geq 1} a_n$ converge si y sólo si $\sum_{k \geq 0} 2^k a_{2^k}$ converge.
\end{theorem}

\begin{proof}
  Sean $s_n = a_1 + \cdots + a_n$ y $t_k = a_1 + 2 a_2 + \cdots + 2^k a_{2^k}$.

  Si $n < 2^k$, agrupando en bloques y usando que la sucesión es decreciente,
  \begin{align*}
    s_n & \leq a_1 + (a_2 + a_3) + \cdots + (a_{2^{k-1}} + \cdots + a_{2^k - 1}) \\
        & \leq a_1 + 2 a_2 + \cdots + 2^{k-1} a_{2^{k-1}} \leq t_k.
  \end{align*}
  Por otro lado, si $n \geq 2^k$,
  \begin{align*}
    s_n & \geq a_1 + a_2 + (a_3 + a_4) + \cdots + (a_{2^{k-1} + 1} + \cdots + a_{2^k}) \\
        & \geq a_1 + a_2 + 2 a_4 + \cdots + 2^{k-1} a_{2^k} \geq \frac{1}{2} t_k.
  \end{align*}
  Luego las sucesiones ${(s_n)}$ y ${(t_k)}$ (crecientes) son ambas acotadas o ambas no acotadas, y el resultado se sigue del teorema sobre series de términos no negativos.
\end{proof}


\clase{24/09/2024}{Convergencia de series}
\section{Un ejemplo de convergencia condicional}

\begin{eg}
  La serie \(\sum_{n \geq 1} \dfrac{(-1)^{n+1}}{n}\) converge, pero no lo hace absolutamente.
\end{eg}

\begin{proof}
  No converge absolutamente porque \(\sum \frac{1}{n}\) diverge. Veamos que converge.

  \emph{(i) Sumas parciales pares.} Se tiene
  \[
    S_{2n} = \left( 1 - \frac{1}{2} \right) + \left( \frac{1}{3} - \frac{1}{4} \right) + \cdots + \left( \frac{1}{2n-1} - \frac{1}{2n} \right),
  \]
  de modo que \(S_2 < S_4 < \cdots < S_{2n} < \cdots\) es creciente. Además,
  \[
    S_{2n} = 1 - \left( \frac{1}{2} - \frac{1}{3} \right) - \cdots - \left( \frac{1}{2n-2} - \frac{1}{2n-1} \right) - \frac{1}{2n} < 1,
  \]
  luego \({(S_{2n})}\) es creciente y acotada, y por lo tanto converge a algún \(\tilde{s}\).

  \emph{(ii) Sumas parciales impares.} Como \(S_{2n+1} = S_{2n} + \frac{1}{2n+1}\) y
  \[
    S_{2n+1} = 1 - \left( \frac{1}{2} - \frac{1}{3} \right) - \cdots - \left( \frac{1}{2n} - \frac{1}{2n+1} \right),
  \]
  se tiene \(S_1 > S_3 > \cdots > S_{2n+1} > \cdots > 0\), es decir, \({(S_{2n+1})}\) es decreciente y acotada, y converge a algún \(s'\).

  Como \(S_{2n+1} - S_{2n} = \frac{1}{2n+1} \to 0\), resulta \(\tilde{s} = s'\). Entonces \({(S_n)}\) converge, pues sus subsucesiones de índices pares e impares convergen al mismo límite.
\end{proof}

\section{Corolarios del criterio de comparación}

\begin{corollary}
  Sea \(\sum b_n\) una serie convergente de términos positivos. Si existen \(k > 0\) y \(n_0 \in \N\) tales que \(|a_n| \leq k\, b_n\) para todo \(n > n_0\), entonces \(\sum a_n\) converge absolutamente.
\end{corollary}

\begin{proof}
  Es una aplicación directa del criterio de comparación a \(|a_n|\) y \(b_n\).
\end{proof}

\begin{corollary}
  Si existen \(k > 0\), \(0 < c < 1\) y \(n_0 \in \N\) tales que \(|a_n| \leq k\, c^n\) para todo \(n > n_0\), entonces \(\sum a_n\) converge absolutamente.
\end{corollary}

\begin{proof}
  Como \(0 < c < 1\), la serie geométrica \(\sum c^n\) converge, y se aplica el corolario anterior.
\end{proof}

\begin{note}
  Tomando \(k = 1\), la condición \(|a_n| \leq c^n\) equivale a \(\sqrt[n]{|a_n|} \leq c < 1\). Si esto vale para todo \(n > n_0\), entonces \(\limsup_{n \to \infty} \sqrt[n]{|a_n|} \leq c < 1\). Recíprocamente, si \(\limsup_{n \to \infty} \sqrt[n]{|a_n|} < 1\), existen \(c < 1\) y \(n_0\) con \(\sqrt[n]{|a_n|} < c\) para \(n > n_0\).
\end{note}

\section{Criterios de convergencia}

\begin{theorem}[Criterio de la raíz]
  Sea \(\sum a_n\) una serie y \(\alpha = \limsup_{n \to \infty} \sqrt[n]{|a_n|}\).
  \begin{enumerate}[label=(\arabic*)]
    \item Si \(\alpha < 1\), entonces \(\sum a_n\) converge absolutamente.
    \item Si \(\alpha > 1\), entonces \(\sum a_n\) diverge.
    \item Si \(\alpha = 1\), el criterio no decide: tanto \(\sum \frac{1}{n}\) (que diverge) como \(\sum \frac{1}{n^2}\) (que converge) cumplen \(\sqrt[n]{|a_n|} \to 1\).
  \end{enumerate}
\end{theorem}

\begin{proof}
  (1) Sea \(\beta \in (\alpha, 1)\). Como \(\alpha = \inf_n \sup_{k \geq n} \sqrt[k]{|a_k|}\), existe \(N\) tal que \(\sqrt[n]{|a_n|} < \beta\) para todo \(n > N\), es decir, \(|a_n| < \beta^n\). Como \(0 < \beta < 1\), la serie \(\sum \beta^n\) converge y, por el criterio de comparación, \(\sum a_n\) converge absolutamente.

  (2) Si \(\alpha > 1\), existen infinitos índices con \(\sqrt[n]{|a_n|} > 1\) (si la sucesión de raíces no es acotada, también los hay), y entonces \(|a_n| > 1\) para infinitos \(n\). Luego \(a_n \not\to 0\) y la serie diverge.

  (3) Se tiene \(\sqrt[n]{1/n} \to 1\) por un ejemplo anterior y, por la misma razón, \(\sqrt[n]{1/n^2} = {(\sqrt[n]{1/n})}^2 \to 1\). La serie \(\sum \frac{1}{n}\) diverge, mientras que \(\sum \frac{1}{n^2}\) converge (por ejemplo, por el criterio de condensación, pues \(\sum 2^k \frac{1}{4^k} = \sum \frac{1}{2^k}\) converge).
\end{proof}

\begin{theorem}[Criterio de comparación de cocientes]
  Sea \(\sum a_n\) una serie con \(a_n \neq 0\) para todo \(n\), y sea \(\sum b_n\) una serie convergente de términos positivos. Si existe \(n_0 \in \N\) tal que
  \[
    \frac{|a_{n+1}|}{|a_n|} \leq \frac{b_{n+1}}{b_n} \quad \text{para todo } n > n_0,
  \]
  entonces \(\sum a_n\) converge absolutamente.
\end{theorem}

\begin{proof}
  Multiplicando las desigualdades correspondientes a \(n_0 + 1, n_0 + 2, \dots, n - 1\) se obtiene, para \(n > n_0 + 1\),
  \[
    \frac{|a_n|}{|a_{n_0 + 1}|} \leq \frac{b_n}{b_{n_0 + 1}}, \qquad \text{es decir,} \qquad |a_n| \leq \frac{|a_{n_0 + 1}|}{b_{n_0 + 1}} \, b_n.
  \]
  Por el criterio de comparación, \(\sum |a_n|\) converge.
\end{proof}

\begin{corollary}[Criterio de D'Alembert]
  Si existe \(c \in (0, 1)\) tal que \(\dfrac{|a_{n+1}|}{|a_n|} \leq c\) para todo \(n > n_0\), entonces \(\sum a_n\) converge absolutamente. En particular, esto ocurre si \(\limsup \dfrac{|a_{n+1}|}{|a_n|} < 1\). Si el cociente vale \(1\) el criterio no decide, y si \(\dfrac{|a_{n+1}|}{|a_n|} \geq 1\) para todo \(n > n_0\) la serie diverge (pues \(a_n \not\to 0\)).
\end{corollary}

\begin{proof}
  Se toma \(b_n = c^n\) en el teorema anterior, pues \(\sum c^n\) converge si \(c \in (0, 1)\). Si \(\limsup \frac{|a_{n+1}|}{|a_n|} < 1\), se elige \(c\) entre ese valor y \(1\) como en el criterio de la raíz.
\end{proof}

\begin{theorem}
  Sea \({(a_n)}\) una sucesión de términos positivos tal que \({\left( \frac{a_{n+1}}{a_n} \right)}\) es acotada. Entonces
  \[
    \liminf_{n \to \infty} \frac{a_{n+1}}{a_n} \leq \liminf_{n \to \infty} \sqrt[n]{a_n} \leq \limsup_{n \to \infty} \sqrt[n]{a_n} \leq \limsup_{n \to \infty} \frac{a_{n+1}}{a_n}.
  \]
  En particular, si existe \(\lim_{n \to \infty} \frac{a_{n+1}}{a_n}\), entonces existe \(\lim_{n \to \infty} \sqrt[n]{a_n}\) y ambos son iguales.
\end{theorem}

\begin{proof}
  La desigualdad del medio es evidente. Veamos que \(\limsup \sqrt[n]{a_n} \leq \limsup \frac{a_{n+1}}{a_n}\). Supongamos que no es así; entonces existe \(c\) con
  \[
    \limsup_{n \to \infty} \frac{a_{n+1}}{a_n} < c < \limsup_{n \to \infty} \sqrt[n]{a_n}.
  \]
  Por la definición de \(\limsup\), existe \(n_0\) tal que \(\frac{a_{n+1}}{a_n} < c\) para todo \(n \geq n_0\). Multiplicando estas desigualdades para \(n_0, \dots, n - 1\),
  \[
    \frac{a_n}{a_{n_0}} < c^{n - n_0} \quad \Longrightarrow \quad a_n < k\, c^n, \qquad k = \frac{a_{n_0}}{c^{n_0}}.
  \]
  Como \(\sqrt[n]{k} \to 1\), resulta
  \[
    \limsup_{n \to \infty} \sqrt[n]{a_n} \leq \limsup_{n \to \infty} \sqrt[n]{k\, c^n} = c \cdot \lim_{n \to \infty} \sqrt[n]{k} = c,
  \]
  lo cual es absurdo. La primera desigualdad se prueba de manera análoga, invirtiendo las desigualdades.

  Si el límite del cociente existe, \(\liminf\) y \(\limsup\) del cociente coinciden, y entonces los de la raíz también coinciden con él, de modo que la sucesión de raíces converge a ese valor.
\end{proof}

\begin{eg}
  Puede existir el límite de las raíces sin que exista el del cociente.
\end{eg}

\begin{proof}
  Sean \(0 < a < b\) y sea \({(x_n)}\) la sucesión que se obtiene multiplicando alternadamente por \(a\) y por \(b\):
  \[
    x_1 = a, \quad x_2 = a b, \quad x_3 = a^2 b, \quad x_4 = a^2 b^2, \quad \dots
  \]
  Entonces
  \[
    \frac{x_{n+1}}{x_n} =
    \begin{cases}
      b & \text{si } n \text{ es impar}, \\
      a & \text{si } n \text{ es par},
    \end{cases}
  \]
  de modo que no existe \(\lim_{n \to \infty} \frac{x_{n+1}}{x_n}\) (su \(\limsup\) es \(b\) y su \(\liminf\) es \(a\)). Por otro lado, \(x_{2k} = a^k b^k\) y \(x_{2k-1} = a^k b^{k-1}\), de modo que
  \[
    \sqrt[2k]{x_{2k}} = \sqrt{ab}, \qquad \sqrt[2k-1]{x_{2k-1}} = a^{\frac{k}{2k-1}} b^{\frac{k-1}{2k-1}} = \sqrt{ab} \cdot a^{\frac{1}{2(2k-1)}} b^{-\frac{1}{2(2k-1)}} \longrightarrow \sqrt{ab}.
  \]
  Como ambas subsucesiones convergen a \(\sqrt{ab}\), se tiene \(\lim_{n \to \infty} \sqrt[n]{x_n} = \sqrt{ab}\).
\end{proof}

\begin{theorem}[Criterio de Dirichlet]
  Sea \(\sum a_n\) una serie (no necesariamente convergente) cuyas sumas parciales \(S_n\) forman una sucesión acotada, y sea \({(b_n)}\) una sucesión decreciente de números positivos con \(\lim_{n \to \infty} b_n = 0\). Entonces \(\sum a_n b_n\) converge.
\end{theorem}

\begin{proof}
  Por sumación por partes, con \(S_0 = 0\),
  \begin{align*}
    a_1 b_1 + \cdots + a_n b_n & = S_1 (b_1 - b_2) + S_2 (b_2 - b_3) + \cdots + S_{n-1} (b_{n-1} - b_n) + S_n b_n \\
                               & = \sum_{i=1}^{n-1} S_i (b_i - b_{i+1}) + S_n b_n.
  \end{align*}
  Por hipótesis existe \(k > 0\) con \(|S_n| \leq k\) para todo \(n\). Como \(b_i - b_{i+1} \geq 0\), y la serie telescópica \(\sum_{i \geq 1} (b_i - b_{i+1})\) converge (a \(b_1\)), la serie \(\sum S_i (b_i - b_{i+1})\) converge absolutamente por comparación, y en particular converge. Además \(S_n b_n \to 0\), por ser \((S_n)\) acotada y \(b_n \to 0\). Luego la sucesión de sumas parciales de \(\sum a_n b_n\) converge.
\end{proof}

\begin{corollary}[Criterio de Abel]
  Si \(\sum a_n\) converge y \({(b_n)}\) es una sucesión decreciente de términos positivos, entonces \(\sum a_n b_n\) converge.
\end{corollary}

\begin{proof}
  Como \({(b_n)}\) es decreciente y acotada inferiormente por \(0\), converge a algún \(c \geq 0\). La sucesión \({(b_n - c)}\) es decreciente, de términos no negativos, y tiende a \(0\). Como \(\sum a_n\) converge, sus sumas parciales son acotadas, y por el criterio de Dirichlet \(\sum a_n (b_n - c)\) converge. Entonces
  \[
    \sum a_n b_n = \sum a_n (b_n - c) + c \sum a_n
  \]
  es suma de dos series convergentes.
\end{proof}

\begin{corollary}[Criterio de Leibniz]
  Si \({(b_n)}\) es decreciente y \(b_n \to 0\), entonces \(\sum (-1)^n b_n\) converge.
\end{corollary}

\begin{proof}
  La serie \(\sum (-1)^n\) no converge, pero sus sumas parciales toman los valores \(-1\) y \(0\), y son acotadas. Se aplica el criterio de Dirichlet.
\end{proof}

\section{Parte positiva y parte negativa}

\begin{definition}
  Dada una serie \(\sum a_n\), para cada \(n \in \N\) definimos
  \[
    p_n =
    \begin{cases}
      a_n & \text{si } a_n > 0,    \\
      0   & \text{en otro caso},
    \end{cases}
    \qquad
    q_n =
    \begin{cases}
      -a_n & \text{si } a_n < 0,    \\
      0    & \text{en otro caso}.
    \end{cases}
  \]
  Entonces \(p_n, q_n \geq 0\), \(a_n = p_n - q_n\) y \(|a_n| = p_n + q_n = a_n + 2 q_n = 2 p_n - a_n\).
\end{definition}

\begin{note}
  Si \(\sum a_n\) converge absolutamente, para todo \(k\) se tiene
  \[
    \sum_{n \geq 1} |a_n| \geq \sum_{n=1}^{k} |a_n| = \sum_{n=1}^{k} p_n + \sum_{n=1}^{k} q_n,
  \]
  de modo que las sumas parciales de \(\sum p_n\) y \(\sum q_n\) son crecientes y acotadas por \(\sum |a_n|\), y ambas series convergen. Recíprocamente, si \(\sum p_n\) y \(\sum q_n\) convergen, entonces \(\sum |a_n| = \sum (p_n + q_n)\) converge, y \(\sum a_n\) converge absolutamente.
\end{note}

\begin{theorem}
  Si \(\sum a_n\) converge condicionalmente, entonces \(\sum p_n\) y \(\sum q_n\) divergen.
\end{theorem}

\begin{proof}
  Supongamos que \(\sum q_n\) converge, con suma \(c\). Como \(|a_n| = a_n + 2 q_n\),
  \[
    \sum_{n=1}^{k} |a_n| = \sum_{n=1}^{k} a_n + 2 \sum_{n=1}^{k} q_n \longrightarrow \sum_{n \geq 1} a_n + 2c \quad (k \to \infty),
  \]
  de modo que \(\sum |a_n|\) converge, lo cual es absurdo pues \(\sum a_n\) no converge absolutamente. Luego \(\sum q_n\) diverge. Análogamente, usando \(|a_n| = 2 p_n - a_n\), se prueba que \(\sum p_n\) diverge.
\end{proof}

\section{Reordenamientos}

\begin{definition}
  Reordenar una serie \(\sum a_n\) significa tomar una biyección \(\varphi\colon \N \to \N\) y considerar la serie \(\sum b_n\) con \(b_n = a_{\varphi(n)}\).
\end{definition}

\begin{theorem}
  Todo reordenamiento de una serie absolutamente convergente converge al mismo valor que la serie original.
\end{theorem}

\begin{proof}
  \emph{Caso \(a_n \geq 0\).} Sea \(\varphi\colon \N \to \N\) una biyección y \(b_n = a_{\varphi(n)}\). Sean \(S_n\) y \(T_n\) las sumas parciales de \(\sum a_n\) y \(\sum b_n\). Para cada \(n\), sea \(m = \max\{\varphi(1), \dots, \varphi(n)\}\), de modo que \(\{\varphi(1), \dots, \varphi(n)\} \subset \{1, \dots, m\}\). Como \(a_j \geq 0\),
  \[
    T_n = \sum_{i=1}^{n} a_{\varphi(i)} \leq \sum_{j=1}^{m} a_j = S_m \leq \sum_{j \geq 1} a_j.
  \]
  Entonces \(\sum b_n\) converge y \(\sum b_n \leq \sum a_n\). Como \(a_n = b_{\varphi^{-1}(n)}\) es un reordenamiento de \(\sum b_n\), el mismo argumento da \(\sum a_n \leq \sum b_n\). Luego son iguales.

  \emph{Caso general.} Como \(\sum a_n\) converge absolutamente, \(\sum a_n = \sum p_n - \sum q_n\) con \(\sum p_n\) y \(\sum q_n\) convergentes y de términos no negativos. Las partes positiva y negativa de \(b_n = a_{\varphi(n)}\) son \(p_{\varphi(n)}\) y \(q_{\varphi(n)}\), que son reordenamientos de \(\sum p_n\) y \(\sum q_n\). Por el caso anterior,
  \[
    \sum b_n = \sum p_{\varphi(n)} - \sum q_{\varphi(n)} = \sum p_n - \sum q_n = \sum a_n. \qedhere
  \]
\end{proof}

\begin{theorem}[Riemann]
  Sea \(\sum a_n\) una serie condicionalmente convergente y sea \(c \in \R\). Entonces existe un reordenamiento \(\sum b_n\) de \(\sum a_n\) tal que \(\sum b_n = c\).
\end{theorem}

\begin{proof}
  Como \(\sum a_n\) converge, \(a_n \to 0\). Sean \(P_1, P_2, \dots\) los términos positivos de la serie y \(-Q_1, -Q_2, \dots\) los negativos, cada uno en su orden original. Entonces \(P_k \to 0\) y \(Q_k \to 0\), y por el teorema anterior \(\sum P_k = \sum Q_k = +\infty\). (Los términos nulos, si los hay, pueden intercalarse en cualquier lugar sin alterar el argumento.)

  Construimos el reordenamiento por etapas:
  \begin{enumerate}[label=(\arabic*)]
    \item Se toman \(P_1, P_2, \dots\) hasta que la suma parcial supere por primera vez a \(c\).
    \item Se agregan \(-Q_1, -Q_2, \dots\) hasta que la suma parcial sea por primera vez menor que \(c\).
    \item Se agregan los siguientes \(P_k\) hasta superar nuevamente a \(c\), luego los siguientes \(-Q_k\) hasta quedar por debajo de \(c\), y así sucesivamente.
  \end{enumerate}
  El proceso nunca se detiene, pues \(\sum P_k\) y \(\sum Q_k\) divergen, y usa todos los términos de la serie original, ya que en cada etapa se consumen términos nuevos de cada tipo.

  Veamos que las sumas parciales \(T_N\) convergen a \(c\). Al terminar una etapa de suma de positivos, con último término \(P_j\), la suma parcial supera a \(c\) pero antes de sumar \(P_j\) no lo hacía, de modo que \(0 < T - c \leq P_j\). Análogamente, al terminar una etapa de negativos con último término \(-Q_j\), \(0 < c - T \leq Q_j\). Dentro de cada etapa, las sumas parciales varían monótonamente entre dos de estos valores, de modo que están entre ambos. Entonces, para \(N\) en una etapa avanzada, \(|T_N - c| \leq \max\{P_j, Q_{j'}\}\), donde \(P_j\) y \(Q_{j'}\) son los últimos términos usados en los extremos de dicha etapa. Como \(P_j \to 0\) y \(Q_{j'} \to 0\) cuando el número de etapas tiende a infinito, resulta \(T_N \to c\).
\end{proof}


\clase{27/09/2024}{Topología euclídea}
\section{Topología euclídea}

\begin{definition}[Norma euclídea]
  La \emph{norma euclídea} en $\R^n$ es la función $\norm{\cdot}\colon \R^n \to \R$ dada por
  \[
    \norm{x} = \sqrt{\sum_{i=1}^{n} x_i^2}, \qquad x = (x_1, \dots, x_n).
  \]
\end{definition}

Se cumplen las siguientes propiedades, para todo $x, y \in \R^n$ y todo $\lambda \in \R$:
\begin{enumerate}[label=(\arabic*)]
  \item $\norm{x} \geq 0$, y $\norm{x} = 0 \iff x = 0$.
  \item $\norm{\lambda x} = |\lambda| \, \norm{x}$.
  \item $\norm{x + y} \leq \norm{x} + \norm{y}$ (desigualdad triangular).
\end{enumerate}
Las dos primeras son inmediatas; la tercera se deduce de la siguiente desigualdad.

\begin{prop}[Desigualdad de Cauchy--Schwarz]
  Dados $x, y \in \R^n$,
  \[
    \left| \sum_{i=1}^{n} x_i y_i \right| \leq \norm{x} \, \norm{y}.
  \]
\end{prop}

\begin{proof}
  Se verifica la identidad de Lagrange
  \[
    \left( \sum_{k=1}^{n} x_k^2 \right) \left( \sum_{k=1}^{n} y_k^2 \right) - {\left( \sum_{k=1}^{n} x_k y_k \right)}^2 = \sum_{1 \leq k < j \leq n} {(x_k y_j - x_j y_k)}^2.
  \]
  En efecto, el lado izquierdo es
  \[
    \sum_{k, j} x_k^2 y_j^2 - \sum_{k, j} x_k y_k x_j y_j = \frac{1}{2} \sum_{k, j} \bigl( x_k^2 y_j^2 + x_j^2 y_k^2 - 2 x_k y_k x_j y_j \bigr) = \frac{1}{2} \sum_{k, j} {(x_k y_j - x_j y_k)}^2,
  \]
  y los términos con $k = j$ se anulan, mientras que los pares $(k, j)$ y $(j, k)$ aportan lo mismo. Como el lado derecho es una suma de cuadrados, es no negativo, de modo que ${(\sum x_k y_k)}^2 \leq \norm{x}^2 \norm{y}^2$.
\end{proof}

\begin{proof}[Demostración de la desigualdad triangular]
  Por Cauchy--Schwarz,
  \[
    \norm{x + y}^2 = \norm{x}^2 + 2 \sum_i x_i y_i + \norm{y}^2 \leq \norm{x}^2 + 2 \norm{x} \norm{y} + \norm{y}^2 = {\bigl( \norm{x} + \norm{y} \bigr)}^2. \qedhere
  \]
\end{proof}

\begin{definition}[Distancia euclídea]
  La \emph{distancia euclídea} es la función $d\colon \R^n \times \R^n \to \R$, $d(x, y) = \norm{x - y}$.
\end{definition}

\begin{prop}
  Para todos $x, y, z \in \R^n$:
  \begin{enumerate}[label=(\arabic*)]
    \item $d(x, y) \geq 0$, y $d(x, y) = 0 \iff x = y$.
    \item $d(x, y) = d(y, x)$.
    \item $d(x, z) \leq d(x, y) + d(y, z)$.
  \end{enumerate}
\end{prop}

\begin{proof}
  (1) Es la primera propiedad de la norma. (2) $d(x, y) = \norm{x - y} = \norm{(-1)(y - x)} = \norm{y - x} = d(y, x)$. (3) $d(x, z) = \norm{(x - y) + (y - z)} \leq \norm{x - y} + \norm{y - z} = d(x, y) + d(y, z)$.
\end{proof}

\section{Conjuntos abiertos}

\begin{definition}[Bola abierta]
  Si $x \in \R^n$ y $r > 0$, la \emph{bola abierta} de centro $x$ y radio $r$ es el conjunto
  \[
    B_r(x) = B(x, r) = \{ y \in \R^n : d(x, y) < r \}.
  \]
\end{definition}

\begin{definition}[Conjunto abierto]
  Un subconjunto $U \subset \R^n$ es \emph{abierto} si para todo $x \in U$ existe $r > 0$ tal que $B_r(x) \subset U$.
\end{definition}

\begin{lemma}
  La bola abierta $B_r(x)$ es un conjunto abierto.
\end{lemma}

\begin{proof}
  Sea $y \in B_r(x)$ y $s = r - d(x, y) > 0$. Si $z \in B_s(y)$, entonces
  \[
    d(x, z) \leq d(x, y) + d(y, z) < d(x, y) + s = r,
  \]
  es decir, $z \in B_r(x)$. Luego $B_s(y) \subset B_r(x)$.
\end{proof}

\begin{prop}
  Un subconjunto $U \subsetneq \R^n$ es abierto si y sólo si para todo $x \in U$ se tiene $\inf \{ d(x, y) : y \in \R^n \setminus U \} > 0$. El ínfimo existe pues el conjunto es no vacío y acotado inferiormente por $0$.
\end{prop}

\begin{proof}
  ($\Rightarrow$) Sea $x \in U$ y $r > 0$ con $B_r(x) \subset U$. Si $y \in \R^n \setminus U$, entonces $y \notin B_r(x)$, es decir, $d(x, y) \geq r$. Luego $r$ es cota inferior del conjunto y el ínfimo es mayor o igual que $r > 0$.

  ($\Leftarrow$) Sea $x \in U$ y $s = \inf \{ d(x, y) : y \in \R^n \setminus U \} > 0$. Si $z \in B_s(x)$ y fuera $z \notin U$, sería $d(x, z) \geq s$, lo que contradice $d(x, z) < s$. Luego $B_s(x) \subset U$.
\end{proof}

\begin{prop}
  Valen las siguientes propiedades:
  \begin{enumerate}[label=(\arabic*)]
    \item $\varnothing$ y $\R^n$ son abiertos.
    \item La intersección de una familia finita y no vacía de abiertos de $\R^n$ es abierta.
    \item La unión de una familia cualquiera de abiertos de $\R^n$ es abierta.
  \end{enumerate}
\end{prop}

\begin{proof}
  (1) Es inmediato de la definición.

  (2) Sean $U_1, \dots, U_k$ abiertos y $x \in U = \bigcap_{i=1}^{k} U_i$. Para cada $i$ existe $r_i > 0$ con $B_{r_i}(x) \subset U_i$. Con $r = \min\{r_1, \dots, r_k\} > 0$ se tiene $B_r(x) \subset U_i$ para todo $i$, es decir, $B_r(x) \subset U$.

  (3) Sea $\mathscr{U}$ una familia de abiertos y $U$ su unión. Dado $x \in U$, existe $V \in \mathscr{U}$ con $x \in V$. Como $V$ es abierto, existe $r > 0$ con $B_r(x) \subset V \subset U$.
\end{proof}

La hipótesis de que la familia sea finita en (2) es necesaria: si $U_n = \left( 1 - \frac{1}{n}, 1 + \frac{1}{n} \right)$ para $n \in \N$, entonces $\bigcap_{n \geq 1} U_n = \{1\}$, que no es abierto.

\begin{definition}[Topología euclídea]
  La \emph{topología euclídea} de $\R^n$ es el conjunto de todos los abiertos de $\R^n$.
\end{definition}

\begin{corollary}
  Un subconjunto de $\R^n$ es abierto si y sólo si es unión de una familia de bolas abiertas.
\end{corollary}

\begin{proof}
  ($\Rightarrow$) Sea $U$ abierto y sea $\mathscr{B}$ la familia de bolas abiertas contenidas en $U$, con unión $\tilde{U}$. Claramente $\tilde{U} \subset U$. Recíprocamente, dado $x \in U$, existe $r > 0$ con $B_r(x) \subset U$, es decir, $B_r(x) \in \mathscr{B}$, y entonces $x \in B_r(x) \subset \tilde{U}$. Luego $U = \tilde{U}$.

  ($\Leftarrow$) Las bolas abiertas son abiertas, y toda unión de abiertos es abierta.
\end{proof}

\section{Interior de un conjunto}

\begin{definition}[Interior]
  Sea $X \subset \R^n$. Un punto $x \in \R^n$ es \emph{interior} a $X$ si existe $r > 0$ con $B_r(x) \subset X$. El \emph{interior} de $X$, que notamos $X^{\circ}$, es el conjunto de los puntos interiores de $X$.
\end{definition}

\begin{note}
  Siempre vale $X^{\circ} \subset X$, y $X$ es abierto si y sólo si $X^{\circ} = X$.
\end{note}

\begin{prop}
  Sea $X \subset \R^n$. Entonces $X^{\circ}$ es la unión de todos los subconjuntos abiertos de $X$; en particular es abierto, y es el mayor abierto contenido en $X$.
\end{prop}

\begin{proof}
  Sea $U$ la unión de todos los abiertos contenidos en $X$. Veamos que $U = X^{\circ}$.

  ($X^{\circ} \subset U$) Si $x \in X^{\circ}$, existe $r > 0$ con $B_r(x) \subset X$; como $B_r(x)$ es abierto, $B_r(x) \subset U$, y entonces $x \in U$.

  ($U \subset X^{\circ}$) Si $x \in U$, existe un abierto $V \subset X$ con $x \in V$, y existe $r > 0$ con $B_r(x) \subset V \subset X$. Luego $x \in X^{\circ}$.

  Como $U$ es unión de abiertos, es abierto, y contiene a todo abierto contenido en $X$.
\end{proof}

\begin{corollary}
  Si $X \subset Y$, entonces $X^{\circ} \subset Y^{\circ}$.
\end{corollary}

\begin{proof}
  $X^{\circ}$ es un abierto contenido en $X \subset Y$, y $Y^{\circ}$ es el mayor abierto contenido en $Y$.
\end{proof}

\subsection*{Propiedades del interior}

\begin{prop}
  ${(X \cap Y)}^{\circ} = X^{\circ} \cap Y^{\circ}$.
\end{prop}

\begin{proof}
  Como $X^{\circ} \cap Y^{\circ}$ es abierto (intersección finita de abiertos) y está contenido en $X \cap Y$, se tiene $X^{\circ} \cap Y^{\circ} \subset {(X \cap Y)}^{\circ}$. Por otro lado, $X \cap Y \subset X$ implica ${(X \cap Y)}^{\circ} \subset X^{\circ}$, y análogamente ${(X \cap Y)}^{\circ} \subset Y^{\circ}$. Luego ${(X \cap Y)}^{\circ} \subset X^{\circ} \cap Y^{\circ}$.
\end{proof}

\begin{prop}
  $X^{\circ} \cup Y^{\circ} \subset {(X \cup Y)}^{\circ}$.
\end{prop}

\begin{proof}
  $X^{\circ} \cup Y^{\circ}$ es abierto y está contenido en $X \cup Y$, de modo que está contenido en el mayor abierto contenido en $X \cup Y$.
\end{proof}

\begin{eg}
  La inclusión anterior puede ser estricta. Sean $X = (-1, 1)$ e $Y = [1, 2]$. Entonces $X^{\circ} \cup Y^{\circ} = (-1, 1) \cup (1, 2)$ no contiene a $1$, pero $X \cup Y = (-1, 2]$ tiene interior $(-1, 2) \ni 1$.
\end{eg}

\section{Entornos}

\begin{definition}[Entorno]
  Sea $x \in \R^n$. Un \emph{entorno} de $x$ es un subconjunto $N \subset \R^n$ tal que $x \in N^{\circ}$. Un \emph{entorno abierto} de $x$ es un entorno que es abierto, es decir, un abierto que contiene a $x$.
\end{definition}

\begin{prop}
  Un subconjunto de $\R^n$ es abierto si y sólo si es entorno de cada uno de sus puntos.
\end{prop}

\begin{proof}
  Si $U$ es abierto, $U = U^{\circ}$ y entonces $x \in U^{\circ}$ para todo $x \in U$. Recíprocamente, si $U$ es entorno de cada uno de sus puntos, $U \subset U^{\circ} \subset U$, de modo que $U = U^{\circ}$ y $U$ es abierto.
\end{proof}

\begin{prop}
  El interior de $X \subset \R^n$ es el conjunto de puntos de los que $X$ es entorno.
\end{prop}

\begin{proof}
  Es la definición: $X$ es entorno de $x$ si y sólo si $x \in X^{\circ}$.
\end{proof}

\section{Conjuntos cerrados}

\begin{definition}[Conjunto cerrado]
  Un subconjunto $F \subset \R^n$ es \emph{cerrado} si su complemento $\R^n \setminus F$ es abierto.
\end{definition}

\begin{prop}
  Valen las siguientes propiedades:
  \begin{enumerate}[label=(\arabic*)]
    \item $\varnothing$ y $\R^n$ son cerrados.
    \item La unión de una familia finita de cerrados es cerrada.
    \item La intersección de una familia arbitraria de cerrados es cerrada.
  \end{enumerate}
\end{prop}

\begin{proof}
  (1) $\R^n \setminus \R^n = \varnothing$ y $\R^n \setminus \varnothing = \R^n$ son abiertos.

  (2) Si $F_1, \dots, F_k$ son cerrados, $\R^n \setminus (F_1 \cup \cdots \cup F_k) = (\R^n \setminus F_1) \cap \cdots \cap (\R^n \setminus F_k)$ es una intersección finita de abiertos.

  (3) Si $\mathscr{F}$ es una familia de cerrados, $\R^n \setminus \bigcap_{F \in \mathscr{F}} F = \bigcup_{F \in \mathscr{F}} (\R^n \setminus F)$ es unión de abiertos.
\end{proof}

\begin{definition}[Bola cerrada]
  Si $x \in \R^n$ y $r > 0$, la \emph{bola cerrada} de centro $x$ y radio $r$ es
  \[
    \cl{B}_r(x) = \{ y \in \R^n : d(x, y) \leq r \}.
  \]
\end{definition}

\begin{prop}
  La bola cerrada $\cl{B}_r(x)$ es un conjunto cerrado.
\end{prop}

\begin{proof}
  Sea $y \in \R^n \setminus \cl{B}_r(x)$, es decir, $d(x, y) > r$, y sea $r_1 = d(x, y) - r > 0$. Si $z \in B_{r_1}(y)$, por la desigualdad triangular
  \[
    d(x, z) \geq d(x, y) - d(z, y) > d(x, y) - r_1 = r,
  \]
  de modo que $z \notin \cl{B}_r(x)$. Luego $B_{r_1}(y) \subset \R^n \setminus \cl{B}_r(x)$, y el complemento es abierto.
\end{proof}

\section{Clausura}

\begin{definition}[Clausura]
  La \emph{clausura} de $X \subset \R^n$ es la intersección $\cl{X}$ de todos los subconjuntos cerrados de $\R^n$ que contienen a $X$.
\end{definition}

\begin{prop}
  $\cl{X}$ es el menor subconjunto cerrado de $\R^n$ que contiene a $X$, en el sentido de que está contenido en todo cerrado que contiene a $X$.
\end{prop}

\begin{proof}
  $\cl{X}$ es cerrado por ser intersección de cerrados, contiene a $X$, y por definición está contenido en todo cerrado que contiene a $X$.
\end{proof}

\begin{corollary}
  Un subconjunto $X \subset \R^n$ es cerrado si y sólo si $\cl{X} = X$.
\end{corollary}

\begin{proof}
  Si $X$ es cerrado, es un cerrado que contiene a $X$, luego $\cl{X} \subset X$; como $X \subset \cl{X}$, vale la igualdad. Recíprocamente, si $X = \cl{X}$, entonces $X$ es cerrado pues $\cl{X}$ lo es.
\end{proof}

\begin{prop}
  Si $X \subset Y$, entonces $\cl{X} \subset \cl{Y}$.
\end{prop}

\begin{proof}
  $X \subset Y \subset \cl{Y}$, y $\cl{X}$ es el menor cerrado que contiene a $X$.
\end{proof}

\begin{prop}
  $\cl{X \cup Y} = \cl{X} \cup \cl{Y}$.
\end{prop}

\begin{proof}
  Como $X \subset X \cup Y$ e $Y \subset X \cup Y$, se tiene $\cl{X} \subset \cl{X \cup Y}$ y $\cl{Y} \subset \cl{X \cup Y}$, de donde $\cl{X} \cup \cl{Y} \subset \cl{X \cup Y}$. Por otro lado, $\cl{X} \cup \cl{Y}$ es cerrado (unión finita de cerrados) y contiene a $X \cup Y$, de modo que $\cl{X \cup Y} \subset \cl{X} \cup \cl{Y}$.
\end{proof}

\begin{prop}
  $\cl{X \cap Y} \subset \cl{X} \cap \cl{Y}$.
\end{prop}

\begin{proof}
  $\cl{X} \cap \cl{Y}$ es cerrado y contiene a $X \cap Y$, y $\cl{X \cap Y}$ es el menor cerrado que contiene a $X \cap Y$.
\end{proof}

\begin{prop}
  Sean $X \subset \R^n$ y $x \in \R^n$. Son equivalentes:
  \begin{enumerate}[label=(\arabic*)]
    \item $x \in \cl{X}$.
    \item Para todo $r > 0$, $B_r(x) \cap X \neq \varnothing$.
    \item Para todo entorno abierto $N$ de $x$, $N \cap X \neq \varnothing$.
    \item Para todo entorno $N$ de $x$, $N \cap X \neq \varnothing$.
  \end{enumerate}
\end{prop}

\begin{proof}
  (4) $\Rightarrow$ (3) $\Rightarrow$ (2) son inmediatos, pues los entornos abiertos son entornos y las bolas $B_r(x)$ son entornos abiertos de $x$.

  (2) $\Rightarrow$ (1) Por contrarrecíproco, supongamos $x \notin \cl{X}$. Entonces existe un cerrado $F$ con $X \subset F$ y $x \notin F$. Luego $x \in \R^n \setminus F$, que es abierto, y existe $r > 0$ con $B_r(x) \subset \R^n \setminus F$. Entonces $B_r(x) \cap X \subset (\R^n \setminus F) \cap F = \varnothing$.

  (1) $\Rightarrow$ (4) Por contrarrecíproco, supongamos que existe un entorno $N$ de $x$ con $N \cap X = \varnothing$. El conjunto $U = N^{\circ}$ es abierto, $x \in U$ y $U \subset N \subset \R^n \setminus X$. Entonces $X \subset \R^n \setminus U$, que es cerrado, de modo que $\cl{X} \subset \R^n \setminus U$ y, como $x \in U$, $x \notin \cl{X}$.
\end{proof}

\begin{eg}
  Sea $\Q^n$ el conjunto de puntos de $\R^n$ con todas sus coordenadas racionales. Entonces $\cl{\Q^n} = \R^n$; se dice que $\Q^n$ es \emph{denso} en $\R^n$.
\end{eg}

\begin{proof}
  Sea $x = (x_1, \dots, x_n) \in \R^n$ y $r > 0$. Para cada $i$, como $\Q$ es denso en $\R$, existe $q_i \in \Q$ con $|x_i - q_i| < \frac{r}{\sqrt{n}}$. Entonces $q = (q_1, \dots, q_n) \in \Q^n$ y
  \[
    d(x, q) = \sqrt{\sum_{i=1}^{n} {(x_i - q_i)}^2} < \sqrt{n \cdot \frac{r^2}{n}} = r,
  \]
  es decir, $q \in B_r(x) \cap \Q^n$. Por la equivalencia (1)~$\Leftrightarrow$~(2) de la proposición anterior, $x \in \cl{\Q^n}$.
\end{proof}


\clase{01/10/2024}{Compacidad}
\section{Puntos de acumulación}

\begin{definition}[Punto de acumulación]
  Sea \(X \subset \R^n\). Un punto \(x \in \R^n\) (no necesariamente de \(X\)) es un \emph{punto de acumulación} de \(X\) si \(x \in \cl{X \setminus \{x\}}\).
\end{definition}

\begin{definition}[Conjunto derivado]
  El \emph{conjunto derivado} de \(X\) es el conjunto \(X'\) de todos los puntos de acumulación de \(X\).
\end{definition}

\begin{prop}
  Dados \(X \subset \R^n\) y \(x \in \R^n\), son equivalentes:
  \begin{enumerate}[label=(\arabic*)]
    \item \(x\) es un punto de acumulación de \(X\).
    \item Para todo \(r > 0\), \(B_r(x)\) contiene un punto de \(X\) distinto de \(x\).
    \item Todo entorno abierto de \(x\) contiene un punto de \(X\) distinto de \(x\).
    \item Todo entorno de \(x\) contiene un punto de \(X\) distinto de \(x\).
  \end{enumerate}
\end{prop}

\begin{proof}
  Es la caracterización de la clausura (proposición de la sección de clausura) aplicada al conjunto \(X \setminus \{x\}\).
\end{proof}

\begin{prop}
  Si \(X \subset \R^n\), entonces \(\cl{X} = X \cup X'\).
\end{prop}

\begin{proof}
  (\(\supset\)) Se tiene \(X \subset \cl{X}\). Si \(x \in X'\), entonces \(x \in \cl{X \setminus \{x\}} \subset \cl{X}\).

  (\(\subset\)) Sea \(x \in \cl{X}\) y supongamos \(x \notin X\). Entonces \(X \setminus \{x\} = X\), de modo que \(x \in \cl{X} = \cl{X \setminus \{x\}}\), es decir, \(x \in X'\).
\end{proof}

\begin{prop}
  Sea \(X \subset \R^n\). Un punto \(x\) es de acumulación de \(X\) si y sólo si todo entorno abierto de \(x\) contiene infinitos puntos de \(X\).
\end{prop}

\begin{proof}
  (\(\Leftarrow\)) Si todo entorno abierto de \(x\) contiene infinitos puntos de \(X\), contiene alguno distinto de \(x\), y entonces \(x \in X'\).

  (\(\Rightarrow\)) Sea \(x \in X'\) y \(N\) un entorno abierto de \(x\); existe \(r > 0\) con \(B_r(x) \subset N\). Construimos recursivamente puntos \(x_1, x_2, \dots \in N \cap X\), todos distintos de \(x\), con \(d(x_1, x) > d(x_2, x) > \cdots > 0\). Como \(x \in X'\), \(B_r(x)\) contiene algún \(x_1 \in X\) con \(x_1 \neq x\). Dado \(x_m\), la bola \(B_{d(x_m, x)}(x)\) es un entorno abierto de \(x\), y contiene algún \(x_{m+1} \in X\) distinto de \(x\); entonces \(0 < d(x_{m+1}, x) < d(x_m, x)\). Como las distancias a \(x\) son estrictamente decrecientes, los \(x_m\) son distintos dos a dos. Luego \(N \cap X\) es infinito.
\end{proof}

\begin{prop}
  Si \(X \subset \R^n\), entonces \({(X')}' \subset X'\).
\end{prop}

\begin{proof}
  Sea \(x \in {(X')}'\) y \(N\) un entorno abierto de \(x\). Existe \(y \in X'\) con \(y \in N\) e \(y \neq x\). Como \(N\) es un entorno abierto de \(y\) y \(y \in X'\), \(N\) contiene infinitos puntos de \(X\). Por la proposición anterior, \(x \in X'\).
\end{proof}

\begin{corollary}
  El conjunto derivado de todo subconjunto de \(\R^n\) es cerrado.
\end{corollary}

\begin{proof}
  \(\cl{X'} = X' \cup {(X')}' = X'\), luego \(X'\) es cerrado.
\end{proof}

\begin{definition}[Conjunto perfecto]
  Un conjunto \(X \subset \R^n\) es \emph{perfecto} si \(X = X'\).
\end{definition}

\begin{definition}[Punto aislado]
  Sea \(X \subset \R^n\). Un punto \(x \in X\) es \emph{aislado} en \(X\) si existe un entorno \(N\) de \(x\) tal que \(N \cap X = \{x\}\).
\end{definition}

\begin{definition}[Conjunto discreto]
  Un conjunto es \emph{discreto} si todos sus puntos son aislados.
\end{definition}

\begin{eg}
  \(\Z \subset \R\) es discreto: para \(k \in \Z\), el entorno \(\left( k - \frac{1}{2}, k + \frac{1}{2} \right)\) cumple \(\left( k - \frac{1}{2}, k + \frac{1}{2} \right) \cap \Z = \{k\}\).
\end{eg}

\section{Sucesiones en varias dimensiones}

En esta sección trabajamos en \(\R^m\), para no confundir la dimensión con el índice de las sucesiones. Una sucesión \({(x_n)}_{n \in \N} \subset \R^m\) \emph{converge} a \(L \in \R^m\) si
\[
  (\forall \e > 0)(\exists n_0 \in \N) \colon (\forall n > n_0) \quad d(x_n, L) < \e.
\]

\begin{lemma}
  Si una sucesión en \(\R^m\) tiene límite, éste es único.
\end{lemma}

\begin{proof}
  Si \(x_n \to L\) y \(x_n \to L'\), dado \(\e > 0\) existe \(n_0\) tal que para \(n > n_0\) vale \(d(x_n, L) < \e\) y \(d(x_n, L') < \e\). Entonces \(d(L, L') \leq d(L, x_n) + d(x_n, L') < 2\e\). Como \(\e\) es arbitrario, \(d(L, L') = 0\) y \(L = L'\).
\end{proof}

\begin{prop}
  Sean \({(x_n)}_{n \in \N} \subset \R^m\) y \(L \in \R^m\). Son equivalentes:
  \begin{enumerate}[label=(\arabic*)]
    \item \(L = \lim_{n \to \infty} x_n\).
    \item Para todo \(r > 0\) existe \(n_0\) tal que \(x_n \in B_r(L)\) para todo \(n > n_0\).
    \item Para todo entorno abierto \(N\) de \(L\) existe \(n_0\) tal que \(x_n \in N\) para todo \(n > n_0\).
    \item Para todo entorno \(N\) de \(L\) existe \(n_0\) tal que \(x_n \in N\) para todo \(n > n_0\).
  \end{enumerate}
\end{prop}

\begin{proof}
  (4) \(\Rightarrow\) (3) \(\Rightarrow\) (2) \(\Rightarrow\) (1) son inmediatos de las definiciones. (1) \(\Rightarrow\) (4): sea \(N\) un entorno de \(L\). Como \(L \in N^{\circ}\), existe \(r > 0\) con \(B_r(L) \subset N^{\circ} \subset N\). Por (1) existe \(n_0\) con \(d(L, x_n) < r\) para \(n > n_0\), es decir, \(x_n \in B_r(L) \subset N\).
\end{proof}

\begin{prop}
  Toda sucesión convergente en \(\R^m\) es acotada.
\end{prop}

\begin{proof}
  Sea \(x_n \to x\). Existe \(n_0\) tal que \(d(x_n, x) < 1\) para \(n > n_0\), y entonces \(\norm{x_n} \leq \norm{x_n - x} + \norm{x} < 1 + \norm{x}\). Luego \(\norm{x_n} \leq M\) para todo \(n\), donde \(M = \max\{ \norm{x_1}, \dots, \norm{x_{n_0}}, 1 + \norm{x} \}\).
\end{proof}

\begin{prop}
  Si \(a_n \to a\) y \(b_n \to b\) en \(\R^m\), entonces \(a_n + b_n \to a + b\).
\end{prop}

\begin{proof}
  Dado \(\e > 0\), existen \(n_0\) y \(n_1\) tales que \(d(a_n, a) < \frac{\e}{2}\) para \(n > n_0\) y \(d(b_n, b) < \frac{\e}{2}\) para \(n > n_1\). Para \(n > \max\{n_0, n_1\}\),
  \[
    d(a_n + b_n, a + b) = \norm{(a_n - a) + (b_n - b)} \leq d(a_n, a) + d(b_n, b) < \e. \qedhere
  \]
\end{proof}

\begin{prop}
  Si \(x_n \to x\) en \(\R^m\) y \(\lambda_n \to \lambda\) en \(\R\), entonces \(\lambda_n x_n \to \lambda x\).
\end{prop}

\begin{proof}
  Como ambas sucesiones convergen, son acotadas: existe \(M > 0\) con \(\norm{x_n} \leq M\) para todo \(n\) y \(|\lambda| \leq M\). Dado \(\e > 0\), existen \(n_0\) y \(n_1\) tales que \(\norm{x_n - x} < \frac{\e}{2M}\) para \(n > n_0\) y \(|\lambda_n - \lambda| < \frac{\e}{2M}\) para \(n > n_1\). Para \(n > \max\{n_0, n_1\}\),
  \[
    \norm{\lambda_n x_n - \lambda x} \leq \norm{x_n} \, |\lambda_n - \lambda| + |\lambda| \, \norm{x_n - x} < M \cdot \frac{\e}{2M} + M \cdot \frac{\e}{2M} = \e. \qedhere
  \]
\end{proof}

\begin{prop}
  Sea \({(x_n)}_{n \in \N} \subset \R^m\), con \(x_n = \bigl( x_n^{(1)}, \dots, x_n^{(m)} \bigr)\). Son equivalentes:
  \begin{enumerate}[label=(\arabic*)]
    \item La sucesión \({(x_n)}\) converge en \(\R^m\).
    \item Para todo \(i \in \{1, \dots, m\}\), la sucesión \({\bigl( x_n^{(i)} \bigr)}_{n \in \N}\) converge en \(\R\).
  \end{enumerate}
  Además, en ese caso, \(\lim_{n \to \infty} x_n = (L_1, \dots, L_m)\), donde \(L_i = \lim_{n \to \infty} x_n^{(i)}\).
\end{prop}

\begin{proof}
  (1) \(\Rightarrow\) (2) Sea \(x_n \to L = (L_1, \dots, L_m)\) y \(\e > 0\). Existe \(n_0\) tal que \(\norm{x_n - L} < \e\) para \(n > n_0\). Para cada \(i\),
  \[
    \bigl| x_n^{(i)} - L_i \bigr| \leq \sqrt{\sum_{j=1}^{m} {\bigl( x_n^{(j)} - L_j \bigr)}^2} = \norm{x_n - L} < \e,
  \]
  de modo que \(x_n^{(i)} \to L_i\).

  (2) \(\Rightarrow\) (1) Sea \(x_n^{(i)} \to L_i\) para cada \(i\), y \(L = (L_1, \dots, L_m)\). Dado \(\e > 0\), para cada \(i\) existe \(n_i\) con \(\bigl| x_n^{(i)} - L_i \bigr| < \frac{\e}{\sqrt{m}}\) para \(n > n_i\). Sea \(n_0 = \max\{n_1, \dots, n_m\}\). Para \(n > n_0\),
  \[
    \norm{x_n - L} = \sqrt{\sum_{i=1}^{m} {\bigl( x_n^{(i)} - L_i \bigr)}^2} < \sqrt{m \cdot \frac{\e^2}{m}} = \e. \qedhere
  \]
\end{proof}

\section{Conjuntos compactos}

\begin{theorem}
  Un subconjunto \(F \subset \R^m\) es cerrado si y sólo si para toda sucesión \({(x_n)} \subset F\) tal que \(x_n \to x\) se cumple \(x \in F\).
\end{theorem}

\begin{proof}
  (\(\Rightarrow\)) Sea \(F\) cerrado, \({(x_n)} \subset F\) y \(x_n \to x\). Si fuera \(x \notin F\), entonces \(x \in \R^m \setminus F\), que es abierto, y existiría \(r > 0\) con \(B_r(x) \subset \R^m \setminus F\). Como \(x_n \to x\), existe \(n_0\) con \(x_n \in B_r(x)\) para \(n > n_0\), y entonces \(x_n \notin F\), lo cual es absurdo. Luego \(x \in F\).

  (\(\Leftarrow\)) Supongamos que \(F\) no es cerrado. Entonces \(\cl{F} \neq F\), y existe \(x \in \cl{F} \setminus F\). Para cada \(n \in \N\), por la caracterización de la clausura existe \(x_n \in F \cap B_{1/n}(x)\). Entonces \({(x_n)} \subset F\) y \(x_n \to x \notin F\), lo que contradice la hipótesis. Luego \(F\) es cerrado.
\end{proof}

\begin{eg}
  El intervalo \((0, 1]\) no es cerrado: \(\frac{1}{n} \in (0, 1]\) para todo \(n\) y \(\frac{1}{n} \to 0 \notin (0, 1]\).
\end{eg}

\begin{definition}[Relativamente compacto]
  Un subconjunto \(F \subset \R^m\) es \emph{relativamente compacto} si toda sucesión en \(F\) tiene una subsucesión convergente (en \(\R^m\)).
\end{definition}

\begin{definition}[Compacto]
  Un subconjunto \(F \subset \R^m\) es \emph{compacto} si toda sucesión en \(F\) tiene una subsucesión que converge a un punto de \(F\). Equivalentemente, \(F\) es relativamente compacto y cerrado.
\end{definition}

\begin{note}
  La equivalencia se debe a que, si \(F\) es compacto, toda sucesión convergente en \(F\) converge a un punto de \(F\) (pues una subsucesión lo hace), de modo que \(F\) es cerrado por el teorema anterior. Recíprocamente, si \(F\) es cerrado y relativamente compacto, el límite de la subsucesión convergente pertenece a \(F\) por el mismo teorema.
\end{note}

\begin{prop}
  Un subconjunto \(F \subset \R^m\) es relativamente compacto si y sólo si es acotado.
\end{prop}

\begin{proof}
  (\(\Rightarrow\)) Por contrarrecíproco, supongamos que \(F\) no es acotado. Para cada \(n \in \N\) existe \(y_n \in F\) con \(\norm{y_n} > n\). Si \(f\colon \N \to \N\) es estrictamente creciente, \(\norm{y_{f(n)}} > f(n) \geq n\), de modo que \({(y_{f(n)})}\) no es acotada y por lo tanto no converge. Luego ninguna subsucesión de \({(y_n)}\) converge, y \(F\) no es relativamente compacto.

  (\(\Leftarrow\)) Sea \(F\) acotado y \({(x_k)}_{k \in \N} \subset F\). Como \(\norm{x_k} \leq K\) para algún \(K\), cada sucesión de coordenadas \({\bigl( x_k^{(i)} \bigr)}_{k}\) es acotada en \(\R\). Por el teorema de Bolzano--Weierstrass en \(\R\) existe una función estrictamente creciente \(h_1\colon \N \to \N\) tal que \({\bigl( x_{h_1(k)}^{(1)} \bigr)}_{k}\) converge. Del mismo modo, como \({\bigl( x_{h_1(k)}^{(2)} \bigr)}_{k}\) es acotada, existe \(g_2\) estrictamente creciente tal que, con \(h_2 = h_1 \circ g_2\), la sucesión \({\bigl( x_{h_2(k)}^{(2)} \bigr)}_{k}\) converge; la primera coordenada de \({(x_{h_2(k)})}\) sigue convergiendo por ser subsucesión de una convergente. Repitiendo \(m\) veces se obtiene una subsucesión \({(x_{h_m(k)})}_{k}\) en la cual todas las coordenadas convergen, y por la proposición anterior converge en \(\R^m\).
\end{proof}

\begin{corollary}[Heine--Borel]
  Un subconjunto \(F \subset \R^m\) es compacto si y sólo si es cerrado y acotado.
\end{corollary}

\begin{proof}
  Es consecuencia de la definición de compacto y de la proposición anterior.
\end{proof}

\section{Cubrimientos y compacidad}

Los resultados siguientes corresponden a la clase siguiente, pero se agrupan aquí para reunir lo referente a conjuntos compactos. En esta sección volvemos a trabajar en \(\R^n\).

\begin{definition}[Cubrimiento]
  Sea \(X \subset \R^n\). Una familia \(\mathscr{U}\) de subconjuntos de \(\R^n\) es un \emph{cubrimiento} de \(X\) si \(X \subset \bigcup_{U \in \mathscr{U}} U\). Es un \emph{cubrimiento abierto} si todos sus elementos son abiertos.
\end{definition}

\begin{definition}[Subcubrimiento]
  Un \emph{subcubrimiento} de \(\mathscr{U}\) es una subfamilia \(\mathscr{V} \subset \mathscr{U}\) que también es un cubrimiento de \(X\).
\end{definition}

\begin{theorem}[Lindelöf]
  Todo cubrimiento abierto de un subconjunto de \(\R^n\) admite un subcubrimiento a lo sumo numerable.
\end{theorem}

\begin{proof}
  Sea \(X \subset \R^n\) y \(\mathscr{U}\) un cubrimiento abierto de \(X\). Si \(X = \varnothing\), la familia vacía es un subcubrimiento. Si \(X \neq \varnothing\), entonces \(\mathscr{U} \neq \varnothing\) y fijamos \(U_0 \in \mathscr{U}\).

  Para cada \(q \in \Q^n\) y cada \(s \in \Q\), \(s > 0\), elegimos \(U(q, s) \in \mathscr{U}\) así: si existen elementos de \(\mathscr{U}\) que contienen a \(B_s(q)\), elegimos uno cualquiera; si no, ponemos \(U(q, s) = U_0\).

  Sea \(x \in X\). Existen \(U \in \mathscr{U}\) con \(x \in U\) y \(r > 0\) con \(B_r(x) \subset U\). Como \(\Q^n\) es denso en \(\R^n\), existe \(q \in \Q^n\) con \(d(x, q) < \frac{r}{2}\), y existe \(s \in \Q\) con \(d(x, q) < s < \frac{r}{2}\). Entonces \(x \in B_s(q)\), y si \(y \in B_s(q)\), \(d(y, x) \leq d(y, q) + d(q, x) < 2s < r\), es decir, \(B_s(q) \subset B_r(x) \subset U\). Luego hay elementos de \(\mathscr{U}\) que contienen a \(B_s(q)\), de modo que \(x \in B_s(q) \subset U(q, s)\).

  Esto prueba que \(\mathscr{V} = \{ U(q, s) : q \in \Q^n, \ s \in \Q, \ s > 0 \}\) es un subcubrimiento de \(\mathscr{U}\). Es a lo sumo numerable porque es la imagen de \(\Q^n \times (\Q \cap (0, \infty))\), que es numerable, por la función \((q, s) \mapsto U(q, s)\).
\end{proof}

\begin{theorem}[Heine--Borel--Lebesgue]
  Un subconjunto \(F \subset \R^n\) es compacto si y sólo si todo cubrimiento abierto de \(F\) admite un subcubrimiento finito.
\end{theorem}

\begin{proof}
  (\(\Rightarrow\)) Sea \(F\) compacto y \(\mathscr{U}\) un cubrimiento abierto de \(F\). Por el teorema de Lindelöf existe un subcubrimiento \(\mathscr{V}\) a lo sumo numerable; si es finito no hay nada que probar, de modo que supongamos \(\mathscr{V} = \{ V_1, V_2, \dots \}\) numerable. Supongamos, por el absurdo, que \(\mathscr{U}\) no admite subcubrimiento finito. En particular, para cada \(m\) la familia \(\{ V_1, \dots, V_m \}\) no cubre \(F\), y existe \(x_m \in F \setminus \bigcup_{i=1}^{m} V_i\). Como \(F\) es compacto, existe una subsucesión \(x_{h(m)} \to L \in F\) con \(h\) estrictamente creciente. Como \(\mathscr{V}\) cubre \(F\), existe \(s\) con \(L \in V_s\). Como \(V_s\) es abierto, existe \(m_0\) tal que \(x_{h(m)} \in V_s\) para \(m > m_0\). Pero para \(m > \max\{m_0, s\}\) se tiene \(h(m) \geq m > s\), y por construcción \(x_{h(m)} \notin V_i\) para todo \(i \leq h(m)\), en particular \(x_{h(m)} \notin V_s\). Absurdo.

  (\(\Leftarrow\)) Supongamos que todo cubrimiento abierto de \(F\) admite un subcubrimiento finito, pero que \(F\) no es compacto. Entonces existe \({(x_m)} \subset F\) sin subsucesiones convergentes a puntos de \(F\). Sea \(x \in F\). Como \(x\) no es límite de ninguna subsucesión, por el mismo argumento que en \(\R\) (un punto es límite de una subsucesión si y sólo si todo entorno contiene infinitos términos) existe \(r_x > 0\) tal que \(S_x = \{ m \in \N : x_m \in B_{r_x}(x) \}\) es finito. La familia \(\{ B_{r_x}(x) : x \in F \}\) es un cubrimiento abierto de \(F\), y por hipótesis existen \(y_1, \dots, y_k \in F\) con \(F \subset \bigcup_{i=1}^{k} B_{r_{y_i}}(y_i)\). Como \(x_m \in F\) para todo \(m\), cada \(m \in \N\) pertenece a algún \(S_{y_i}\), es decir, \(\N \subset S_{y_1} \cup \cdots \cup S_{y_k}\), que es finito. Absurdo.
\end{proof}

\begin{prop}[Número de Lebesgue]
  Sea \(K \subset \R^n\) compacto y \(\mathscr{U}\) un cubrimiento abierto de \(K\). Entonces existe \(\delta > 0\) tal que si \(x, y \in K\) y \(d(x, y) < \delta\), existe \(U \in \mathscr{U}\) con \(x, y \in U\). Un tal \(\delta\) se llama \emph{número de Lebesgue} del cubrimiento.
\end{prop}

\begin{proof}
  Para cada \(x \in K\) elegimos \(U_x \in \mathscr{U}\) con \(x \in U_x\) y \(r_x > 0\) con \(B_{r_x}(x) \subset U_x\). La familia \(\bigl\{ B_{r_x/2}(x) : x \in K \bigr\}\) es un cubrimiento abierto de \(K\), y por el teorema anterior existen \(x_1, \dots, x_k \in K\) con \(K \subset \bigcup_{i=1}^{k} B_{r_{x_i}/2}(x_i)\). Sea \(\delta = \frac{1}{2} \min\{ r_{x_1}, \dots, r_{x_k} \}\).

  Sean \(x, y \in K\) con \(d(x, y) < \delta\). Existe \(i\) con \(x \in B_{r_{x_i}/2}(x_i)\), y entonces
  \[
    d(x_i, y) \leq d(x_i, x) + d(x, y) < \frac{r_{x_i}}{2} + \delta \leq \frac{r_{x_i}}{2} + \frac{r_{x_i}}{2} = r_{x_i}.
  \]
  Luego \(x, y \in B_{r_{x_i}}(x_i) \subset U_{x_i} \in \mathscr{U}\).
\end{proof}


\clase{04/10/2024}{Límite de una función}
\section{Límite de una función}

\begin{definition}[Límite]
  Sean \(A \subset \R^n\), \(x_0 \in A'\) (punto de acumulación de \(A\)) y \(f\colon A \to \R^k\) una función. Decimos que \(y \in \R^k\) es el \emph{límite} de \(f\) en \(x_0\), y escribimos \(\lim_{x \to x_0} f(x) = y\), si
  \[
    (\forall \e > 0)(\exists \delta > 0) \colon (\forall x \in A) \quad 0 < \norm{x - x_0} < \delta \Rightarrow \norm{f(x) - y} < \e.
  \]
\end{definition}

\begin{lemma}
  Si \(A \subset \R^n\) y \(x_0 \in A'\), una función \(f\colon A \to \R^k\) tiene a lo sumo un límite en \(x_0\).
\end{lemma}

\begin{proof}
  Sean \(y, y' \in \R^k\) límites de \(f\) en \(x_0\) y supongamos \(y \neq y'\). Sea \(\e = \frac{1}{2} \norm{y - y'} > 0\). Existen \(\delta_1, \delta_2 > 0\) tales que
  \[
    0 < \norm{x - x_0} < \delta_1 \Rightarrow \norm{f(x) - y} < \e, \qquad 0 < \norm{x - x_0} < \delta_2 \Rightarrow \norm{f(x) - y'} < \e.
  \]
  Sea \(\delta = \min\{\delta_1, \delta_2\}\). Como \(x_0 \in A'\), existe \(x \in A\) con \(0 < \norm{x - x_0} < \delta\), y para ese \(x\),
  \[
    \norm{y - y'} \leq \norm{y - f(x)} + \norm{f(x) - y'} < 2\e = \norm{y - y'},
  \]
  lo cual es absurdo. Luego \(y = y'\).
\end{proof}

\begin{prop}
  Sean \(A \subset \R^n\), \(x_0 \in A'\), \(f\colon A \to \R^k\) e \(y \in \R^k\). Son equivalentes:
  \begin{enumerate}[label=(\arabic*)]
    \item \(\lim_{x \to x_0} f(x) = y\).
    \item Para todo \(\e > 0\) existe \(\delta > 0\) tal que \(f\bigl( B_{\delta}(x_0) \cap (A \setminus \{x_0\}) \bigr) \subset B_{\e}(y)\).
    \item Para todo entorno \(V\) de \(y\) existe un entorno \(U\) de \(x_0\) tal que \(f\bigl( U \cap (A \setminus \{x_0\}) \bigr) \subset V\).
  \end{enumerate}
\end{prop}

\begin{proof}
  (1) \(\Rightarrow\) (2) Dado \(\e > 0\), sea \(\delta > 0\) como en la definición. Si \(x \in B_{\delta}(x_0) \cap (A \setminus \{x_0\})\), entonces \(0 < \norm{x - x_0} < \delta\), luego \(\norm{f(x) - y} < \e\), es decir, \(f(x) \in B_{\e}(y)\).

  (2) \(\Rightarrow\) (3) Sea \(V\) un entorno de \(y\). Existe \(\e > 0\) con \(B_{\e}(y) \subset V\). Por (2) existe \(\delta > 0\) con \(f\bigl( B_{\delta}(x_0) \cap (A \setminus \{x_0\}) \bigr) \subset B_{\e}(y) \subset V\). Basta tomar \(U = B_{\delta}(x_0)\), que es un entorno de \(x_0\).

  (3) \(\Rightarrow\) (1) Sea \(\e > 0\). Como \(B_{\e}(y)\) es un entorno de \(y\), existe un entorno \(U\) de \(x_0\) con \(f\bigl( U \cap (A \setminus \{x_0\}) \bigr) \subset B_{\e}(y)\). Como \(U\) es entorno de \(x_0\), existe \(\delta > 0\) con \(B_{\delta}(x_0) \subset U\). Si \(x \in A\) cumple \(0 < \norm{x - x_0} < \delta\), entonces \(x \in U \cap (A \setminus \{x_0\})\), luego \(f(x) \in B_{\e}(y)\), es decir, \(\norm{f(x) - y} < \e\).
\end{proof}

\begin{prop}[Caracterización por sucesiones]
  Sean \(A \subset \R^n\), \(x_0 \in A'\), \(f\colon A \to \R^k\) e \(y \in \R^k\). Entonces \(\lim_{x \to x_0} f(x) = y\) si y sólo si para toda sucesión \({(x_m)}_{m \in \N} \subset A \setminus \{x_0\}\) con \(x_m \to x_0\) se cumple \(f(x_m) \to y\).
\end{prop}

\begin{proof}
  (\(\Rightarrow\)) Sea \({(x_m)} \subset A \setminus \{x_0\}\) con \(x_m \to x_0\) y sea \(\e > 0\). Existe \(\delta > 0\) tal que \(0 < \norm{x - x_0} < \delta \Rightarrow \norm{f(x) - y} < \e\). Como \(x_m \to x_0\), existe \(m_0\) con \(\norm{x_m - x_0} < \delta\) para \(m > m_0\), y como \(x_m \neq x_0\), resulta \(\norm{f(x_m) - y} < \e\) para \(m > m_0\).

  (\(\Leftarrow\)) Por contrarrecíproco, supongamos que \(\lim_{x \to x_0} f(x) \neq y\). Entonces existe \(\e > 0\) tal que para todo \(\delta > 0\) hay un \(x \in A\) con \(0 < \norm{x - x_0} < \delta\) y \(\norm{f(x) - y} \geq \e\). Tomando \(\delta = \frac{1}{m}\) obtenemos \(x_m \in A \setminus \{x_0\}\) con \(\norm{x_m - x_0} < \frac{1}{m}\), de modo que \(x_m \to x_0\), pero \(\norm{f(x_m) - y} \geq \e\) para todo \(m\), es decir, \(f(x_m) \not\to y\).
\end{proof}

\begin{prop}[Álgebra de límites]
  Sean \(A \subset \R^n\), \(x_0 \in A'\), \(f, g\colon A \to \R^k\) y \(h\colon A \to \R\).
  \begin{enumerate}[label=(\arabic*)]
    \item Si \(\lim_{x \to x_0} f(x) = y\) y \(\lim_{x \to x_0} g(x) = z\), entonces \(\lim_{x \to x_0} \bigl( f(x) + g(x) \bigr) = y + z\).
    \item Si \(\lim_{x \to x_0} h(x) = r\) y \(\lim_{x \to x_0} f(x) = y\), entonces \(\lim_{x \to x_0} h(x) f(x) = r y\).
    \item Si \(\lim_{x \to x_0} h(x) = r \neq 0\) y \(h(x) \neq 0\) para todo \(x \in A\), entonces \(\lim_{x \to x_0} \dfrac{1}{h(x)} = \dfrac{1}{r}\).
  \end{enumerate}
\end{prop}

\begin{proof}
  Sea \({(x_m)} \subset A \setminus \{x_0\}\) con \(x_m \to x_0\). Por la caracterización por sucesiones, \(f(x_m) \to y\), \(g(x_m) \to z\) y \(h(x_m) \to r\). Por las propiedades de límites de sucesiones, \(f(x_m) + g(x_m) \to y + z\), \(h(x_m) f(x_m) \to r y\) y \(\frac{1}{h(x_m)} \to \frac{1}{r}\). Como esto vale para toda tal sucesión, se obtienen (1), (2) y (3).
\end{proof}

\begin{prop}
  Sean \(A \subset \R^n\), \(x_0 \in A'\), \(f\colon A \to \R^k\) y \(\lim_{x \to x_0} f(x) = y\).
  \begin{enumerate}[label=(\arabic*)]
    \item Si \(z \neq y\), existe \(r > 0\) tal que \(z \notin f\bigl( B_r(x_0) \cap (A \setminus \{x_0\}) \bigr)\).
    \item Si \(k = 1\) e \(y > 0\), existe \(r > 0\) tal que \(f(x) > 0\) para todo \(x \in B_r(x_0) \cap (A \setminus \{x_0\})\).
    \item Existe \(r > 0\) tal que \(f\) es acotada en \(B_r(x_0) \cap A\), es decir, el conjunto \(\{ \norm{f(x)} : x \in B_r(x_0) \cap A \}\) está acotado superiormente.
  \end{enumerate}
\end{prop}

\begin{proof}
  (1) Como \(\norm{y - z} > 0\), tomando \(\e = \norm{y - z}\) existe \(r > 0\) tal que \(0 < \norm{x - x_0} < r \Rightarrow \norm{f(x) - y} < \norm{y - z}\). En particular \(f(x) \neq z\) para tales \(x\).

  (2) Tomando \(\e = y > 0\), existe \(r > 0\) con \(|f(x) - y| < y\), es decir, \(f(x) > 0\), para todo \(x \in A\) con \(0 < \norm{x - x_0} < r\).

  (3) Tomando \(\e = 1\), existe \(r > 0\) tal que para \(x \in A\) con \(0 < \norm{x - x_0} < r\) se tiene \(\norm{f(x)} \leq \norm{f(x) - y} + \norm{y} < 1 + \norm{y}\). Si \(x_0 \notin A\), entonces \(1 + \norm{y}\) es cota de \(\norm{f(x)}\) en \(B_r(x_0) \cap A\). Si \(x_0 \in A\), el punto \(x_0\) no está contemplado por la definición de límite, y una cota es \(\max\{ 1 + \norm{y}, \norm{f(x_0)} \}\).
\end{proof}

\begin{theorem}[Teorema del sándwich]
  Sean \(A \subset \R^n\), \(x_0 \in A'\) y \(f, g, h\colon A \to \R\) tales que \(f(x) \leq g(x) \leq h(x)\) para todo \(x \in A\). Si \(\lim_{x \to x_0} f(x) = \lim_{x \to x_0} h(x) = y\), entonces \(\lim_{x \to x_0} g(x) = y\).
\end{theorem}

\begin{proof}
  Dado \(\e > 0\), existen \(\delta_1, \delta_2 > 0\) tales que
  \[
    0 < \norm{x - x_0} < \delta_1 \Rightarrow |f(x) - y| < \e, \qquad 0 < \norm{x - x_0} < \delta_2 \Rightarrow |h(x) - y| < \e.
  \]
  Sea \(\delta = \min\{\delta_1, \delta_2\}\). Si \(x \in A\) cumple \(0 < \norm{x - x_0} < \delta\), entonces
  \[
    y - \e < f(x) \leq g(x) \leq h(x) < y + \e,
  \]
  es decir, \(|g(x) - y| < \e\). Luego \(\lim_{x \to x_0} g(x) = y\).
\end{proof}


\clase{08/10/2024}{Continuidad}
\section{Límites en varias dimensiones}

\begin{prop}
  Sean $A \subset \R^n$, $x_0 \in A'$, $f\colon A \to \R^k$, $f = (f_1, \dots, f_k)$ e $y = (y_1, \dots, y_k) \in \R^k$. Son equivalentes:
  \begin{enumerate}[label=(\arabic*)]
    \item $\lim_{x \to x_0} f(x) = y$.
    \item Para todo $i \in \{1, \dots, k\}$, $\lim_{x \to x_0} f_i(x) = y_i$.
  \end{enumerate}
\end{prop}

\begin{proof}
  (1) $\Rightarrow$ (2) Dado $\e > 0$ existe $\delta > 0$ tal que, si $x \in A$ y $0 < \norm{x - x_0} < \delta$, entonces $\norm{f(x) - y} < \e$. Para cada $i$, para esos mismos $x$,
  \[
    |f_i(x) - y_i| \leq \sqrt{\sum_{j=1}^{k} {\bigl( f_j(x) - y_j \bigr)}^2} = \norm{f(x) - y} < \e,
  \]
  luego $\lim_{x \to x_0} f_i(x) = y_i$.

  (2) $\Rightarrow$ (1) Dado $\e > 0$, para cada $i$ existe $\delta_i > 0$ tal que $0 < \norm{x - x_0} < \delta_i$ implica $|f_i(x) - y_i| < \frac{\e}{\sqrt{k}}$. Sea $\delta = \min\{\delta_1, \dots, \delta_k\}$. Si $x \in A$ cumple $0 < \norm{x - x_0} < \delta$, entonces
  \[
    \norm{f(x) - y} = \sqrt{\sum_{i=1}^{k} {\bigl( f_i(x) - y_i \bigr)}^2} < \sqrt{\sum_{i=1}^{k} \frac{\e^2}{k}} = \e. \qedhere
  \]
\end{proof}

\begin{definition}[Límites laterales]
  Sean $A \subset \R$, $f\colon A \to \R^k$ y $x_0 \in \R$.
  \begin{enumerate}[label=(\arabic*)]
    \item Sea $A_{x_0}^{+} = A \cap (x_0, +\infty)$. Si $x_0 \in {(A_{x_0}^{+})}'$ y la restricción $f|_{A_{x_0}^{+}}$ tiene límite $y$ en $x_0$, decimos que $y$ es el \emph{límite por derecha} de $f$ en $x_0$ y escribimos $\lim_{x \to x_0^{+}} f(x) = y$.
    \item Sea $A_{x_0}^{-} = A \cap (-\infty, x_0)$. Si $x_0 \in {(A_{x_0}^{-})}'$ y la restricción $f|_{A_{x_0}^{-}}$ tiene límite $y$ en $x_0$, decimos que $y$ es el \emph{límite por izquierda} de $f$ en $x_0$ y escribimos $\lim_{x \to x_0^{-}} f(x) = y$.
  \end{enumerate}
\end{definition}

\section{Continuidad en un punto}

\begin{definition}[Continuidad en un punto]
  Sean $A \subset \R^n$ y $x_0 \in A$. Una función $f\colon A \to \R^k$ es \emph{continua en $x_0$} si
  \[
    (\forall \e > 0)(\exists \delta > 0) \colon (\forall x \in A) \quad \norm{x - x_0} < \delta \Rightarrow \norm{f(x) - f(x_0)} < \e.
  \]
\end{definition}

\begin{lemma}
  Sean $A \subset \R^n$ y $x_0 \in A$. Una función $f\colon A \to \R^k$ es continua en $x_0$ si y sólo si $x_0$ es un punto aislado de $A$ o bien $\lim_{x \to x_0} f(x) = f(x_0)$.
\end{lemma}

\begin{proof}
  ($\Rightarrow$) Sea $f$ continua en $x_0$ y supongamos que $x_0$ no es aislado en $A$. Entonces toda bola centrada en $x_0$ contiene puntos de $A$ distintos de $x_0$, es decir, $x_0 \in A'$. Dado $\e > 0$, existe $\delta > 0$ tal que $\norm{x - x_0} < \delta \Rightarrow \norm{f(x) - f(x_0)} < \e$; en particular vale para $0 < \norm{x - x_0} < \delta$, y entonces $\lim_{x \to x_0} f(x) = f(x_0)$.

  ($\Leftarrow$) Sea $\e > 0$. Si $x_0$ es aislado en $A$, existe $\delta > 0$ con $B_{\delta}(x_0) \cap A = \{x_0\}$, y para $x \in A$ con $\norm{x - x_0} < \delta$ se tiene $x = x_0$ y $\norm{f(x) - f(x_0)} = 0 < \e$. Si $\lim_{x \to x_0} f(x) = f(x_0)$, existe $\delta > 0$ tal que $0 < \norm{x - x_0} < \delta \Rightarrow \norm{f(x) - f(x_0)} < \e$; para $x = x_0$ la desigualdad vale trivialmente, de modo que vale para todo $x \in A$ con $\norm{x - x_0} < \delta$.
\end{proof}

\begin{definition}[Entorno relativo]
  Dado $A \subset \R^n$, un conjunto $V$ es un \emph{entorno de $x$ relativo a $A$} si existe un entorno $N \subset \R^n$ de $x$ tal que $V = A \cap N$.
\end{definition}

\begin{definition}[Abierto relativo]
  Dado $A \subset \R^n$, un conjunto $U$ es \emph{abierto relativo a $A$} si existe un abierto $\tilde{U} \subset \R^n$ tal que $U = \tilde{U} \cap A$.
\end{definition}

\begin{definition}[Cerrado relativo]
  Dado $A \subset \R^n$, un conjunto $F$ es \emph{cerrado relativo a $A$} si existe un cerrado $\tilde{F} \subset \R^n$ tal que $F = \tilde{F} \cap A$.
\end{definition}

\begin{prop}
  Sean $A \subset \R^n$, $x_0 \in A$ y $f\colon A \to \R^k$. Son equivalentes:
  \begin{enumerate}[label=(\arabic*)]
    \item $f$ es continua en $x_0$.
    \item Para todo entorno $V$ de $f(x_0)$ existe un entorno $U$ de $x_0$ relativo a $A$ tal que $f(U) \subset V$.
    \item Para todo entorno $V$ de $f(x_0)$, el conjunto $f^{-1}(V)$ es un entorno de $x_0$ relativo a $A$.
  \end{enumerate}
\end{prop}

\begin{proof}
  (1) $\Rightarrow$ (2) Sea $V$ un entorno de $f(x_0)$; existe $\e > 0$ con $B_{\e}(f(x_0)) \subset V$. Por continuidad existe $\delta > 0$ tal que, para $x \in A$, $\norm{x - x_0} < \delta \Rightarrow \norm{f(x) - f(x_0)} < \e$. El conjunto $U = A \cap B_{\delta}(x_0)$ es un entorno de $x_0$ relativo a $A$, y si $x \in U$ entonces $f(x) \in B_{\e}(f(x_0)) \subset V$. Luego $f(U) \subset V$.

  (2) $\Rightarrow$ (3) Sea $V$ un entorno de $f(x_0)$. Por hipótesis existe un entorno $U$ de $x_0$ relativo a $A$ con $f(U) \subset V$, es decir, $U \subset f^{-1}(V)$. Existe un entorno $N$ de $x_0$ con $U = A \cap N$. Como $U \subset f^{-1}(V) \subset A$, resulta $f^{-1}(V) = A \cap \bigl( N \cup f^{-1}(V) \bigr)$, y $N \cup f^{-1}(V)$ es un entorno de $x_0$ en $\R^n$. Luego $f^{-1}(V)$ es un entorno de $x_0$ relativo a $A$.

  (3) $\Rightarrow$ (1) Sea $\e > 0$. Como $B_{\e}(f(x_0))$ es un entorno de $f(x_0)$, el conjunto $f^{-1}\bigl( B_{\e}(f(x_0)) \bigr)$ es un entorno de $x_0$ relativo a $A$: existe un entorno $N$ de $x_0$ en $\R^n$ con $f^{-1}\bigl( B_{\e}(f(x_0)) \bigr) = A \cap N$. Sea $\delta > 0$ con $B_{\delta}(x_0) \subset N$. Si $x \in A$ y $\norm{x - x_0} < \delta$, entonces $x \in A \cap N = f^{-1}\bigl( B_{\e}(f(x_0)) \bigr)$, es decir, $\norm{f(x) - f(x_0)} < \e$.
\end{proof}

\begin{prop}[Caracterización por sucesiones]
  Sean $A \subset \R^n$, $x_0 \in A$ y $f\colon A \to \R^k$. Entonces $f$ es continua en $x_0$ si y sólo si para toda sucesión ${(x_m)}_{m \in \N} \subset A$ con $x_m \to x_0$ se cumple $f(x_m) \to f(x_0)$.
\end{prop}

\begin{proof}
  ($\Rightarrow$) Sea $f$ continua en $x_0$, ${(x_m)} \subset A$ con $x_m \to x_0$ y $\e > 0$. Existe $\delta > 0$ tal que $\norm{x - x_0} < \delta \Rightarrow \norm{f(x) - f(x_0)} < \e$ para $x \in A$. Como $x_m \to x_0$, existe $m_0$ con $\norm{x_m - x_0} < \delta$ para $m > m_0$, y entonces $\norm{f(x_m) - f(x_0)} < \e$ para $m > m_0$.

  ($\Leftarrow$) Por contrarrecíproco, supongamos que $f$ no es continua en $x_0$. Existe $\e > 0$ tal que para todo $\delta > 0$ hay un $x \in A$ con $\norm{x - x_0} < \delta$ y $\norm{f(x) - f(x_0)} \geq \e$. Tomando $\delta = \frac{1}{m}$ obtenemos ${(x_m)} \subset A$ con $x_m \to x_0$ pero $\norm{f(x_m) - f(x_0)} \geq \e$ para todo $m$, es decir, $f(x_m) \not\to f(x_0)$.
\end{proof}

\begin{prop}
  Sean $A \subset \R^n$, $x_0 \in A$ y $f\colon A \to \R^k$ continua en $x_0$. Si $B \subset A$ y $x_0 \in B$, entonces $f|_B\colon B \to \R^k$ es continua en $x_0$.
\end{prop}

\begin{proof}
  Dado $\e > 0$, existe $\delta > 0$ tal que para todo $x \in A$ con $\norm{x - x_0} < \delta$ vale $\norm{f(x) - f(x_0)} < \e$. En particular vale para todo $x \in B \subset A$.
\end{proof}

\begin{prop}
  Sean $A \subset \R^n$ y $x_0 \in A$.
  \begin{enumerate}[label=(\arabic*)]
    \item Si $f, g\colon A \to \R^k$ y $h\colon A \to \R$ son continuas en $x_0$, entonces $f + g\colon A \to \R^k$ y $h \cdot f\colon A \to \R^k$ son continuas en $x_0$.
    \item Si $h\colon A \to \R$ es continua en $x_0$ y $h(x) \neq 0$ para todo $x \in A$, entonces $\frac{1}{h}\colon A \to \R$ es continua en $x_0$.
  \end{enumerate}
\end{prop}

\begin{proof}
  Sea ${(x_m)} \subset A$ con $x_m \to x_0$. Por la caracterización por sucesiones, $f(x_m) \to f(x_0)$, $g(x_m) \to g(x_0)$ y $h(x_m) \to h(x_0)$. Por las propiedades de límites de sucesiones, $f(x_m) + g(x_m) \to f(x_0) + g(x_0)$, $h(x_m) f(x_m) \to h(x_0) f(x_0)$ y, como $h(x_m) \neq 0$ y $h(x_0) \neq 0$, $\frac{1}{h(x_m)} \to \frac{1}{h(x_0)}$. Aplicando nuevamente la caracterización por sucesiones se obtienen (1) y (2).
\end{proof}

\begin{prop}
  Sean $A \subset \R^n$, $x_0 \in A$ y $f\colon A \to \R^k$ continua en $x_0$.
  \begin{enumerate}[label=(\arabic*)]
    \item Si $y_0 \neq f(x_0)$, existe un entorno $U$ de $x_0$ relativo a $A$ tal que $y_0 \notin f(U)$.
    \item Si $k = 1$ y $c \in \R$ cumple $f(x_0) < c$, existe un entorno $U$ de $x_0$ relativo a $A$ tal que $f(x) < c$ para todo $x \in U$.
    \item Existe un entorno $U$ de $x_0$ relativo a $A$ tal que $f$ es acotada en $U$, es decir, $f(U)$ es un conjunto acotado.
  \end{enumerate}
\end{prop}

\begin{proof}
  (1) Sea $r = \norm{f(x_0) - y_0} > 0$. Como $B_r(f(x_0))$ es un entorno de $f(x_0)$, por la continuidad existe un entorno $U$ de $x_0$ relativo a $A$ con $f(U) \subset B_r(f(x_0))$. Como $y_0 \notin B_r(f(x_0))$, resulta $y_0 \notin f(U)$.

  (2) El intervalo $(-\infty, c)$ es un entorno de $f(x_0)$, luego existe un entorno $U$ de $x_0$ relativo a $A$ con $f(U) \subset (-\infty, c)$, es decir, $f(x) < c$ para todo $x \in U$.

  (3) La bola $B_1(f(x_0))$ es un entorno de $f(x_0)$ y existe un entorno $U$ de $x_0$ relativo a $A$ con $f(U) \subset B_1(f(x_0))$, que es un conjunto acotado.
\end{proof}

\section{Composición de funciones}

\begin{theorem}
  Sean $A \subset \R^n$, $B \subset \R^k$, $f\colon A \to \R^k$ y $g\colon B \to \R^p$ con $f(A) \subset B$. Si $f$ es continua en $x_0 \in A$ y $g$ es continua en $f(x_0)$, entonces $g \circ f\colon A \to \R^p$ es continua en $x_0$.
\end{theorem}

\begin{proof}
  Sea $W$ un entorno de $g(f(x_0))$. Como $g$ es continua en $f(x_0)$, existe un entorno $V$ de $f(x_0)$ relativo a $B$ con $g(V) \subset W$. Existe un entorno $V_0$ de $f(x_0)$ en $\R^k$ con $V = V_0 \cap B$. Como $f$ es continua en $x_0$, existe un entorno $U$ de $x_0$ relativo a $A$ con $f(U) \subset V_0$. Como además $f(A) \subset B$, se tiene $f(U) \subset V_0 \cap B = V$, y entonces $(g \circ f)(U) \subset g(V) \subset W$. Luego $g \circ f$ es continua en $x_0$.
\end{proof}

\begin{prop}
  Sean $A \subset \R^n$, $x_0 \in A$ y $f\colon A \to \R^k$, $f = (f_1, \dots, f_k)$. Entonces $f$ es continua en $x_0$ si y sólo si cada componente $f_1, \dots, f_k\colon A \to \R$ es continua en $x_0$.
\end{prop}

\begin{proof}
  ($\Rightarrow$) Para todo $x \in A$ y todo $i$,
  \[
    |f_i(x) - f_i(x_0)| \leq \sqrt{\sum_{j=1}^{k} {\bigl( f_j(x) - f_j(x_0) \bigr)}^2} = \norm{f(x) - f(x_0)}.
  \]
  Dado $\e > 0$, el $\delta$ de la continuidad de $f$ sirve para cada $f_i$.

  ($\Leftarrow$) Dado $\e > 0$, para cada $i$ existe $\delta_i > 0$ tal que $\norm{x - x_0} < \delta_i \Rightarrow |f_i(x) - f_i(x_0)| < \frac{\e}{\sqrt{k}}$. Con $\delta = \min\{\delta_1, \dots, \delta_k\}$ y $\norm{x - x_0} < \delta$,
  \[
    \norm{f(x) - f(x_0)} = \sqrt{\sum_{i=1}^{k} {\bigl( f_i(x) - f_i(x_0) \bigr)}^2} < \sqrt{k \cdot \frac{\e^2}{k}} = \e. \qedhere
  \]
\end{proof}

\section{Continuidad global}

\begin{definition}[Continuidad global]
  Sea $A \subset \R^n$. Una función $f\colon A \to \R^k$ es \emph{continua} si es continua en todo punto de $A$.
\end{definition}

\begin{prop}
  Sean $A \subset \R^n$ y $f\colon A \to \R^k$. Entonces $f$ es continua si y sólo si para toda sucesión ${(x_m)}_{m \in \N} \subset A$ con $x_m \to x \in A$ se tiene $\lim_{m \to \infty} f(x_m) = f(x)$.
\end{prop}

\begin{proof}
  Es la caracterización por sucesiones de la continuidad en un punto, aplicada en cada $x \in A$.
\end{proof}

\begin{prop}
  Sean $A \subset \R^n$ y $f\colon A \to \R^k$. Son equivalentes:
  \begin{enumerate}[label=(\arabic*)]
    \item $f$ es continua.
    \item Para todo abierto $V \subset \R^k$, $f^{-1}(V)$ es abierto relativo a $A$.
    \item Para todo cerrado $F \subset \R^k$, $f^{-1}(F)$ es cerrado relativo a $A$.
  \end{enumerate}
\end{prop}

\begin{proof}
  (1) $\Rightarrow$ (2) Sea $V \subset \R^k$ abierto y $p \in f^{-1}(V)$. Existe $\e > 0$ con $B_{\e}(f(p)) \subset V$, y por la continuidad de $f$ en $p$ existe $\delta_p > 0$ tal que $x \in A$, $\norm{x - p} < \delta_p$ implica $f(x) \in B_{\e}(f(p)) \subset V$, es decir, $A \cap B_{\delta_p}(p) \subset f^{-1}(V)$. Entonces
  \[
    f^{-1}(V) = A \cap \bigcup_{p \in f^{-1}(V)} B_{\delta_p}(p),
  \]
  pues $\supset$ vale por lo anterior y $\subset$ porque $p \in B_{\delta_p}(p)$. La unión es un abierto de $\R^n$, de modo que $f^{-1}(V)$ es abierto relativo a $A$.

  (2) $\Rightarrow$ (3) Sea $F \subset \R^k$ cerrado. Entonces $\R^k \setminus F$ es abierto y $f^{-1}(\R^k \setminus F) = A \cap \tilde{U}$ para algún abierto $\tilde{U} \subset \R^n$. Como $f^{-1}(F) = A \setminus f^{-1}(\R^k \setminus F)$, resulta $f^{-1}(F) = A \cap (\R^n \setminus \tilde{U})$, que es cerrado relativo a $A$.

  (3) $\Rightarrow$ (1) Sea $x \in A$ y $V$ un entorno de $f(x)$; entonces $f(x) \in V^{\circ}$. Como $\R^k \setminus V^{\circ}$ es cerrado, $f^{-1}(\R^k \setminus V^{\circ}) = A \cap \tilde{F}$ con $\tilde{F}$ cerrado en $\R^n$. Luego $f^{-1}(V^{\circ}) = A \setminus f^{-1}(\R^k \setminus V^{\circ}) = A \cap (\R^n \setminus \tilde{F})$, donde $\R^n \setminus \tilde{F}$ es abierto y contiene a $x$. Entonces $U = A \cap (\R^n \setminus \tilde{F})$ es un entorno de $x$ relativo a $A$ con $f(U) \subset V^{\circ} \subset V$. Por la caracterización de la continuidad en un punto, $f$ es continua en $x$.
\end{proof}

\begin{corollary}
  Sean $f, g\colon A \to \R^k$ continuas.
  \begin{enumerate}[label=(\arabic*)]
    \item El conjunto $\{ x \in A : f(x) \neq g(x) \}$ es abierto relativo a $A$ y $\{ x \in A : f(x) = g(x) \}$ es cerrado relativo a $A$.
    \item Si $k = 1$, $\{ x \in A : f(x) > 0 \}$ es abierto relativo a $A$.
  \end{enumerate}
\end{corollary}

\begin{proof}
  (1) La función $h = f - g\colon A \to \R^k$ es continua, $\R^k \setminus \{0\}$ es abierto y $\{0\}$ es cerrado. Luego $h^{-1}(\R^k \setminus \{0\}) = \{ x \in A : f(x) \neq g(x) \}$ es abierto relativo a $A$ y $h^{-1}(\{0\}) = \{ x \in A : f(x) = g(x) \}$ es cerrado relativo a $A$.

  (2) Como $f$ es continua y $(0, +\infty)$ es abierto, $f^{-1}\bigl( (0, +\infty) \bigr) = \{ x \in A : f(x) > 0 \}$ es abierto relativo a $A$.
\end{proof}

\begin{note}
  La imagen de un conjunto abierto (o cerrado) por una función continua puede no ser abierta ni cerrada. Por ejemplo, $f(x) = \frac{1}{1 + x^2}$ es continua, $\R$ es abierto y cerrado, y $f(\R) = (0, 1]$, que no es abierto ni cerrado.
\end{note}


% TODO: revisar clase 13
\clase{15/10/2024}{Gráficas continuas y convexidad}
\section{Gráficas continuas}

\begin{definition}
  Sea \(A \subset \R, f: A \to \R^k\) el gráfico de \(f\) es el subconjunto de \(A \times \R^k\) dado por \(G_f = \{ (x, f(x)) : x \in A \} \).
\end{definition}

\begin{prop}
  Si \(f\) es continua, \(G_f\) es cerrado en \(A \times \R^k\).
  \begin{proof}
    Si \({(p_n)}_{n \in \N} \subset G_f\) y \(p_n \to p \in A \times \R^k\). Para cada \(n \in \N \) podemos poner \(p_n = (x_n, y_n)\) con \(y_n = f(x_n)\). Si \(p = (x, y)\) tiene que ser \(x_n \to x, y_n \to y\). Como \(f\) es continua \(f(x_n) \to f(x)\), pero como \(y_n = f(x_n) \to y = f(x) \Rightarrow p = (x, y) \in G_f\).
  \end{proof}
\end{prop}

\begin{prop}
  Si \(G_f\) es cerrado en \(A \times \R^k\) y \(f\) es acotada \(\Rightarrow f\) es continua.
  \begin{proof}
    Como \(f\) es acotada \(\exists k : f(x) \in B_k(0) \subset \R^k, \forall x \in A\). Sea \(x_0 \in A\) queremos ver que \(f\) es continua en \(x_0\).
    Sea \({(x_n)}_{n \in \N} \subset A\) tal que \(x_n \to x_0\) y supongamos que \(f(x_n) \not \to f(x_0) \Rightarrow (\exists \e > 0), h: \N \to \N \) estrictamente creciente tal que
    \( \|f(x_{h(n)}) - f(x_0) \| \geq \e, \forall n \in \N \). \\
    Para cada \(n \in \N \) consideremos \(p_n = (x_{h(n)}, f(x_{h(n)})) \in G_f\). Como \(x_{h(n)} \) converge, es acotada, \(\exists L > 0 : x_{h(n)} \in B_L(0), \forall n \in \N \). Luego \({(p_n)}_{n \in \N} \) es acotada pues toma valores en \(B_L(0)\) o \(B_k(0)\), por lo que existe \(\phi \) estrictamente creciente tal que \({(p_{\phi_{h(n)}})}_{n \in \N} \) converge a \(p \in \R^n \times \R^k\).
    Supongamos que \(p = (x_0^{\prime}, y) \Rightarrow x_{\phi_{h(n)}} \to x_0^{\prime} \) y \(f(x_{\phi_{h(n)}}) \to y\). Como \(x_{\phi_{h(n)}} \) es subsubsucesión de \(x_n \to x_0^{\prime} = x_0 \in A\). Se tiene que \(\lim_{n \to +\infty} p_{\phi_{h(n)}} \in A \times \R^k\). Como estamos suponiendo que \(G_f\) es cerrado entonces \(x_0, y \in G_f\), es decir \(y = f(x_0)\), pero entonces \(f(x_{\phi_{h(n)}}) \to f(x_0)\). Absurdo!
  \end{proof}
\end{prop}

\begin{prop}
  Sea \(A \subset \R^n, D \subset A : A \subset \overline{D} \) (\(D\) es denso en \(A\)) y sean \(f, g : A \to \R^k\) continuas si \(f|_D = g|_D \Rightarrow f=g\).

  \begin{proof}
    Sea \(x \in A\), como \(A \subset \overline{D} (\exists {(x_n)}_{n \in \N}) \subset D : x_n \to x\). Como \(f, g\) son continuas en \(x\), \(f(x_n) \to f(x)\), \(g(x_n) \to g(x)\). Como \(f|_D = g|_D \Rightarrow f(x_n) = g(x_n), \forall n \in \N \therefore f(x) = g(x)\).
  \end{proof}
\end{prop}

\clearpage

\begin{theorem}[Weiertrass]
  Sea \(K \subset \R^n\) compacto y no vacío. Una función continua \(f: K \to \R \) es acotada superiormente y más aún \(\exists y \in K : f(x) \leq f(y) (\forall x \in K)\).
  \begin{proof}
    Si \(f\) no es acotada superiormente para cada \(n \in \N, \exists x_n \in K : f(x_n) \geq n\). Como \({(x_n)}_{n \in \N} \subset K\) que es compacto \(\exists \phi: \N \to \N \) estrictamente creciente tal que \(\lim_{n \to +\infty} x_{\phi(n)} = \alpha \in K\). Como \(f\) es continua en \(\alpha, f(x_n) \to f(\alpha)\) y como es convergente, tiene que ser acotada. Absurdo, pues \(f(x_{\phi(n)}) \geq h(n) \geq n (\forall n \in \N) \therefore f\) es acotada superiormente. \\
    Sea ahora \(F = \{ f(x) : x \in K\} \) este conjunto está acotado superiormente y es no vacío, \(\exists s = \sup(F) \Rightarrow \exists {(y_n)}_{n \in \N} \subset K : f(y_n) \to s\). Como \(K\) es compacto \(\exists \phi : \N \to \N \) estrictamente creciente tal que \(y_{\phi(n)} \to y \in K\). Como \(f\) es continua en \(y \Rightarrow f(y_{\phi(n)}) \to f(y) \Rightarrow f(y) = s\). Luego \(f(y)\) es cota superior de \(F\) y \(f(y) \in F \therefore \) alcanza valor máximo.
  \end{proof}
\end{theorem}

\begin{note}
  Vale para acotada inferiormente y valor mínimo.
\end{note}

\begin{theorem}
  Sea \(K \subset \R^n\) compacto y \(f\) una función continua \(\Rightarrow f(K)\) es compacto.
  \begin{proof}
    Sea \({(y_n)}_{n \in \N} \) que toma valores en la imagen \(f(K)\). Para cada \(n \in \N \exists x_n \in K : y_n = f(x_n)\). Como \({(x_n)}_{n \in \N} \subset K\) y \(K\) es compacto, \(\exists \phi : \N \to \N \) estrictamente creciente tal que \(x_{\phi(n)} \to x \in K\). Como \(f\) es continua en \(x \Rightarrow f(x_{\phi(n)}) \to f(x)\). Como \(f(x_{\phi(n)}) = f(y_n)\), probamos que toda sucesión contenida en \(f(K)\) tiene una subsucesión convergente con límite contenido en \(f(K) \therefore f(K)\) es compacto.
  \end{proof}
\end{theorem}

\begin{theorem}[Valor intermedio]
  Sea \(f: [a, b] \to \R \) continua entonces \(\forall y\) entre \(f(a)\) y \(f(b), \exists x \in [a, b] : f(x) = y\).

  \begin{proof}
    Si \(f(a) = f(b)\) no hay nada que probar. \\
    Supongamos que \(f(a) < f(b) \Rightarrow \) Sea \(y \in (f(a), f(b))\) y consideremos a \(U = \{ x \in [a, b] : f(x) < y \} \neq \varnothing \). Luego \(U\) es acotado y entonces \(\exists \alpha = \sup(U)\) y \(\alpha \in [a, b] \Rightarrow \exists {(x_n)}_{n \in \N} \subset U : x_n \to \alpha \Rightarrow f(x_n) \to f(\alpha)\) pues \(f\) es continua. Como \((\forall n \in \N)(x_n \in U)(f(x_n) < y) \Rightarrow \lim_{n \to +\infty} f(x_n) \leq y\). Quiero ver que \(f(\alpha) = y\). Si fuese \(f(\alpha) < y \Rightarrow \alpha \in U\). Además, como \(y < f(b)\) y \(\alpha \in [a, b]\) tiene que ser \(\alpha < b\). Como \(f\) es continua, \(U\) es abierto relativo a \([a, b]\) así que \(\exists r > 0 : B_r(\alpha) \cap [a, b] \subset U\). Si \(s = \dfrac{1}{2} \min(r, b-\alpha) \Rightarrow \alpha + s \in B_r(\alpha)\) y \(a \leq \alpha + s \leq b \Rightarrow \alpha + s \in U \Rightarrow f(\alpha + s) < y\) Absurdo pues \(\alpha = \sup(U) \therefore f(\alpha) = y\).
  \end{proof}
\end{theorem}

\section{Convexidad}

\begin{definition}[Convexidad]
  Decimos que \(C \subset \R^n\) es convexo si siempre que \(x, y \in C\) y \(t \in [0, 1]\) se tiene que \((1-t)x+ty \in C\). Geométricamente, \(C\) es convexo si dados dos puntos cualquiera en \(C\), el segmento que los une también está contenido en \(C\).
\end{definition}

\begin{eg}
  Si \(x \in \R^n, r > 0, B_r(x)\) es convexo.
  \begin{proof}
    \(y, z \in B_r(x), t \in [0,1] \Rightarrow \) \begin{align*}
      \|(1-t)y+tz - x\| & = \| (1-t)y + tz - (1-t)x - tx \|                \\
                        & \leq \|(1 -t) (y-x) \| + \|t(z-x)\|              \\
                        & \leq (1-t) \|y-x\| + t \|z-x\| < (1-t)r + tr = r
    \end{align*} \(\therefore (1-t)y + tz \in B_r(x) (\forall t \in [0, 1])\).
  \end{proof}
\end{eg}

\begin{prop}
  Sea \(C \subset \R^n\) convexo y sea \(f : C \to \R \) continua \(\Rightarrow \) si \(a, b \in C\) e \(y\) está entre \(f(a)\) y \(f(b) \Rightarrow \exists \alpha \in C : f(\alpha) = y\).
  \begin{proof}
    Sea \(\sigma: [0, 1] \to \R^n, \sigma(t) = (1-t)a+tb \Rightarrow \sigma \) es continua y \(\sigma([0, 1]) \subset C \Rightarrow g = f \circ \sigma: [0, 1] \to \R \) es continua y \(g(0) = f(a), g(1) = f(b) \Rightarrow y\) está entre \(g(a)\) y \(g(b)\) así que \(\exists \beta \in [0, 1] : g(\beta) = y\). Si llamo \(\alpha = \sigma(\beta) \Rightarrow f(\alpha) = f(\sigma(\beta)) = g(\beta) = y\).
  \end{proof}
\end{prop}

\begin{definition}
  Decimos que \(C \subset \R^n\) es arco-conexo si \(\forall x, y \in \C, \exists \sigma: [0, 1] \to \R^n\) continua tal que \(Im(\sigma) \subset C\) y \(\sigma(0) = x, \sigma(1) = y\). Geométricamente, puedo unir dos puntos del conjunto por una curva continua contenida en él.
\end{definition}

\begin{note}
  Todo conjunto convexo es arco-conexo en \(\R^n\).
\end{note}

\begin{prop}
  Sea \(A \subset \R^n\) arco-conexo y sea \(f: A \to \R \) continua, si \(a, b \in A\) e \(y\) está entre \(f(a)\) y \(f(b) \Rightarrow \exists \alpha \in A : f(\alpha) = y\).
  \begin{proof}
    Como \(A\) es arco-conexo \(\exists \sigma: [0, 1] \to \R^n\) continua tal que \(\sigma(0) = a, \sigma(1) = b\). Si \(g = f \circ \sigma: [0, 1] \to \R \Rightarrow g\) es continua y \(g(0) = f(a), g(1) = f(b)\). Así que \(\exists \beta \in [0, 1] : g(\beta) = y \Rightarrow \alpha = \sigma(\beta) \Rightarrow f(\alpha) = f(\sigma(\beta)) = g(\beta) = y\).
  \end{proof}
\end{prop}

\begin{note}
  Supongamos que \(A \subset \R^n\) y que \(f: A \to \R^n\) es una función continua para la que no vale el Teorema de valor intermedio.\(\exists a, b \in A\) e \(y\) entre \(f(a)\) y \(f(b) : \nexists \alpha \in A : f(\alpha) = y \Rightarrow \exists U = \{ x \in A : f(x) = y \}, V = \{ x \in A : f(x) > y \} \). Son abiertos en \(A\), no vacíos, disjuntos y \(A = U \cup V\).
\end{note}

\begin{definition}[Conexo]
  \(A \subset \R^n\) es conexo si \(\nexists U, V\) abiertos en \(A\), no vacíos tales que \(U \cap V = \varnothing \) y \(U \cup V = A\).
\end{definition}

\begin{prop}
  \(A \subset \R^n, A \neq \varnothing \) es conexo \(\iff \) el único subconjunto no vacío \(U\) de \(A\) que es simultáneamente abierto y cerrado en \(A\) es \(A\) mismo.
  \begin{proof}
    Si \(A\) no es conexo \(\exists U, V\) abiertos, no vacíos en \(A\) tales que \(U \cap V = \varnothing \) y \(U \cup V = A\). Luego \(U \neq \varnothing, U \neq A\) y \(U\) es abierto y cerrado en \(A\). Por otro lado si \(A\) es conexo y \(U\) es un subconjunto \(U \neq\varnothing \) de \(A\) abierto y cerrado en \(A \Rightarrow V = A - U\) es abierto en \(A\) y cumple que \(U \cap V = \varnothing, U \cup V = A, V = \varnothing \Rightarrow U = A\).
  \end{proof}
\end{prop}

\begin{theorem}
  \(A \subset \R^n\) conexo y \(f: A \to \R^k\) continua entonces \(f(A)\) es conexo.
  \begin{proof}
    \(U, V\) dos abiertos de \(f(A) : f(A) = U \cup V, U \cap V = \varnothing \Rightarrow \exists U_0, V_0 \subset \R^k\) abiertos tales que \(U = U_0 \cap f(A), V = V_0 \cap f(A)\). Como \(f\) es continua, \(f^{-1}(V_0), f^{-1}(U_0)\) son abiertos relativos en \(A\). \(U_0 \cap V_0 = \varnothing \). Además \(f^{-1}(U) \cap f^{-1}(V) =\varnothing \) y \(f^{-1}(U) \cup f^{-1}(V) = A\). Pero como \(A\) es conexo o \(U = \varnothing \) o \(V = \varnothing \therefore f(A)\) es conexo.
  \end{proof}
\end{theorem}

\begin{prop}
  \(A\) una familia de subconjuntos conexos de \(\R^n\). Si cada vez que \(B\text{, } C \in A\) vale que son no disjuntos \(\Rightarrow \bigcup_{a \in A} a\) es conexo.
  \begin{proof}
    Consideremos \( \mathcal{A} = \bigcup_{a \in A} a \).
    Si \( U\text{, } V \) son dos abiertos, disjuntos no vacíos tales que \( A \subseteq U \cup V \) quiero ver que \( A = U \) y \( V = \varnothing \).
    Dado \(a \in A\) como \( a \subseteq \mathcal{A} \subseteq U \cup V \) debe ser (s.p.d.g) \( a = U \), pues \( a \) es conexo. 
    Si \( a = U \) tenemos que \( a = U \quad \forall a \in A \) pues \( a \) era genérico \( \therefore \mathcal{A} \) es conexo, pues \( \mathcal{A} = U \).
  \end{proof}
\end{prop}


\clase{18/10/2024}{Conexión y continuidad uniforme}
\section{Conjuntos conexos y arco-conexos}

\begin{theorem}
  Todo subconjunto no vacío y arco-conexo de $\R^n$ es conexo.
\end{theorem}

\begin{proof}
  Sea $A \subset \R^n$ arco-conexo y no vacío, y sea $a_0 \in A$. Para cada $a \in A$ existe una función continua $\sigma_a\colon [0, 1] \to \R^n$ con valores en $A$ tal que $\sigma_a(0) = a_0$ y $\sigma_a(1) = a$. Como $[0, 1]$ es conexo y $\sigma_a$ es continua, $\sigma_a([0, 1])$ es conexo; además contiene a $a_0$ y a $a$ y está contenido en $A$. Luego $A = \bigcup_{a \in A} \sigma_a([0, 1])$ es unión de conexos que contienen todos a $a_0$, en particular son no disjuntos de a pares, y por lo tanto $A$ es conexo.
\end{proof}

\begin{corollary}
  $\R^n$ es conexo. En particular, los únicos subconjuntos de $\R^n$ que son abiertos y cerrados a la vez son $\varnothing$ y $\R^n$.
\end{corollary}

\begin{proof}
  $\R^n$ es convexo, luego arco-conexo, luego conexo. La segunda afirmación es la caracterización de los conexos por subconjuntos abiertos y cerrados.
\end{proof}

\begin{theorem}
  Un subconjunto abierto y no vacío de $\R^n$ es arco-conexo si y sólo si es conexo.
\end{theorem}

\begin{proof}
  ($\Rightarrow$) Es el teorema anterior.

  ($\Leftarrow$) Sea $A \subset \R^n$ abierto, no vacío y conexo, y sea $x_0 \in A$. Sea $U$ el conjunto de los $x \in A$ para los que existe una función continua $\sigma\colon [0, 1] \to \R^n$ con valores en $A$, $\sigma(0) = x_0$ y $\sigma(1) = x$. Veamos que $U = A$.

  Primero observamos que, si $p \in U$, $q \in A$ y el segmento $\{ (1-s)p + sq : s \in [0, 1] \}$ está contenido en $A$, entonces $q \in U$. En efecto, si $\sigma$ es una curva en $A$ de $x_0$ a $p$, la función
  \[
    \tau(\alpha) = \begin{cases}
      \sigma(2\alpha)                       & \text{si } \alpha \in \left[ 0, \tfrac{1}{2} \right], \\
      (2 - 2\alpha)\, p + (2\alpha - 1)\, q & \text{si } \alpha \in \left[ \tfrac{1}{2}, 1 \right]
    \end{cases}
  \]
  está bien definida, pues ambas expresiones valen $p$ en $\alpha = \frac{1}{2}$, es continua, tiene valores en $A$ y cumple $\tau(0) = x_0$ y $\tau(1) = q$.

  $U$ es abierto: sea $x \in U$. Como $A$ es abierto, existe $r > 0$ con $B_r(x) \subset A$. Si $y \in B_r(x)$, el segmento de $x$ a $y$ está contenido en la bola $B_r(x)$, que es convexa, y por lo tanto en $A$. Luego $y \in U$ y $B_r(x) \subset U$.

  $A \setminus U$ es abierto: sea $u \in A \setminus U$ y $s > 0$ con $B_s(u) \subset A$. Si existiera $y \in B_s(u) \cap U$, el segmento de $y$ a $u$ estaría en $B_s(u) \subset A$ y entonces $u \in U$, absurdo. Luego $B_s(u) \subset A \setminus U$.

  Entonces $A = U \cup (A \setminus U)$ con $U$ y $A \setminus U$ abiertos y disjuntos. Como $A$ es conexo y $x_0 \in U$, debe ser $A \setminus U = \varnothing$, es decir, $U = A$. Esto prueba que $A$ es arco-conexo, pues dados $x, y \in A$ se unen mediante una curva de $x_0$ a $x$ y otra de $x_0$ a $y$; recorriendo la primera al revés y luego la segunda se obtiene una curva de $x$ a $y$ en $A$.
\end{proof}

\section{Continuidad uniforme}

\begin{definition}[Continuidad uniforme]
  Una función $f\colon A \to \R^k$, con $A \subset \R^n$, es \emph{uniformemente continua} si
  \[
    (\forall \e > 0)(\exists \delta > 0) \colon (\forall x, y \in A) \quad \norm{x - y} < \delta \Rightarrow \norm{f(x) - f(y)} < \e.
  \]
\end{definition}

\begin{eg}
  \begin{enumerate}[label=(\arabic*)]
    \item $f(x) = x$ es uniformemente continua: dado $\e > 0$ basta tomar $\delta = \e$.
    \item $f\colon \R \to \R$, $f(x) = \sin x$ es uniformemente continua. Como $f$ es continua en $0$, dado $\e > 0$ existe $\delta' > 0$ tal que $|z| < \delta' \Rightarrow |\sin z| < \frac{\e}{2}$. Sea $\delta = 2\delta'$. Si $|x - y| < \delta$, entonces
          \[
            |\sin x - \sin y| = \left| 2 \sin \frac{x - y}{2} \cos \frac{x + y}{2} \right| \leq 2 \left| \sin \frac{x - y}{2} \right| < 2 \cdot \frac{\e}{2} = \e.
          \]
    \item $f\colon (0, +\infty) \to \R$, $f(x) = \cos \frac{\pi}{x}$ no es uniformemente continua. Sean $x_m = \frac{1}{m}$ e $y_m = \frac{1}{m + 1}$. Entonces $|x_m - y_m| = \frac{1}{m(m+1)} \to 0$, pero $|f(x_m) - f(y_m)| = \bigl| \cos(m\pi) - \cos((m+1)\pi) \bigr| = 2$ para todo $m$. Para $\e = 1$, ningún $\delta > 0$ sirve: existe $m$ con $\frac{1}{m(m+1)} < \delta$ y para ese $m$ se tiene $|x_m - y_m| < \delta$ pero $|f(x_m) - f(y_m)| = 2 \geq \e$.
  \end{enumerate}
\end{eg}

\begin{theorem}
  Si $K \subset \R^n$ es compacto y $f\colon K \to \R^k$ es continua, entonces $f$ es uniformemente continua.
\end{theorem}

\begin{proof}
  Sea $\e > 0$. Como $f$ es continua, para cada $x \in K$ existe $\delta_x > 0$ tal que
  \[
    (\forall z \in K) \quad \norm{z - x} < \delta_x \Rightarrow \norm{f(z) - f(x)} < \frac{\e}{2}.
  \]
  La familia $\mathscr{U} = \{ B_{\delta_x}(x) : x \in K \}$ es un cubrimiento abierto de $K$. Como $K$ es compacto, existe un número de Lebesgue $\delta > 0$ para $\mathscr{U}$: si $x, y \in K$ y $\norm{x - y} < \delta$, existe $z \in K$ con $x, y \in B_{\delta_z}(z)$. Entonces $\norm{f(x) - f(z)} < \frac{\e}{2}$ y $\norm{f(y) - f(z)} < \frac{\e}{2}$, y
  \[
    \norm{f(x) - f(y)} \leq \norm{f(x) - f(z)} + \norm{f(z) - f(y)} < \e.
  \]
  Luego $f$ es uniformemente continua.
\end{proof}

\begin{theorem}[Extensión de funciones uniformemente continuas]
  Sea $A \subset \R^n$ y $f\colon A \to \R^k$ uniformemente continua. Entonces existe una única función continua $\bar{f}\colon \cl{A} \to \R^k$ tal que $\bar{f}(x) = f(x)$ para todo $x \in A$, y además $\bar{f}$ es uniformemente continua.
\end{theorem}

\begin{proof}
  \emph{Paso 1.} Si ${(x_m)} \subset A$ es convergente, entonces ${(f(x_m))}$ es convergente. En efecto, ${(x_m)}$ es de Cauchy. Dado $\e > 0$, sea $\delta > 0$ tal que $\norm{x - y} < \delta \Rightarrow \norm{f(x) - f(y)} < \e$. Existe $m_0$ tal que $\norm{x_m - x_l} < \delta$ para $m, l > m_0$, y entonces $\norm{f(x_m) - f(x_l)} < \e$. Luego ${(f(x_m))}$ es de Cauchy en $\R^k$ y converge.

  \emph{Paso 2.} Si ${(x_m)}, {(y_m)} \subset A$ convergen al mismo punto $x_0$, entonces $\lim f(x_m) = \lim f(y_m)$. En efecto, $\norm{x_m - y_m} \to 0$. Dado $\e > 0$, sea $\delta$ como en la continuidad uniforme; para $m$ grande $\norm{x_m - y_m} < \delta$ y entonces $\norm{f(x_m) - f(y_m)} < \e$. Luego $\norm{f(x_m) - f(y_m)} \to 0$, y los límites, que existen por el paso 1, coinciden.

  \emph{Paso 3: definición de $\bar{f}$.} Dado $x_0 \in \cl{A}$ existe una sucesión ${(x_m)} \subset A$ con $x_m \to x_0$. Definimos $\bar{f}(x_0) = \lim_{m \to \infty} f(x_m)$. Por los pasos 1 y 2, este límite existe y no depende de la sucesión elegida. Si $x \in A$, la sucesión constante $x_m = x$ converge a $x$, de modo que $\bar{f}(x) = f(x)$.

  \emph{Paso 4: $\bar{f}$ es uniformemente continua.} Sea $\e > 0$ y $\delta > 0$ tal que $x, y \in A$, $\norm{x - y} < \delta \Rightarrow \norm{f(x) - f(y)} < \frac{\e}{2}$. Sean $u, v \in \cl{A}$ con $\norm{u - v} < \delta$, y sean ${(u_m)}, {(v_m)} \subset A$ con $u_m \to u$, $v_m \to v$. Entonces $\norm{u_m - v_m} \to \norm{u - v} < \delta$, de modo que $\norm{u_m - v_m} < \delta$ para $m$ grande, y entonces $\norm{f(u_m) - f(v_m)} < \frac{\e}{2}$. Tomando límite cuando $m \to \infty$ se obtiene $\norm{\bar{f}(u) - \bar{f}(v)} \leq \frac{\e}{2} < \e$. En particular, $\bar{f}$ es continua.

  \emph{Paso 5: unicidad.} Sea $g\colon \cl{A} \to \R^k$ continua con $g|_A = f$. Dado $x \in \cl{A}$, sea ${(x_m)} \subset A$ con $x_m \to x$. Por la continuidad de $g$, $g(x_m) \to g(x)$, y por la continuidad de $\bar{f}$, $\bar{f}(x_m) \to \bar{f}(x)$. Como $g(x_m) = f(x_m) = \bar{f}(x_m)$ para todo $m$, por la unicidad del límite $g(x) = \bar{f}(x)$.
\end{proof}

% TODO: Revisar esta sección
\section{Diferenciación}

\begin{definition}
  Sea \(F : A \subset \R^n \to \R^n\) y sea \(p \in A^{\circ} \) si \(T_p : \R^n \to \R^n\) es una transformación lineal que verifica:
  \begin{align*}
    \lim_{x \to p} \dfrac{\| F(x) - F(p) - T_p(x-p) \|}{\|x-p\|} = 0 \Rightarrow
  \end{align*} Decimos que \(F\) es diferenciable en \(p\). A la transformación lineal la llamamos diferencial de \(F\) en \(p\) y lo notamos \(DF_p\).
\end{definition}

Si \(F = (f_1, \ldots, f_n)\) con \(f_i: A \to \R \) \(i \in \{1, \ldots, n\} \), \(F\) es diferenciable en \(p \iff \) cada \(f_i\) es diferenciable en \(p\).
Además la diferencial de \(F\) es la matriz que se obtiene al poner como filas los gradientes de \(f_i\).

\begin{align*}
  DF_p = \begin{pmatrix}
           \nabla f_1(p) \\
           \vdots        \\
           \nabla f_n(p)
         \end{pmatrix} = \begin{pmatrix}
                           \dfrac{\partial f_1}{\partial x_1}(p) & \cdots & \dfrac{\partial f_1}{\partial x_n}(p) \\
                           \vdots                                & \cdots & \vdots                                \\
                           \dfrac{\partial f_n}{\partial x_1}(p) & \cdots & \dfrac{\partial f_n}{\partial x_n}(p)
                         \end{pmatrix}
\end{align*}

También podemos observar que si \(F\) es diferenciable en \(P\) y \(v \in \R^n \Rightarrow \alpha(t) = p + t \cdot v\).

\begin{align*}
  0 & = \lim_{t \to 0} \dfrac{\| F(p + t \cdot v) - F(p) - DF_p(t \cdot v) \|}{|t|}                 \\
    & = \lim_{t \to 0} \| \dfrac{F(p + tv) - F(p)}{t} - DF_p(v) \|                                  \\
    & \iff DF_p(v) = \lim_{t \to 0} \dfrac{F(p+t \cdot v) - F(p)}{t} \text{  } (\forall v \in \R^n)
\end{align*}

\begin{note}
  Si \(G, H: A \subset \R^n \to \R^n\) son diferenciables podemos considerar su producto escalar como \(F(x) = \langle G(x), H(x) \rangle \) que es diferenciable.

  Además:
  \begin{align*}
     & \lim_{t \to 0} \dfrac{F(p + tx) - F(p)}{t} =                                                                                                                 \\
     & \lim_{t \to 0} \dfrac{1}{t} ( \langle G(p + tx), H(p+tx) \rangle - \langle G(p), H(p) \rangle )                                                              \\
     & \lim_{t \to 0} \dfrac{1}{t} ( \langle G(p+tx), H(p+tx) \rangle - \langle G(p), H(p+tx) \rangle + \langle G(p), H(p+tx) \rangle - \langle G(p), H(p) \rangle) \\
     & \lim_{t \to 0} \dfrac{1}{t} (\langle G(p+tx) - G(p), H(p + tx) \rangle + \langle G(p), H(p+tx) - H(p) \rangle)                                               \\
     & \langle DG_p(x), H(p) \rangle + \langle G(p), DH_p \rangle \Rightarrow                                                                                       \\
     & D\langle G, H \rangle_p = \langle DG_p(x), H(p) \rangle + \langle G(p), DH_p(x) \rangle                                                                      \\
  \end{align*}
\end{note}

\begin{lemma}
  Si \(T: \R^n \to \R^n\) es transformación lineal, \(\exists c > 0: \|T(x)\| < c \|x\| \) \((\forall x \in \R^n)\). Al mínimo de estas constantes lo llamamos \( \|T\|_{\infty} \). Vale que \( \|T\|_{\infty} = \max_{\|x\| \leq 1} \|T(x)\| \) en particular
  \( \|T\| \leq \|T\|_{\infty} \|x\| \).
  \begin{proof}
    \(T\) es lineal, luego \(T\) es continua. Tomar norma también es continua así que \(x \mapsto \|Tx\| \) es continua. Como \(\overline{B_1(0)} \) es compacto la función alcanza máximo. Digamos \( \|T\|_{\infty} = \max_{x \in B_r(0)} \|Tx\| \Rightarrow x \neq 0\) \\
    \( \|T(\dfrac{x}{\|x\|})\| \leq \|T\|_{\infty} \Rightarrow \|Tx\| \leq \|x\| \|T\|_{\infty} \) por como la elegimos es la mejor cota.
  \end{proof}
\end{lemma}

\begin{lemma}
  La función \( \| \cdot \|: \R^{n \times m} \to \R \) es una norma.
  \begin{enumerate}
    \item \( \|T\|_{\infty} > 0\), \( \|T\|_{\infty} = 0 \iff T = 0\).
    \item \( \|\lambda \| \|T\|_{\infty} = \|\lambda T\|_{\infty} \).
    \item \( \| T + S\| \leq \|T\|_{\infty} + \|S\|_{\infty} \).
  \end{enumerate}
  \begin{proof}
    \begin{enumerate}
      \item \( \|T\|_{\infty} > 0\) es trivial. Si \(T = 0 \Rightarrow Tx = 0 \Rightarrow \|Tx\|_{\infty} = 0\) y si \( \|T\|_{\infty} \neq 0 \Rightarrow Tx=0 (\forall x \in \overline{B_1(0)}) \) pues por linealidad si
            \begin{align*}
              x \neq 0 \Rightarrow Tx = T(\dfrac{x}{\|x\|} \cdot \|x\|) = \|x\| T(\dfrac{x}{x}) = 0 \therefore T = 0
            \end{align*}
      \item Trivial
      \item \(S + T \in \R^{n \times m} \) dado \(x \in \overline{B_1(0)} \) \begin{align*}
              \|(S+T)(x)\| & = \| Tx + Sx \| \leq \|Tx\| + \|Sx\| \leq  \|x\| \cdot (\|Tx\|_{\infty} + \|Sx\|_{\infty})                   \\
                           & \leq \|Tx\|_{\infty} + \|Sx\|_{\infty} \therefore \|(T+S)(x)\|_{\infty} \leq \|T\|_{\infty} + \|S\|_{\infty}
            \end{align*}
    \end{enumerate}
  \end{proof}
\end{lemma}

\begin{lemma}
  Si \(F: A \subset \R^n \to \R^n\) es diferenciable con \(p \in A^{\circ} \), \(\exists R > 0\) y \(c_r > 0 : x \in B_r(p) \Rightarrow \| F(x) - F(p) \| \leq c_r \| x - p\| \)
  Esto implica la continuidad uniforme en \(B_r(p)\).
  \begin{proof}
    Como \(F\) es diferenciable en \(p\), \(\exists B_r(p)\) donde el cociente de diferenciabilidad es menor que \(1\).\begin{align*}
      \dfrac{\|F(x) - F(p)\|}{\|x - p \|} & \leq \dfrac{\|F(x) - F(p) - DF_p(x-p) \|}{\|x-p\|} + \dfrac{\| DF_p(x-p) \|}{\|x-p\|} \\
                                          & < 1 + \dfrac{\|DF_p(x-p)\|_{\infty} \|x-p\|}{\|x-p\|} = 1 + \|DF_p(x-p)\|_{\infty}
    \end{align*}
    Si \(x \in B_r(p) \Rightarrow \). Luego \(c_r = 1 + \|DF_p(x-p)\|_{\infty} \Rightarrow \|F(X) - F(p)\| \leq c_r \|x-p\| \) \((\forall x \in B_r(p))\).
  \end{proof}
\end{lemma}


\clase{25/10/2024}{Teorema de la función inversa}
\section{Preliminares}

Podíamos definir \(\| T(x) \| \le \| T \|_{\infty} \cdot \| x \| \quad \forall x \in \mathbb{R}^n\).

\begin{theorem}[Regla de la Cadena]
    Sea \(A \subset \mathbb{R}^n\) abierto, \(F: A \subset \mathbb{R}^n \to \mathbb{R}^m\), y \(G: B \subset \mathbb{R}^m \to \mathbb{R}^k\) con \(F(A) \subset B\). Si \(F\) es diferenciable en \(p \in A\) y \(G\) es diferenciable en \(F(p) \in B\), entonces \(G \circ F\) es diferenciable en \(p\) y además:
    \[ D(G \circ F)_p = DG_{F(p)} \cdot DF_p \]
\end{theorem}

\subsection{Teoremas de Análisis I}

\begin{theorem}[Teorema de Fermat]
    Sea \(f: [a,b] \to \mathbb{R}\). Supongamos que \(c \in (a,b)\) es extremo local de \(f\). Si \(f\) es derivable en \(c\), entonces \(f'(c) = 0\).
\end{theorem}

\begin{theorem}[Teorema de Rolle]
    Si \(f\) es continua en \([a,b]\), derivable en \((a,b)\) y \(f(a) = f(b)\), entonces \(\exists c \in (a,b)\) tal que \(f'(c) = 0\).
\end{theorem}

\begin{theorem}[Teorema de Lagrange]
    Sea \(f\) continua en \([a,b]\) y derivable en \((a,b)\). Entonces \(\exists c \in (a,b)\) tal que:
    \[ f'(c) = \frac{f(b) - f(a)}{b - a} \]
\end{theorem}

\subsection{Extensión a \(\mathbb{R}^n\)}

\begin{theorem}[Lagrange en \(\mathbb{R}^n\)]
    Sea \(f: B_r(p) \subset \mathbb{R}^n \to \mathbb{R}\) diferenciable. Para todo \(a, b \in B_r(p)\), en el segmento que une \(a\) con \(b\), existe \(c\) tal que:
    \[ f(b) - f(a) = \nabla f(c) \cdot (b - a) \]
\end{theorem}

\begin{theorem}[Teorema de Fermat en \(\mathbb{R}^n\)]
    Sea \(f: A \subset \mathbb{R}^n \to \mathbb{R}\) con \(A\) abierto. Si \(p \in A\) es extremo local de \(f\), entonces \(\nabla f(p) = 0\).
\end{theorem}

\subsection{Propiedades de Transformaciones Lineales}
\begin{enumerate}
    \item Toda transformación lineal \(T\) es diferenciable y \(DT_p = T\) para todo \(p \in \mathbb{R}^n\).
    \item Si \(T: \mathbb{R}^m \to \mathbb{R}^k\) es lineal y \(F \in C^1\), entonces \(T \circ F\) es diferenciable y \(D(T \circ F)_p = T \cdot DF_p\).
    \item \(T\) es inyectiva si y solo si \(N(T) = \{0\}\).
    \item Si \(T: \mathbb{R}^n \to \mathbb{R}^n\), se tiene que \(T\) es inversible si y solo si es inyectiva (por teorema de la dimensión).
\end{enumerate}

\begin{lemma}
    Sea \(G: A \subset \mathbb{R}^n \to \mathbb{R}^m\) abierto. Si \(G\) es \(C^1\) en \(A\), dados \(p, q \in A\), vale que:
    \[ \| DG_p - DG_q \|_{\infty} \le n \max_{i,j} \left| \frac{\partial g_i}{\partial x_j}(p) - \frac{\partial g_i}{\partial x_j}(q) \right| \]
\end{lemma}

\begin{lemma}
    Sea \(G: B_r(p) \subset \mathbb{R}^n \to \mathbb{R}^m\) diferenciable y supongamos que \(\| DG_x \|_{\infty} \le M\) para todo \(x \in B_r(p)\). Entonces:
    \[ \| G(x) - G(y) \| \le M \| x - y \| \quad \forall x,y \in B_r(p) \]
\end{lemma}

\section{Teorema de la Función Inversa}

\begin{theorem}[Teorema de la Función Inversa]
    Sea \(F: A \subset \mathbb{R}^n \to \mathbb{R}^n\) con \(A\) abierto. Si \(F\) es \(C^1\) en \(A\) y \(DF_p: \mathbb{R}^n \to \mathbb{R}^n\) es una transformación lineal inversible para un \(p \in A\), entonces:
    \begin{enumerate}
        \item \(\exists r > 0\) tal que \(F: B_r(p) \to F(B_r(p))\) es biyectiva y la función inversa es continua en su dominio.
        \item Existen abiertos \(U \subset B_r(p)\) y \(W \subset F(B_r(p))\) tales que \(F^{-1}: W \to U\) es diferenciable.
        \item La diferencial de la función inversa es \(D(F^{-1})_{F(q)} = (DF_q)^{-1}\).
        \item Si además \(F \in C^k(U)\), entonces \(F^{-1} \in C^k(W)\).
    \end{enumerate}
\end{theorem}
\begin{proof}
    1) Supongamos primero que \(DF_p = I\). Usando la continuidad de la diferencial, existe \(r_0 > 0\) tal que \(B_{r_0}(p) \subset A\). Tomamos \(r < r_0\) tal que \(\|DG_x - DF_p\|_\infty \le \frac{1}{2}\) para \(\|x-p\| < r\).
    Definimos \(G(x) = F(x) - x\), entonces \(DG_x = DF_x - I\), y \(\|DG_x\|_\infty \le \frac{1}{2}\). Por el lema anterior, \(\|G(x) - G(y)\| \le \frac{1}{2}\|x-y\|\).
    Desarrollando \(\|x - y\|\) tenemos:
    \[ \|x - y\| = \|x - y - (F(x) - F(y)) + (F(x) - F(y))\| \le \|G(x) - G(y)\| + \|F(x) - F(y)\| \]
    Sustituyendo la cota de \(G\):
    \[ \|x - y\| \le \frac{1}{2}\|x-y\| + \|F(x) - F(y)\| \implies \frac{1}{2}\|x-y\| \le \|F(x) - F(y)\| \]
    Esto prueba que si \(x \neq y\), \(F(x) \neq F(y)\), por ende \(F\) es inyectiva en \(B_r(p)\) y su inversa es continua (es bi-lipschitz).

    2) Demostremos que cubre un abierto. Sea \(S_r(p) = \{x \in \mathbb{R}^n : \|x-p\| = r\} \subset A\). \(S_r(p)\) es compacto. Para un \(y\) en \(W = B_{d/2}(F(p))\) (donde \(d\) es la distancia de \(F(p)\) a \(F(S_r(p))\)), consideramos \(h(x) = \|y - F(x)\|^2\).
    Al ser continua en el compacto \(\overline{B_r}(p)\), alcanza su mínimo en un punto \(p_y\). Evaluando en los bordes se verifica que \(p_y\) debe estar en el interior \(B_r(p)\). Al ser mínimo local diferenciable, evaluando en cualquier dirección \(v\):
    \[ \langle v, DF_{p_y}^t (y - F(p_y)) \rangle = 0 \]
    Como el operador diferencial es inversible por estar en la bola, \(DF_{p_y}^t\) también lo es, deduciéndose que \(y - F(p_y) = 0 \implies F(p_y) = y\).

    3) La diferenciabilidad se deduce mediante una manipulación del límite \(z \to F(q)\) apoyándonos en que \(F\) es diferenciable y que su inversa local acaba de ser probada como continua.

    Finalmente, si \(DF_p \neq I\), tomamos \(\tilde{F} = (DF_p)^{-1} \circ F\) que verifica \(D\tilde{F}_p = I\) y hereda las propiedades demostradas, extendiendo el resultado al caso general.
\end{proof}

\begin{lemma}
    Sean \(A, B: \mathbb{R}^n \to \mathbb{R}^n\) transformaciones lineales.
    \begin{enumerate}
        \item Si \(A\) es inversible y \(\| B - A \|_{\infty} < \frac{1}{\| A^{-1} \|_{\infty}}\), entonces \(B\) es inversible.
        \item El conjunto de operadores lineales inversibles en \(\mathbb{R}^n\) es abierto, y la función \(A \mapsto A^{-1}\) es continua.
    \end{enumerate}
\end{lemma}
\begin{proof}
    1) Llamemos \(\frac{1}{\alpha} = \|A^{-1}\|_\infty\) y \(\beta = \|B-A\|_\infty < \alpha\). Sea \(x \in \mathbb{R}^n\).
    Por un lado, \(\|x\| = \|A^{-1}Ax\| \le \frac{1}{\alpha}\|Ax\| \implies \alpha\|x\| \le \|Ax\|\).
    Además, \(\|Ax\| = \|(A-B+B)x\| \le \|(A-B)x\| + \|Bx\| \le \beta\|x\| + \|Bx\|\).
    Uniendo ambas desigualdades:
    \[ \alpha\|x\| \le \beta\|x\| + \|Bx\| \implies (\alpha-\beta)\|x\| \le \|Bx\| \]
    Como \(\beta < \alpha\), tenemos \(\alpha - \beta > 0\). Si \(x \neq 0\), la desigualdad implica \(\|Bx\| > 0 \implies Bx \neq 0\). Por lo tanto, \(N(B) = \{0\}\), lo que significa que \(B\) es inyectiva. En dimensión finita, esto equivale a ser inversible.

    2) Tomando \(x = B^{-1}y\) en la desigualdad anterior obtenemos \((\alpha - \beta)\|B^{-1}y\| \le \|y\|\), lo que nos da que \(\|B^{-1}\|_\infty \le \frac{1}{\alpha-\beta}\). Para ver la continuidad, usamos la identidad \(B^{-1} - A^{-1} = B^{-1}(A-B)A^{-1}\).
    Tomando normas:
    \[ \|B^{-1} - A^{-1}\|_\infty \le \|B^{-1}\|_\infty \|A-B\|_\infty \|A^{-1}\|_\infty \]
    Esto tiende a 0 a medida que \(B \to A\), probando que la aplicación toma la inversa de manera continua.
\end{proof}

\clase{29/10/2024}{Teorema de la función implícita}
\begin{theorem}[Teorema de la Función Implícita]
    Sea \(S \subset \mathbb{R}^n \times \mathbb{R}^m\) abierto y \(f: S \to \mathbb{R}^m\) de clase \(C^1\).
    Sea \((a,b) \in S\) (con \(a \in \mathbb{R}^n, b \in \mathbb{R}^m\)) tal que \(f(a,b) = 0\) y la matriz
    \[ M = \begin{pmatrix}
            \frac{\partial f_1}{\partial x_{n+1}}(a,b) & \dots  & \frac{\partial f_1}{\partial x_{n+m}}(a,b) \\
            \vdots                                     & \ddots & \vdots                                     \\
            \frac{\partial f_m}{\partial x_{n+1}}(a,b) & \dots  & \frac{\partial f_m}{\partial x_{n+m}}(a,b)
        \end{pmatrix} \]
    es inversible.

    Entonces, existen \(A \subset \mathbb{R}^n\) y \(B \subset \mathbb{R}^m\) abiertos tales que \(A \times B \subset S\), \(a \in A\), \(b \in B\), y existe una única función \(g: A \to B\) de clase \(C^1\) tal que \(g(a) = b\) y
    \[ f(x,g(x)) = 0 \quad \forall x \in A \]
\end{theorem}

\noindent\textbf{Observación:} La hipótesis se puede cambiar por pedir que \(Df(a,b)\) sea sobreyectiva, ya que esto obliga a que la matriz tenga \(m\) columnas linealmente independientes. Se puede suponer que son las últimas \(m\), si no se hace un cambio de variables.

\begin{proof}
    Sea \(F: S \subset \mathbb{R}^n \times \mathbb{R}^m \to \mathbb{R}^n \times \mathbb{R}^m\), definida como \(F(x,y) = (x, f(x,y))\). Como \(f \in C^1\), entonces \(F \in C^1\).
    Evaluando en el punto tenemos \(F(a,b) = (a, f(a,b)) = (a,0)\).

    La matriz jacobiana es:
    \[ DF(a,b) = \begin{pmatrix} I & 0 \\ X & M \end{pmatrix} \]
    la cual es inversible porque \(\det(DF(a,b)) = \det(I) \cdot \det(M) = \det(M) \neq 0\).

    Por el Teorema de la Función Inversa aplicado a \(F\), existen \(V \subset S \subset \mathbb{R}^{n+m}\) abierto y \(W \subset \mathbb{R}^{n+m}\) abierto tales que:
    \((a,b) \in V\), \(F(a,b) = (a,0) \in W\), \(F: V \to W\) es biyectiva y \(F^{-1}: W \to V\) es de clase \(C^1\).

    Escribimos \(V = \tilde{A} \times B\) con \(\tilde{A} \subset \mathbb{R}^n\) y \(B \subset \mathbb{R}^m \implies F^{-1}: W \to \tilde{A} \times B\).
    Podemos escribir \(F^{-1}(x,y) = (F_1^{-1}(x,y), F_2^{-1}(x,y))\) con \(F_1^{-1}: W \to \tilde{A}\) y \(F_2^{-1}: W \to B\).

    Aplicando \(F\) se tiene que:
    \[ (x,y) = F(F^{-1}(x,y)) = F(F_1^{-1}(x,y), F_2^{-1}(x,y)) = (F_1^{-1}(x,y), f(F_1^{-1}(x,y), F_2^{-1}(x,y))) \]
    De aquí se deduce que \(F_1^{-1}(x,y) = x \quad \forall (x,y) \in W\) y además se tiene que:
    \[ f(x, F_2^{-1}(x,y)) = y \quad \forall (x,y) \in W \]
    En definitiva, \(F^{-1}(x,y) = (x, F_2^{-1}(x,y))\) con \(F_2^{-1} \in C^1\).

    Notemos que \(\pi_1: \mathbb{R}^n \times \{0\} \to \mathbb{R}^n\), \(\pi_1(x,0) = x\) es de clase \(C^1\) y biyectiva.
    Tomo \(A = \pi_1( (\mathbb{R}^n \times \{0\}) \cap W )\). Como \(W\) es abierto en \(\mathbb{R}^{n+m}\), la intersección es abierta en \(\mathbb{R}^n \times \{0\}\), y por tanto \(A\) es abierto en \(\mathbb{R}^n\).

    Veamos que \(A \subset \tilde{A}\): si \(x \in A \implies x = \pi_1(x,0)\) con \((x,0) \in W\).
    \[ F^{-1}(x,0) = (x, F_2^{-1}(x,0)) \in \tilde{A} \times B \]
    Por lo tanto, \(x \in \tilde{A}\) y \(F_2^{-1}(x,0) \in B\).
    Luego, restringiendo los dominios adecuadamente, \(g: A \to B\) dada por \(g(x) = F_2^{-1}(x,0)\) es de clase \(C^1\).

    Afirmo que \(g(a) = b\). En efecto:
    \[ F(a,b) = (a,0) \implies (a,b) = F^{-1}(a,0) = (a, F_2^{-1}(a,0)) = (a, g(a)) \]
    Por lo tanto, \(g(a) = b\).

    Nos queda ver que \(f(x,g(x)) = 0\).
    \[ f(x, g(x)) = f(x, F_2^{-1}(x,0)) = 0 \]
    pues vimos que \(f(x, F_2^{-1}(x,y)) = y\), y tomando \(y=0\) se cumple.

    Por último, veamos que \(g\) es única.
    Supongamos que \(\exists \tilde{g}: A \to B\) tal que \(f(x, \tilde{g}(x)) = 0 \quad \forall x \in A\).
    \[ F(x, g(x)) = (x, f(x, g(x))) = (x, 0) = (x, f(x, \tilde{g}(x))) = F(x, \tilde{g}(x)) \]
    Como \(F\) es biyectiva en la región, \(F(x, g(x)) = F(x, \tilde{g}(x)) \implies \tilde{g}(x) = g(x) \quad \forall x \in A\).
\end{proof}

\clase{01/11/2024}{Integral de Riemann}
\section{Integral de Riemann}

Sea $f\colon [a, b] \to \R$ una función acotada, con $a < b$, y sean $m = \inf f([a, b])$ y $M = \sup f([a, b])$. Una \emph{partición} de $[a, b]$ es un subconjunto finito $P \subset [a, b]$ con $a, b \in P$. Escribimos $P = \{ t_0, \dots, t_n \}$ con $a = t_0 < t_1 < \cdots < t_n = b$. Para $i \in \{1, \dots, n\}$ sean $m_i$ y $M_i$ el ínfimo y el supremo de $f$ en $[t_{i-1}, t_i]$. Las \emph{sumas inferior y superior} de $f$ respecto de $P$ son
\begin{align*}
  s(f, P) & = \sum_{i=1}^{n} m_i (t_i - t_{i-1}), \\
  S(f, P) & = \sum_{i=1}^{n} M_i (t_i - t_{i-1}).
\end{align*}
Vale que $m(b - a) \leq s(f, P) \leq S(f, P) \leq M(b - a)$.

\begin{definition}[Refinamiento]
  Si $P$ y $G$ son particiones de un mismo intervalo y $P \subset G$, decimos que $G$ es \emph{más fina} que $P$ (o que refina a $P$).
\end{definition}

\begin{lemma}
  Sea $f\colon [a, b] \to \R$ acotada y sean $P \subset G$ particiones de $[a, b]$. Entonces $s(f, P) \leq s(f, G)$ y $S(f, G) \leq S(f, P)$.
\end{lemma}

\begin{proof}
  Alcanza con probarlo cuando $G$ tiene un único punto más que $P$, pues el caso general se obtiene agregando los puntos de a uno. Sea $P = \{ t_0, \dots, t_n \}$ y $G = P \cup \{ r \}$ con $t_{i-1} < r < t_i$. Sean $m_i = \inf_{[t_{i-1}, t_i]} f$, $m' = \inf_{[t_{i-1}, r]} f$ y $m'' = \inf_{[r, t_i]} f$. Como $[t_{i-1}, r]$ y $[r, t_i]$ están contenidos en $[t_{i-1}, t_i]$, se tiene $m_i \leq m'$ y $m_i \leq m''$. Entonces
  \begin{align*}
    s(f, G) - s(f, P) & = m'(r - t_{i-1}) + m''(t_i - r) - m_i (t_i - t_{i-1})       \\
                      & = (m' - m_i)(r - t_{i-1}) + (m'' - m_i)(t_i - r) \geq 0.
  \end{align*}
  La desigualdad para sumas superiores es análoga.
\end{proof}

\begin{corollary}
  Sea $f\colon [a, b] \to \R$ acotada. Si $P$ y $G$ son particiones de $[a, b]$, entonces $s(f, P) \leq S(f, G)$.
\end{corollary}

\begin{proof}
  Como $P \cup G$ refina a $P$ y a $G$,
  \[
    s(f, P) \leq s(f, P \cup G) \leq S(f, P \cup G) \leq S(f, G). \qedhere
  \]
\end{proof}

\begin{definition}[Integral inferior]
  Si $f\colon [a, b] \to \R$ es acotada, su \emph{integral inferior} es
  \[
    \lowint_a^b f(x) \,\dd x = \sup_P s(f, P),
  \]
  donde el supremo se toma sobre todas las particiones $P$ de $[a, b]$. Es decir, $s(f, P) \leq \lowint_a^b f(x) \,\dd x$ para toda partición $P$, y dado $\e > 0$ existe $P$ con $\lowint_a^b f(x) \,\dd x \leq s(f, P) + \e$.
\end{definition}

\begin{definition}[Integral superior]
  Análogamente, la \emph{integral superior} de $f$ es
  \[
    \upint_a^b f(x) \,\dd x = \inf_P S(f, P).
  \]
\end{definition}

Por el corolario anterior, $\lowint_a^b f(x) \,\dd x \leq \upint_a^b f(x) \,\dd x$.

\begin{lemma}
  Sea $a < c < b$. Si consideramos sólo las particiones de $[a, b]$ que contienen a $c$, los valores de las integrales superior e inferior no cambian.
\end{lemma}

\begin{proof}
  El supremo sobre todas las particiones es mayor o igual que el supremo sobre las que contienen a $c$. Recíprocamente, dada una partición $P$ cualquiera, $P' = P \cup \{c\}$ la refina y contiene a $c$, y $s(f, P) \leq s(f, P')$. Luego el supremo sobre las particiones que contienen a $c$ es mayor o igual que $s(f, P)$ para toda $P$, y entonces ambos supremos coinciden. Para la integral superior se razona igual con ínfimos.
\end{proof}

\begin{theorem}
  Sean $a < c < b$ y $f\colon [a, b] \to \R$ acotada. Entonces
  \begin{align*}
    \lowint_a^b f(x) \,\dd x & = \lowint_a^c f(x) \,\dd x + \lowint_c^b f(x) \,\dd x, \\
    \upint_a^b f(x) \,\dd x  & = \upint_a^c f(x) \,\dd x + \upint_c^b f(x) \,\dd x.
  \end{align*}
\end{theorem}

\begin{proof}
  Sean $A$ y $B$ los conjuntos de sumas inferiores de $f|_{[a, c]}$ y $f|_{[c, b]}$, respectivamente. Las particiones de $[a, b]$ que contienen a $c$ son exactamente las uniones $P_1 \cup P_2$ de una partición $P_1$ de $[a, c]$ y una partición $P_2$ de $[c, b]$, y $s(f, P_1 \cup P_2) = s(f|_{[a, c]}, P_1) + s(f|_{[c, b]}, P_2)$. Por lo tanto $A + B = \{ x + y : x \in A, \ y \in B \}$ es el conjunto de las sumas inferiores de $f$ respecto de particiones de $[a, b]$ que contienen a $c$. Por el lema anterior,
  \[
    \lowint_a^b f(x) \,\dd x = \sup(A + B) = \sup A + \sup B = \lowint_a^c f(x) \,\dd x + \lowint_c^b f(x) \,\dd x.
  \]
  Para las sumas superiores el razonamiento es análogo, usando ínfimos.
\end{proof}

\begin{prop}
  Sean $f, g\colon [a, b] \to \R$ acotadas. Entonces
  \[
    \lowint_a^b f \,\dd x + \lowint_a^b g \,\dd x \leq \lowint_a^b (f + g) \,\dd x \leq \upint_a^b (f + g) \,\dd x \leq \upint_a^b f \,\dd x + \upint_a^b g \,\dd x.
  \]
\end{prop}

\begin{proof}
  Sea $P$ una partición. En cada $[t_{i-1}, t_i]$ valen $m_i(f + g) \geq m_i(f) + m_i(g)$ y $M_i(f + g) \leq M_i(f) + M_i(g)$, de donde
  \[
    s(f, P) + s(g, P) \leq s(f + g, P), \qquad S(f + g, P) \leq S(f, P) + S(g, P).
  \]
  Dadas dos particiones $P$ y $G$, usando el refinamiento $P \cup G$,
  \[
    s(f, P) + s(g, G) \leq s(f, P \cup G) + s(g, P \cup G) \leq s(f + g, P \cup G) \leq \lowint_a^b (f + g) \,\dd x.
  \]
  Tomando supremo sobre $P$ y sobre $G$ se obtiene la primera desigualdad. La tercera se prueba de la misma forma con sumas superiores e ínfimos, y la del medio ya fue observada.
\end{proof}

\begin{prop}
  Sea $f\colon [a, b] \to \R$ acotada y $c \in \R$.
  \begin{itemize}
    \item Si $c \geq 0$, $\lowint_a^b c f \,\dd x = c \lowint_a^b f \,\dd x$ y $\upint_a^b c f \,\dd x = c \upint_a^b f \,\dd x$.
    \item Si $c < 0$, $\lowint_a^b c f \,\dd x = c \upint_a^b f \,\dd x$ y $\upint_a^b c f \,\dd x = c \lowint_a^b f \,\dd x$.
  \end{itemize}
\end{prop}

\begin{proof}
  Si $c \geq 0$, $m_i(cf) = c\, m_i(f)$ y $M_i(cf) = c\, M_i(f)$, y si $c < 0$, $m_i(cf) = c\, M_i(f)$ y $M_i(cf) = c\, m_i(f)$, por las propiedades del supremo y del ínfimo. Luego $s(cf, P) = c\, s(f, P)$ y $S(cf, P) = c\, S(f, P)$ si $c \geq 0$, y $s(cf, P) = c\, S(f, P)$ y $S(cf, P) = c\, s(f, P)$ si $c < 0$. Tomando supremo e ínfimo sobre las particiones (que invierte el orden cuando $c < 0$) se obtiene el resultado.
\end{proof}

\begin{prop}
  Sean $f, g\colon [a, b] \to \R$ acotadas con $f(x) \leq g(x)$ para todo $x \in [a, b]$. Entonces
  \[
    \lowint_a^b f \,\dd x \leq \lowint_a^b g \,\dd x \qquad \text{y} \qquad \upint_a^b f \,\dd x \leq \upint_a^b g \,\dd x.
  \]
\end{prop}

\begin{proof}
  Si $f \leq g$, entonces $m_i(f) \leq m_i(g)$ y $M_i(f) \leq M_i(g)$ para toda partición, y por lo tanto $s(f, P) \leq s(g, P)$ y $S(f, P) \leq S(g, P)$. Tomando supremo e ínfimo se obtienen las desigualdades para las integrales.
\end{proof}

\begin{definition}[Función integrable]
  Una función acotada $f\colon [a, b] \to \R$ es \emph{integrable} si $\lowint_a^b f(x) \,\dd x = \upint_a^b f(x) \,\dd x$. En ese caso, el valor común se denota $\int_a^b f(x) \,\dd x$.
\end{definition}

\begin{eg}
  La función $f\colon [a, b] \to \R$ dada por $f(x) = 1$ si $x \in \Q$ y $f(x) = 0$ si $x \notin \Q$ no es integrable.

  En efecto, como cada intervalo $[t_{i-1}, t_i]$ contiene racionales e irracionales, $m_i = 0$ y $M_i = 1$. Luego $s(f, P) = 0$ y $S(f, P) = b - a$ para toda partición $P$, y entonces $\lowint_a^b f \,\dd x = 0 \neq b - a = \upint_a^b f \,\dd x$.
\end{eg}

\begin{note}
  Sea $f\colon [a, b] \to \R$ acotada, y sean $\Sigma_{-}$ y $\Sigma_{+}$ los conjuntos de sumas inferiores y superiores de $f$. Que $f$ sea integrable significa que $\sup \Sigma_{-} = \inf \Sigma_{+}$. Esto ocurre si y sólo si para todo $\e > 0$ existen $s \in \Sigma_{-}$ y $S \in \Sigma_{+}$ con $S - s < \e$.
\end{note}

\begin{proof}
  ($\Rightarrow$) Si $I$ es el valor común, existen $s \in \Sigma_{-}$ y $S \in \Sigma_{+}$ con $s > I - \frac{\e}{2}$ y $S < I + \frac{\e}{2}$, de modo que $S - s < \e$.

  ($\Leftarrow$) Siempre $s \leq \sup \Sigma_{-} \leq \inf \Sigma_{+} \leq S$, de modo que $0 \leq \inf \Sigma_{+} - \sup \Sigma_{-} \leq S - s < \e$ para todo $\e > 0$, y entonces $\sup \Sigma_{-} = \inf \Sigma_{+}$.
\end{proof}

\begin{definition}[Oscilación]
  Dada una función acotada $f\colon [a, b] \to \R$ y un conjunto no vacío $X \subset [a, b]$, la \emph{oscilación} de $f$ en $X$ es $\omega(f, X) = \sup f(X) - \inf f(X)$.
\end{definition}

\begin{note}
  $\omega(f, X) = \sup \{ |f(x) - f(y)| : x, y \in X \}$.
\end{note}

\begin{proof}
  Para $x, y \in X$ vale $f(x) - f(y) \leq \sup f(X) - \inf f(X)$, y lo mismo intercambiando $x$ e $y$; luego $|f(x) - f(y)| \leq \omega(f, X)$. Recíprocamente, dado $\e > 0$ existen $x, y \in X$ con $f(x) > \sup f(X) - \frac{\e}{2}$ y $f(y) < \inf f(X) + \frac{\e}{2}$, y entonces $|f(x) - f(y)| > \omega(f, X) - \e$.
\end{proof}

\begin{theorem}
  Sea $f\colon [a, b] \to \R$ acotada. Son equivalentes:
  \begin{enumerate}[label=(\arabic*)]
    \item $f$ es integrable.
    \item Para todo $\e > 0$ existen particiones $P$ y $G$ de $[a, b]$ tales que $S(f, G) - s(f, P) < \e$.
    \item Para todo $\e > 0$ existe una partición $P$ de $[a, b]$ tal que $S(f, P) - s(f, P) < \e$.
    \item Para todo $\e > 0$ existe una partición $P = \{ t_0, \dots, t_n \}$ de $[a, b]$ tal que $\sum_{i=1}^{n} \omega_i (t_i - t_{i-1}) < \e$, donde $\omega_i = \omega(f, [t_{i-1}, t_i])$.
  \end{enumerate}
\end{theorem}

\begin{proof}
  (1) $\Leftrightarrow$ (2) Es la nota anterior, pues $s(f, P)$ recorre las sumas inferiores y $S(f, G)$ las superiores.

  (3) $\Leftrightarrow$ (4) Vale porque $\sum_{i=1}^{n} \omega_i (t_i - t_{i-1}) = S(f, P) - s(f, P)$.

  (3) $\Rightarrow$ (2) Es trivial, tomando $G = P$.

  (2) $\Rightarrow$ (3) Sean $P$ y $G$ con $S(f, G) - s(f, P) < \e$ y sea $P_0 = P \cup G$. Entonces
  \[
    s(f, P) \leq s(f, P_0) \leq S(f, P_0) \leq S(f, G),
  \]
  de donde $S(f, P_0) - s(f, P_0) \leq S(f, G) - s(f, P) < \e$.
\end{proof}

\section{Continuidad e integración}

\begin{theorem}
  Sea $f\colon [a, b] \to \R$ acotada y $a < c < b$. Entonces $f$ es integrable en $[a, b]$ si y sólo si $f|_{[a, c]}$ y $f|_{[c, b]}$ son integrables, y en ese caso
  \[
    \int_a^b f(x) \,\dd x = \int_a^c f(x) \,\dd x + \int_c^b f(x) \,\dd x.
  \]
\end{theorem}

\begin{proof}
  Por el teorema de aditividad de las integrales inferior y superior,
  \[
    \upint_a^b f \,\dd x - \lowint_a^b f \,\dd x = \left( \upint_a^c f \,\dd x - \lowint_a^c f \,\dd x \right) + \left( \upint_c^b f \,\dd x - \lowint_c^b f \,\dd x \right).
  \]
  Los dos términos entre paréntesis son no negativos, de modo que el miembro izquierdo es cero si y sólo si ambos lo son. Esto prueba la equivalencia, y la fórmula resulta de nuevo de la aditividad.
\end{proof}

\begin{theorem}
  Sean $f, g\colon [a, b] \to \R$ integrables y $c \in \R$.
  \begin{enumerate}[label=(\arabic*)]
    \item $c f$ es integrable y $\int_a^b c f(x) \,\dd x = c \int_a^b f(x) \,\dd x$.
    \item $f + g$ es integrable y $\int_a^b \bigl( f(x) + g(x) \bigr) \dd x = \int_a^b f(x) \,\dd x + \int_a^b g(x) \,\dd x$.
    \item Si $f(x) \leq g(x)$ para todo $x \in [a, b]$, entonces $\int_a^b f(x) \,\dd x \leq \int_a^b g(x) \,\dd x$.
    \item $|f|$ es integrable y $\left| \int_a^b f(x) \,\dd x \right| \leq \int_a^b |f(x)| \,\dd x$.
    \item $f g$ es integrable.
  \end{enumerate}
\end{theorem}

\begin{proof}
  (1), (2) y (3) se deducen de las proposiciones anteriores sobre integrales inferiores y superiores. Para (2), por ejemplo,
  \[
    \int_a^b f + \int_a^b g \leq \lowint_a^b (f + g) \leq \upint_a^b (f + g) \leq \int_a^b f + \int_a^b g,
  \]
  de modo que todas las desigualdades son igualdades. Para (1) se usa la proposición del producto por una constante, y para (3) la monotonía.

  (4) Para cualesquiera $x, y$ se tiene $\bigl| |f(x)| - |f(y)| \bigr| \leq |f(x) - f(y)|$. Por la nota sobre oscilaciones, $\omega(|f|, X) \leq \omega(f, X)$ para todo $X \subset [a, b]$ no vacío. En particular, para toda partición $P$, $S(|f|, P) - s(|f|, P) \leq S(f, P) - s(f, P)$, y por el teorema de caracterización $|f|$ es integrable. Como $-|f| \leq f \leq |f|$, por (1) y (3),
  \[
    -\int_a^b |f(x)| \,\dd x \leq \int_a^b f(x) \,\dd x \leq \int_a^b |f(x)| \,\dd x,
  \]
  es decir, $\left| \int_a^b f \,\dd x \right| \leq \int_a^b |f| \,\dd x$.

  (5) Sea $k > 0$ con $|f(x)| \leq k$ y $|g(x)| \leq k$ para todo $x \in [a, b]$. Si $x, y \in [t_{i-1}, t_i]$,
  \begin{align*}
    |f(x) g(x) - f(y) g(y)| & = |f(x) \bigl( g(x) - g(y) \bigr) + g(y) \bigl( f(x) - f(y) \bigr)|          \\
                            & \leq k\, \omega_i(g) + k\, \omega_i(f).
  \end{align*}
  Luego $\omega_i(fg) \leq k \bigl( \omega_i(f) + \omega_i(g) \bigr)$. Dado $\e > 0$, existen particiones $P_1$ y $P_2$ con $S(f, P_1) - s(f, P_1) < \frac{\e}{2k}$ y $S(g, P_2) - s(g, P_2) < \frac{\e}{2k}$. Para $P = P_1 \cup P_2$, que refina a ambas, las mismas desigualdades valen, y entonces
  \[
    S(fg, P) - s(fg, P) \leq k \bigl( S(f, P) - s(f, P) \bigr) + k \bigl( S(g, P) - s(g, P) \bigr) < \e. \qedhere
  \]
\end{proof}

\begin{theorem}
  Toda función continua $f\colon [a, b] \to \R$ es integrable.
\end{theorem}

\begin{proof}
  Como $[a, b]$ es compacto, $f$ es acotada y uniformemente continua. Dado $\e > 0$ existe $\delta > 0$ tal que $|x - y| < \delta \Rightarrow |f(x) - f(y)| < \frac{\e}{2(b - a)}$. Sea $n \in \N$ con $\frac{b - a}{n} < \delta$ y sea $P$ la partición $t_i = a + i \frac{b - a}{n}$, con $i \in \{0, \dots, n\}$. Si $x, y \in [t_{i-1}, t_i]$, entonces $|x - y| < \delta$ y $|f(x) - f(y)| < \frac{\e}{2(b - a)}$, de donde $\omega_i \leq \frac{\e}{2(b - a)}$. Luego
  \[
    \sum_{i=1}^{n} \omega_i (t_i - t_{i-1}) \leq \frac{\e}{2(b - a)} \sum_{i=1}^{n} (t_i - t_{i-1}) = \frac{\e}{2} < \e,
  \]
  y $f$ es integrable.
\end{proof}

\begin{theorem}
  Sea $f\colon [a, b] \to \R$ acotada. Si $f|_{[a, c]}$ es integrable para todo $c \in (a, b)$, entonces $f$ es integrable. Análogamente, si $f|_{[c, b]}$ es integrable para todo $c \in (a, b)$, entonces $f$ es integrable.
\end{theorem}

\begin{proof}
  Probamos la primera afirmación; la segunda es simétrica. Sea $k > 0$ con $|f(x)| \leq k$ para todo $x \in [a, b]$, y sea $\e > 0$. Elegimos $c \in (a, b)$ con $k (b - c) < \frac{\e}{4}$. Como $f|_{[a, c]}$ es integrable, existe una partición $\{ t_0, \dots, t_n \}$ de $[a, c]$ tal que $\sum_{i=1}^{n} \omega_i (t_i - t_{i-1}) < \frac{\e}{2}$. Agregando $t_{n+1} = b$ se obtiene una partición de $[a, b]$ y $\omega_{n+1} \leq 2k$, de modo que $\omega_{n+1} (t_{n+1} - t_n) \leq 2k (b - c) < \frac{\e}{2}$. Luego $\sum_{i=1}^{n+1} \omega_i (t_i - t_{i-1}) < \e$.
\end{proof}

\begin{corollary}
  Sea $f\colon [a, b] \to \R$ acotada. Si $f|_{[c, d]}$ es integrable para todos $c, d$ con $a < c < d < b$, entonces $f$ es integrable.
\end{corollary}

\begin{proof}
  Sea $m \in (a, b)$. Para cada $c \in (a, m)$, $f|_{[c, m]}$ es integrable, de modo que $f|_{[a, m]}$ es integrable por la segunda parte del teorema anterior. Análogamente, $f|_{[m, d]}$ es integrable para todo $d \in (m, b)$, y por la primera parte $f|_{[m, b]}$ es integrable. Por aditividad, $f$ es integrable en $[a, b]$.
\end{proof}

\begin{corollary}
  Toda función acotada $f\colon [a, b] \to \R$ con finitas discontinuidades es integrable.
\end{corollary}

\begin{proof}
  Sean $a = s_0 < s_1 < \cdots < s_r = b$ tales que todos los puntos de discontinuidad de $f$ pertenecen a $\{ s_0, \dots, s_r \}$. En cada $[s_{j-1}, s_j]$, la restricción $f|_{[c, d]}$ es continua, y por lo tanto integrable, para $s_{j-1} < c < d < s_j$. Por el corolario anterior $f|_{[s_{j-1}, s_j]}$ es integrable, y aplicando $r - 1$ veces la aditividad, $f$ es integrable en $[a, b]$.
\end{proof}

\begin{eg}
  Sea $f\colon [a, b] \to \R$ dada por $f(x) = 0$ si $x$ es irracional o $x = 0$, y $f(x) = \frac{1}{q}$ si $x = \frac{p}{q}$ con $p \in \Z$, $q \in \N$ y $\mathrm{mcd}(p, q) = 1$. Entonces $f$ es discontinua en todos los racionales no nulos de $[a, b]$ (un conjunto infinito), pero es integrable y $\int_a^b f(x) \,\dd x = 0$.

  Dado $\e > 0$, sea $\eta = \frac{\e}{b - a}$ y $F = \{ x \in [a, b] : f(x) \geq \eta \}$. El conjunto $F$ es finito: si $x = \frac{p}{q} \in F$ entonces $q \leq \frac{1}{\eta}$, y para cada $q$ hay finitos $p$ con $\frac{p}{q} \in [a, b]$. Sea $N = |F|$ y tomemos la partición uniforme de $[a, b]$ en $n$ intervalos con $\frac{2N (b - a)}{n} < \e$. Cada punto de $F$ pertenece a lo sumo a dos intervalos, de modo que los intervalos que intersecan a $F$ tienen longitud total menor que $\e$, y en ellos $M_i \leq 1$. En los demás intervalos todos los valores de $f$ son menores que $\eta$, luego $M_i \leq \eta$. Entonces
  \[
    S(f, P) = \sum M_i (t_i - t_{i-1}) < \e + \eta (b - a) = 2\e.
  \]
  Por lo tanto $\upint_a^b f \,\dd x = 0$. Como $f \geq 0$, $0 \leq \lowint_a^b f \,\dd x \leq \upint_a^b f \,\dd x = 0$, y entonces $f$ es integrable con integral $0$.
\end{eg}


\clase{05/11/2024}{Teorema fundamental del cálculo e integración múltiple}
\section{Teorema fundamental del cálculo}

Para \(f\) integrable en \([a, b]\) y \(c, d \in [a, b]\) con \(c > d\) convenimos \(\int_c^d f(x) \,\dd x = -\int_d^c f(x) \,\dd x\), y \(\int_c^c f(x) \,\dd x = 0\). Con esta convención vale \(\int_a^b f = \int_a^c f + \int_c^b f\) para \(a, b, c\) cualesquiera en el dominio.

\begin{theorem}[Teorema fundamental del cálculo]
  Sea \(f\colon [a, b] \to \R\) integrable. Si \(f\) es continua en \(c \in [a, b]\), entonces \(F\colon [a, b] \to \R\), \(F(x) = \int_a^x f(t) \,\dd t\), es derivable en \(c\) (con derivadas laterales en los extremos) y \(F'(c) = f(c)\).
\end{theorem}

\begin{proof}
  Sea \(\e > 0\). Como \(f\) es continua en \(c\), existe \(\delta > 0\) tal que \(t \in [a, b]\) y \(|t - c| < \delta\) implican \(|f(t) - f(c)| < \e\). Sea \(h \neq 0\) con \(|h| < \delta\) y \(c + h \in [a, b]\). Como \(F(c + h) - F(c) = \int_c^{c+h} f(t) \,\dd t\) y \(\int_c^{c+h} f(c) \,\dd t = h f(c)\),
  \[
    \left| \frac{F(c + h) - F(c)}{h} - f(c) \right| = \frac{1}{|h|} \left| \int_c^{c+h} \bigl( f(t) - f(c) \bigr) \dd t \right| \leq \frac{1}{|h|} \cdot \e |h| = \e.
  \]
  Luego \(\lim_{h \to 0} \frac{F(c + h) - F(c)}{h} = f(c)\).
\end{proof}

\begin{corollary}[Regla de Barrow]
  Sea \(f\colon [a, b] \to \R\) integrable y sea \(F\colon [a, b] \to \R\) derivable con \(F' = f\). Entonces \(\int_a^b f(x) \,\dd x = F(b) - F(a)\).
\end{corollary}

\begin{proof}
  Sea \(P = \{ t_0, \dots, t_n \}\) una partición de \([a, b]\). Por el teorema del valor medio, para cada \(i\) existe \(\alpha_i \in (t_{i-1}, t_i)\) con \(F(t_i) - F(t_{i-1}) = f(\alpha_i)(t_i - t_{i-1})\). Entonces
  \[
    F(b) - F(a) = \sum_{i=1}^{n} \bigl( F(t_i) - F(t_{i-1}) \bigr) = \sum_{i=1}^{n} f(\alpha_i)(t_i - t_{i-1}).
  \]
  Como \(m_i \leq f(\alpha_i) \leq M_i\), resulta \(s(f, P) \leq F(b) - F(a) \leq S(f, P)\) para toda partición \(P\). Luego \(\lowint_a^b f \,\dd x \leq F(b) - F(a) \leq \upint_a^b f \,\dd x\), y como \(f\) es integrable, \(F(b) - F(a) = \int_a^b f(x) \,\dd x\).
\end{proof}

\begin{definition}[Norma de una partición]
  Sea \(P = \{ t_0, \dots, t_n \}\) una partición de \([a, b]\). La \emph{norma} de \(P\) es \(\norm{P} = \max_{1 \leq i \leq n} (t_i - t_{i-1})\).
\end{definition}

\begin{theorem}
  Sea \(f\colon [a, b] \to \R\) acotada. Para todo \(\e > 0\) existe \(\delta > 0\) tal que si \(\norm{P} < \delta\), entonces \(S(f, P) < \upint_a^b f(x) \,\dd x + \e\).
\end{theorem}

\begin{proof}
  Supongamos primero \(f \geq 0\). Sea \(M = \sup f\). Dado \(\e > 0\) existe una partición \(P_0 = \{ t_0, \dots, t_n \}\) con \(S(f, P_0) < \upint_a^b f \,\dd x + \frac{\e}{2}\). Sea \(\delta = \frac{\e}{2n (M + 1)}\) y sea \(P\) una partición con \(\norm{P} < \delta\).

  Dividimos los intervalos de \(P\) en dos tipos: los contenidos en algún \([t_{i-1}, t_i]\) y los restantes. Cada intervalo del segundo tipo contiene algún \(t_i\) en su interior, así que hay a lo sumo \(n\) de ellos, y cada uno tiene longitud menor que \(\delta\) y \(M_k \leq M\). Para cada \(i\), los intervalos del primer tipo contenidos en \([t_{i-1}, t_i]\) tienen supremo \(\leq M_i\) y longitudes que suman a lo sumo \(t_i - t_{i-1}\), de modo que su aporte a \(S(f, P)\) es a lo sumo \(M_i (t_i - t_{i-1})\) (pues \(M_i \geq 0\)). Entonces
  \[
    S(f, P) \leq \sum_{i=1}^{n} M_i (t_i - t_{i-1}) + n M \delta < S(f, P_0) + \frac{\e}{2} < \upint_a^b f \,\dd x + \e.
  \]

  Para \(f\) acotada cualquiera, existe \(c \in \R\) con \(g = f + c \geq 0\). Como \(S(g, P) = S(f, P) + c(b - a)\) para toda partición y \(\upint_a^b g \,\dd x = \upint_a^b f \,\dd x + c(b - a)\), el resultado para \(g\) implica el de \(f\).
\end{proof}

\begin{corollary}
  Si \(f\colon [a, b] \to \R\) es acotada, entonces
  \[
    \lim_{\norm{P} \to 0} S(f, P) = \upint_a^b f(x) \,\dd x, \qquad \lim_{\norm{P} \to 0} s(f, P) = \lowint_a^b f(x) \,\dd x.
  \]
  En particular, si \(f\) es integrable, \(\int_a^b f(x) \,\dd x = \lim_{\norm{P} \to 0} S(f, P) = \lim_{\norm{P} \to 0} s(f, P)\).
\end{corollary}

\begin{proof}
  La primera igualdad es el teorema anterior, junto con \(S(f, P) \geq \upint_a^b f \,\dd x\) para toda \(P\). Para la segunda, \(s(f, P) = -S(-f, P)\) y \(\lowint_a^b f \,\dd x = -\upint_a^b (-f) \,\dd x\).
\end{proof}

\begin{definition}[Partición punteada]
  Una \emph{partición punteada} \(P^{*}\) de \([a, b]\) es una partición \(P = \{ t_0, \dots, t_n \}\) junto con una elección de puntos \(\alpha_i \in [t_{i-1}, t_i]\) para \(i \in \{1, \dots, n\}\). Su norma es \(\norm{P^{*}} = \norm{P}\).
\end{definition}

\begin{definition}[Suma de Riemann]
  Si \(f\colon [a, b] \to \R\) es acotada y \(P^{*}\) es una partición punteada, la \emph{suma de Riemann} de \(f\) asociada a \(P^{*}\) es
  \[
    R(f, P^{*}) = \sum_{i=1}^{n} f(\alpha_i)(t_i - t_{i-1}).
  \]
  Decimos que \(I = \lim_{\norm{P} \to 0} R(f, P^{*})\) si para todo \(\e > 0\) existe \(\delta > 0\) tal que \(|R(f, P^{*}) - I| < \e\) para toda partición punteada \(P^{*}\) con \(\norm{P^{*}} < \delta\).
\end{definition}

\begin{theorem}
  Sea \(f\colon [a, b] \to \R\) acotada. Existe \(I = \lim_{\norm{P} \to 0} R(f, P^{*})\) si y sólo si \(f\) es integrable, y en ese caso \(I = \int_a^b f(x) \,\dd x\).
\end{theorem}

\begin{proof}
  (\(\Leftarrow\)) Si \(f\) es integrable, por el corolario anterior \(\lim_{\norm{P} \to 0} s(f, P) = \lim_{\norm{P} \to 0} S(f, P) = \int_a^b f \,\dd x\). Como \(s(f, P) \leq R(f, P^{*}) \leq S(f, P)\) para toda partición punteada, \(\lim_{\norm{P} \to 0} R(f, P^{*}) = \int_a^b f \,\dd x\).

  (\(\Rightarrow\)) Supongamos que el límite \(I\) existe y sea \(\e > 0\). Existe \(\delta > 0\) tal que \(|R(f, P^{*}) - I| < \frac{\e}{4}\) para toda partición punteada con \(\norm{P^{*}} < \delta\). Fijemos una partición \(P = \{ t_0, \dots, t_n \}\) con \(\norm{P} < \delta\). Elegimos \(\alpha_i \in [t_{i-1}, t_i]\) con \(f(\alpha_i) < m_i + \frac{\e}{4n(t_i - t_{i-1})}\); entonces
  \[
    R(f, P^{*}) < \sum_{i=1}^{n} m_i (t_i - t_{i-1}) + \frac{\e}{4} = s(f, P) + \frac{\e}{4}.
  \]
  Análogamente, eligiendo otros puntos \(\tilde{\alpha}_i\) con \(f(\tilde{\alpha}_i) > M_i - \frac{\e}{4n(t_i - t_{i-1})}\), obtenemos una partición punteada \(\tilde{P}^{*}\) con \(R(f, \tilde{P}^{*}) > S(f, P) - \frac{\e}{4}\). Luego
  \[
    I - \frac{\e}{2} < R(f, P^{*}) - \frac{\e}{4} < s(f, P) \leq S(f, P) < R(f, \tilde{P}^{*}) + \frac{\e}{4} < I + \frac{\e}{2},
  \]
  de donde \(S(f, P) - s(f, P) < \e\). Por el teorema de caracterización, \(f\) es integrable. Además \(s(f, P) \leq \int_a^b f \,\dd x \leq S(f, P)\) y \(I\) pertenece al mismo intervalo \([s(f, P), S(f, P)]\), de longitud menor que \(\e\). Como \(\e\) es arbitrario, \(I = \int_a^b f(x) \,\dd x\).
\end{proof}

\begin{eg}
  Sea \(f\colon [1, 2] \to \R\), \(f(x) = \frac{1}{x}\). Por la regla de Barrow, \(\int_1^2 \frac{1}{x} \,\dd x = \ln 2 - \ln 1 = \ln 2\). Para cada \(n \in \N\) sea \(P_n\) la partición \(t_i = \frac{n + i}{n}\), con \(i \in \{0, \dots, n\}\), y elijamos \(\alpha_i = t_i\) en cada \([t_{i-1}, t_i]\). Entonces \(t_i - t_{i-1} = \frac{1}{n}\), \(f(\alpha_i) = \frac{n}{n + i}\) y
  \[
    R(f, P_n^{*}) = \sum_{i=1}^{n} \frac{n}{n + i} \cdot \frac{1}{n} = \sum_{i=1}^{n} \frac{1}{n + i}.
  \]
  Como \(\norm{P_n} = \frac{1}{n} \to 0\) y \(f\) es integrable,
  \[
    \lim_{n \to \infty} \left( \frac{1}{n + 1} + \cdots + \frac{1}{2n} \right) = \ln 2.
  \]
\end{eg}

\section{Integral de Riemann en varias dimensiones}

Un \emph{intervalo \(n\)-dimensional} es un conjunto \(I = I_1 \times \cdots \times I_n \subset \R^n\), donde \(I_1, \dots, I_n\) son intervalos de \(\R\). Se admite el caso degenerado en que algunos \(I_k\) se reducen a un punto. Si todos los \(I_k\) son abiertos (respectivamente cerrados, acotados), entonces \(I\) es abierto (respectivamente cerrado, acotado) en \(\R^n\). Notamos \(|I| = |I_1| \cdots |I_n|\), donde \(|I_k|\) es la longitud de \(I_k\).

Sea \(I = [a_1, b_1] \times \cdots \times [a_n, b_n]\) un intervalo compacto. Si \(P_k\) es una partición de \([a_k, b_k]\) para cada \(k\), entonces \(P = P_1 \times \cdots \times P_n\) es una \emph{partición} de \(I\), que determina subintervalos compactos \(Q_1, \dots, Q_N\) (los productos de subintervalos de las \(P_k\)), cuyos interiores son disjuntos dos a dos y cuya unión es \(I\). Decimos que una partición \(P' = P'_1 \times \cdots \times P'_n\) es \emph{más fina} que \(P\) si \(P_k \subset P'_k\) para todo \(k\).

\begin{definition}[Suma de Riemann e integrabilidad]
  Sea \(f\colon I \to \R\) acotada, \(P\) una partición de \(I\) con subintervalos \(Q_1, \dots, Q_N\) y \(\alpha_j \in Q_j\) para cada \(j\). La \emph{suma de Riemann} asociada es \(R(f, P^{*}) = \sum_{j=1}^{N} f(\alpha_j) |Q_j|\). Decimos que \(f\) es \emph{integrable Riemann} en \(I\) si existe \(A \in \R\) tal que para todo \(\e > 0\) existe una partición \(P_{\e}\) de \(I\) con \(|R(f, P^{*}) - A| < \e\) para toda partición \(P\) más fina que \(P_{\e}\) y toda elección de puntos \(\alpha_j \in Q_j\). En ese caso escribimos
  \[
    A = \int_I f(x) \,\dd x = \int \cdots \int_I f(x_1, \dots, x_n) \,\dd x_1 \cdots \dd x_n.
  \]
\end{definition}

\begin{definition}[Sumas e integrales de Darboux]
  Sea \(I \subset \R^n\) un intervalo compacto, \(f\colon I \to \R\) acotada y \(P\) una partición de \(I\) con subintervalos \(Q_1, \dots, Q_N\). Sean \(m_j = \inf f(Q_j)\) y \(M_j = \sup f(Q_j)\). Las sumas inferior y superior son \(s(f, P) = \sum_{j=1}^{N} m_j |Q_j|\) y \(S(f, P) = \sum_{j=1}^{N} M_j |Q_j|\), y las integrales inferior y superior son
  \[
    \lowint_I f(x) \,\dd x = \sup_P s(f, P), \qquad \upint_I f(x) \,\dd x = \inf_P S(f, P).
  \]
  Estas cantidades mantienen las propiedades vistas en \(\R\).
\end{definition}

\begin{theorem}
  Sea \(I \subset \R^n\) un intervalo compacto y \(f\colon I \to \R\) acotada. Son equivalentes:
  \begin{enumerate}[label=(\arabic*)]
    \item \(f\) es integrable en \(I\).
    \item Para todo \(\e > 0\) existe una partición \(P_{\e}\) de \(I\) tal que \(S(f, P) - s(f, P) < \e\) para toda \(P\) más fina que \(P_{\e}\).
    \item \(\lowint_I f(x) \,\dd x = \upint_I f(x) \,\dd x\).
  \end{enumerate}
\end{theorem}

\begin{proof}
  La demostración es análoga a la del caso de una variable y se deja como ejercicio.
\end{proof}

\begin{prop}
  Si \(I \subset \R^n\) es un intervalo compacto y \(f\colon I \to \R\) es continua, entonces \(f\) es integrable en \(I\).
\end{prop}

\begin{proof}
  Podemos suponer \(|I| > 0\). Como \(I\) es compacto, \(f\) es acotada y uniformemente continua. Dado \(\e > 0\) existe \(\delta > 0\) tal que \(\norm{x - y} < \delta\) implica \(|f(x) - f(y)| < \frac{\e}{2|I|}\). Sea \(P\) una partición cuyos subintervalos tienen diámetro menor que \(\delta\) (basta subdividir cada lado en partes pequeñas). Entonces \(M_j - m_j \leq \frac{\e}{2|I|}\) para todo \(j\) y
  \[
    S(f, P) - s(f, P) = \sum_{j} (M_j - m_j) |Q_j| \leq \frac{\e}{2|I|} \sum_j |Q_j| = \frac{\e}{2} < \e.
  \]
  Como \(s(f, P) \leq \lowint_I f \leq \upint_I f \leq S(f, P)\), resulta \(0 \leq \upint_I f - \lowint_I f < \e\) para todo \(\e > 0\), es decir, \(f\) es integrable.
\end{proof}

\begin{theorem}
  Sea \(Q = [a, b] \times [c, d] \subset \R^2\) y \(f\colon Q \to \R\) acotada. Entonces
  \[
    \lowint_Q f \leq \lowint_a^b \left( \upint_c^d f(x, y) \,\dd y \right) \dd x \leq \upint_a^b \left( \upint_c^d f(x, y) \,\dd y \right) \dd x \leq \upint_Q f.
  \]
\end{theorem}

\begin{proof}
  Sea \(F(x) = \upint_c^d f(x, y) \,\dd y\) para \(x \in [a, b]\), que está acotada por \(M(d - c)\) si \(|f| \leq M\). Sean \(P_1 = \{ x_0, \dots, x_p \}\) y \(P_2 = \{ y_0, \dots, y_q \}\) particiones de \([a, b]\) y \([c, d]\), y \(P = P_1 \times P_2\), con subrectángulos \(Q_{ij} = [x_{i-1}, x_i] \times [y_{j-1}, y_j]\). Sean \(m_{ij} = \inf f(Q_{ij})\) y \(M_{ij} = \sup f(Q_{ij})\).

  Por la aditividad de la integral superior, \(F(x) = \sum_{j=1}^{q} \upint_{y_{j-1}}^{y_j} f(x, y) \,\dd y\). Si \(x \in [x_{i-1}, x_i]\), entonces \(m_{ij} \leq f(x, y) \leq M_{ij}\) para \(y \in [y_{j-1}, y_j]\), y por lo tanto
  \[
    \sum_{j=1}^{q} m_{ij} (y_j - y_{j-1}) \leq F(x) \leq \sum_{j=1}^{q} M_{ij} (y_j - y_{j-1}).
  \]
  Tomando ínfimo y supremo de \(F\) en \([x_{i-1}, x_i]\) y multiplicando por \(x_i - x_{i-1}\), y sumando en \(i\), obtenemos
  \[
    s(f, P) \leq s(F, P_1) \qquad \text{y} \qquad S(F, P_1) \leq S(f, P).
  \]
  Entonces \(s(f, P) \leq s(F, P_1) \leq \lowint_a^b F \,\dd x \leq \upint_a^b F \,\dd x \leq S(F, P_1) \leq S(f, P)\). Como esto vale para toda partición \(P = P_1 \times P_2\) de \(Q\), tomando supremo sobre \(P\) en el miembro izquierdo e ínfimo en el derecho se obtiene el resultado.
\end{proof}

\begin{note}
  Con la misma demostración se obtiene lo siguiente.
  \begin{enumerate}[label=(\arabic*)]
    \item Vale el teorema reemplazando \(\upint_c^d\) por \(\lowint_c^d\):
          \[
            \lowint_Q f \leq \lowint_a^b \bigl( \lowint_c^d f \,\dd y \bigr) \dd x \leq \upint_a^b \bigl( \lowint_c^d f \,\dd y \bigr) \dd x \leq \upint_Q f.
          \]
    \item Vale intercambiando los papeles de \(x\) e \(y\):
          \[
            \lowint_Q f \leq \lowint_c^d \bigl( \upint_a^b f \,\dd x \bigr) \dd y \leq \upint_c^d \bigl( \upint_a^b f \,\dd x \bigr) \dd y \leq \upint_Q f,
          \]
          y también con \(\lowint_a^b\) en el interior.
    \item Si \(f\) es integrable en \(Q\), los extremos de todas estas cadenas de desigualdades coinciden con \(\iint_Q f\), y por lo tanto todas las integrales iteradas (con cualquier combinación de integrales inferiores o superiores) son iguales a \(\iint_Q f\).
  \end{enumerate}
\end{note}

\section{Teorema de Fubini}

\begin{corollary}[Teorema de Fubini]
  Si \(f\colon Q \to \R\) es continua en \(Q = [a, b] \times [c, d]\), entonces
  \[
    \iint_Q f = \int_a^b \left( \int_c^d f(x, y) \,\dd y \right) \dd x = \int_c^d \left( \int_a^b f(x, y) \,\dd x \right) \dd y.
  \]
\end{corollary}

\begin{proof}
  Por la proposición anterior \(f\) es integrable en \(Q\). Para cada \(x\), la función \(y \mapsto f(x, y)\) es continua, y por lo tanto integrable, de modo que \(\upint_c^d f(x, y) \,\dd y = \int_c^d f(x, y) \,\dd y\). Por la nota anterior (3), \(\lowint_a^b F \,\dd x = \upint_a^b F \,\dd x = \iint_Q f\), donde \(F(x) = \int_c^d f(x, y) \,\dd y\). Es decir, \(F\) es integrable y \(\iint_Q f = \int_a^b F(x) \,\dd x\). La otra igualdad se prueba intercambiando los papeles de \(x\) e \(y\).
\end{proof}

\section{Contenido de Jordan}

\begin{definition}[Contenidos interior y exterior]
  Sea \(S\) un subconjunto de un intervalo compacto \(I \subset \R^n\). Para cada partición \(P\) de \(I\) sea \(\underline{J}(S, P)\) la suma de las medidas de los subintervalos de \(P\) contenidos en el interior de \(S\), y \(\overline{J}(S, P)\) la suma de las medidas de los subintervalos de \(P\) que intersecan a \(\cl{S}\). El \emph{contenido interior} y el \emph{contenido exterior} de \(S\) son
  \[
    \underline{c}(S) = \sup_P \underline{J}(S, P), \qquad \overline{c}(S) = \inf_P \overline{J}(S, P).
  \]
\end{definition}

\begin{definition}[Contenido de Jordan]
  Si \(\underline{c}(S) = \overline{c}(S)\), este valor común se denota \(c(S)\) y se llama \emph{contenido de Jordan} de \(S\). En ese caso decimos que \(S\) es \emph{medible Jordan}.
\end{definition}

\begin{note}
  \begin{enumerate}[label=(\arabic*)]
    \item Se puede ver que \(\underline{c}(S)\) y \(\overline{c}(S)\) dependen únicamente de \(S\) y no del intervalo \(I\) elegido.
    \item \(0 \leq \underline{c}(S) \leq \overline{c}(S)\).
    \item Si \(c(S) = 0\), entonces \(\overline{c}(S) = 0\) y, por definición de ínfimo, para todo \(\e > 0\) el conjunto \(S\) puede cubrirse con finitos intervalos cuyas medidas suman menos que \(\e\).
  \end{enumerate}
\end{note}


\clase{08/11/2024}{Contenido y medida de Jordan}
\section{Ejemplos de contenido de Jordan}

\begin{eg}
  El conjunto $X = \Q \cap [a, b]$, con $a < b$, no tiene contenido nulo.
\end{eg}

\begin{proof}
  Como $\cl{X} = [a, b]$, para toda partición $P$ de un intervalo $I \supseteq [a, b]$ los subintervalos de $P$ que intersecan a $\cl{X}$ cubren $[a, b]$, y sus medidas suman al menos $b - a$. Luego $\overline{c}(X) \geq b - a > 0$.
\end{proof}

\begin{lemma}
  Sea $X \subseteq \R$ acotado. Si para todo $\e > 0$ existen un conjunto finito $F \subseteq X$ e intervalos abiertos $I_1, \dots, I_k$ tales que $X \setminus F \subseteq I_1 \cup \cdots \cup I_k$ y $|I_1| + \cdots + |I_k| < \e$, entonces $c(X) = 0$.
\end{lemma}

\begin{proof}
  Dado $\e > 0$, tomamos $F, I_1, \dots, I_k$ como en el enunciado con $\sum_{i=1}^k |I_i| < \frac{\e}{2}$. Si $F = \{x_1, \dots, x_q\}$, podemos cubrir $F$ con intervalos abiertos $I_{k+1}, \dots, I_n$ de longitud menor que $\frac{\e}{2q}$ cada uno. Entonces $X \subseteq I_1 \cup \cdots \cup I_n$ y $\sum_{i=1}^n |I_i| < \e$.

  Sean $\bar{I}_i = [\alpha_i, \beta_i]$ y $[a, b]$ un intervalo que contiene a $X$ y a todos los $\bar{I}_i$. Para $\eta > 0$, sea $P$ la partición de $[a, b]$ cuyos puntos son $a$, $b$ y todos los $\alpha_i$, $\beta_i$, $\alpha_i - \eta$, $\beta_i + \eta$ que están en $[a, b]$. Un subintervalo $R$ de $P$ que interseca a $\cl{X} \subseteq \bigcup_i \bar{I}_i$ interseca algún $\bar{I}_i$. Como $\alpha_i$ y $\beta_i$ son puntos de $P$, o bien $R \subseteq \bar{I}_i$, o bien $R$ comparte con $\bar{I}_i$ un único extremo, y en ese caso $R = [\beta_i, \beta_i']$ o $R = [\alpha_i', \alpha_i]$ con longitud a lo sumo $\eta$. Hay a lo sumo $2n$ subintervalos del segundo tipo, y los del primer tipo suman a lo sumo $\sum_i |I_i|$, pues sus interiores son disjuntos. Luego
  \[
    \overline{c}(X) \leq \overline{J}(X, P) \leq \sum_{i=1}^n |I_i| + 2n\eta < \e + 2n\eta.
  \]
  Como $\eta$ es arbitrario, $\overline{c}(X) \leq \e$, y como $\e$ es arbitrario, $\overline{c}(X) = 0$. Entonces $c(X) = 0$.
\end{proof}

\begin{eg}
  El conjunto de Cantor $C$ tiene contenido nulo.
\end{eg}

\begin{proof}
  Sea $F_n$ el conjunto obtenido en el paso $n$ de la construcción de $C$, unión de $2^n$ intervalos cerrados disjuntos de longitud $3^{-n}$. Entonces $C \subseteq F_n$ y la suma de las longitudes de los intervalos de $F_n$ es $1 - \frac{1}{3} \sum_{i=0}^{n-1} {\left( \frac{2}{3} \right)}^i = {\left( \frac{2}{3} \right)}^n$. Cada intervalo cerrado de $F_n$ está contenido en un intervalo abierto de longitud $2 \cdot 3^{-n}$, y la suma de las longitudes de éstos es $2 {\left( \frac{2}{3} \right)}^n$, que es menor que $\e$ para $n$ suficientemente grande. Por el lema (con $F = \varnothing$), $c(C) = 0$.
\end{proof}

\begin{theorem}
  Sea $S \subset \R^n$ acotado. Entonces $\overline{c}(\partial S) = \overline{c}(S) - \underline{c}(S)$. En particular, $S$ es medible Jordan si y sólo si $\partial S$ tiene contenido de Jordan nulo.
\end{theorem}

\begin{proof}
  Sea $I \subset \R^n$ un intervalo compacto que contiene a $\cl{S}$. Veamos primero que para toda partición $P$ de $I$,
  \[
    \overline{J}(\partial S, P) = \overline{J}(S, P) - \underline{J}(S, P).
  \]
  Un subintervalo $R$ de $P$ está contenido en $\interior S$ si y sólo si interseca a $\cl{S}$ y no interseca a $\partial S$. En efecto, si $R \subseteq \interior S$ entonces $R \cap \partial S = \varnothing$ y $R \cap \cl{S} \neq \varnothing$. Recíprocamente, si $R \cap \partial S = \varnothing$, entonces $R \subseteq \interior S \cup (\R^n \setminus \cl{S})$, unión de dos abiertos disjuntos; por ser $R$ conexo está contenido en uno de ellos, y si interseca a $\cl{S}$, está contenido en $\interior S$. Luego los subintervalos que intersecan a $\partial S = \cl{\partial S}$ son exactamente los que intersecan a $\cl{S}$ y no están contenidos en $\interior S$, lo que da la igualdad.

  Como $\overline{J}(S, P) \geq \overline{c}(S)$ y $\underline{J}(S, P) \leq \underline{c}(S)$, resulta $\overline{J}(\partial S, P) \geq \overline{c}(S) - \underline{c}(S)$ para toda $P$, y entonces $\overline{c}(\partial S) \geq \overline{c}(S) - \underline{c}(S)$.

  Para la desigualdad opuesta, dado $\e > 0$ sea $P_1$ tal que $\overline{J}(S, P_1) < \overline{c}(S) + \frac{\e}{2}$ y $P_2$ tal que $\underline{J}(S, P_2) > \underline{c}(S) - \frac{\e}{2}$. Sea $P = P_1 \cup P_2$ (refinamiento común). Como refinar achica las sumas exteriores y agranda las interiores,
  \begin{align*}
    \overline{c}(\partial S) \leq \overline{J}(\partial S, P) & = \overline{J}(S, P) - \underline{J}(S, P) \leq \overline{J}(S, P_1) - \underline{J}(S, P_2) \\
                                                              & < \overline{c}(S) - \underline{c}(S) + \e.
  \end{align*}
  Como $\e > 0$ es arbitrario, $\overline{c}(\partial S) \leq \overline{c}(S) - \underline{c}(S)$.

  Finalmente, $S$ es medible Jordan si y sólo si $\overline{c}(S) = \underline{c}(S)$, es decir, si y sólo si $\overline{c}(\partial S) = 0$, que equivale a $c(\partial S) = 0$ pues $0 \leq \underline{c}(\partial S) \leq \overline{c}(\partial S)$.
\end{proof}

\section{Medida de Lebesgue nula}

\begin{definition}[Medida de Lebesgue nula]
  Un conjunto $S \subseteq \R^n$ tiene \emph{medida de Lebesgue nula} si para todo $\e > 0$ puede cubrirse con una cantidad numerable de intervalos cuyas medidas suman menos que $\e$. Notamos $|S| = 0$.
\end{definition}

\begin{eg}
  Si $c(X) = 0$, entonces $|X| = 0$, pues existe una partición $P$ tal que los subintervalos que intersecan a $\cl{X}$ cubren a $X$ y suman menos que $\e$. Por ejemplo, el conjunto de Cantor tiene medida nula.
\end{eg}

\begin{prop}
  Toda unión numerable de conjuntos de medida nula tiene medida nula. En particular, todo conjunto numerable tiene medida nula; por ejemplo, $|\Q| = 0$.
\end{prop}

\begin{proof}
  Sean $F_1, F_2, \dots$ conjuntos de medida nula y $\e > 0$. Para cada $k \in \N$ cubrimos $F_k$ con una cantidad numerable de intervalos cuyas medidas suman menos que $\frac{\e}{2^k}$. La unión de todos estos intervalos es una familia numerable que cubre $\bigcup_k F_k$, y la suma de sus medidas es menor que $\sum_k \frac{\e}{2^k} = \e$.

  Todo conjunto con un único punto tiene medida nula (se cubre con un intervalo de medida menor que $\e$), de modo que todo conjunto numerable, por ejemplo $\Q = \{r_1, r_2, \dots\}$, tiene medida nula. Explícitamente, cubriendo a $r_k$ con el intervalo $\left( r_k - \frac{\e}{2^{k+2}}, r_k + \frac{\e}{2^{k+2}} \right)$ se obtiene una suma total $\sum_k \frac{\e}{2^{k+1}} = \frac{\e}{2} < \e$.
\end{proof}

TODO: terminar esta clase.

\section{Integrabilidad y medida de Jordan}

\section{Integración en conjuntos}


\clase{12/11/2024}{Sucesiones de funciones}
\section{Sucesiones de funciones}

\begin{definition}[Convergencia puntual]
  Sean $X \subset \R$ y $f_n\colon X \to \R$ para cada $n \in \N$. La sucesión ${(f_n)}_{n \in \N}$ \emph{converge puntualmente} a $f\colon X \to \R$ si para todo $x \in X$,
  \[
    \lim_{n \to +\infty} f_n(x) = f(x).
  \]
\end{definition}

\begin{eg}
  Las funciones $f_n\colon [0, 1] \to \R$, $f_n(x) = x^n$, convergen puntualmente a la función $f\colon [0, 1] \to \R$ dada por
  \[
    f(x) = \begin{cases}
      0 & \text{si } 0 \leq x < 1, \\
      1 & \text{si } x = 1.
    \end{cases}
  \]
  En efecto, si $x \in [0, 1)$, $\lim_{n \to \infty} x^n = 0$, y si $x = 1$, $\lim_{n \to \infty} 1^n = 1$.

  \begin{center}
    \begin{tikzpicture}
      \begin{axis}[
          width=10cm, height=7cm,
          xlabel={$x$},
          ylabel={$f_n(x)$},
          xmin=0, xmax=1,
          ymin=0, ymax=1,
          legend pos=north west,
          grid=major,
          title={Funciones $f_n(x) = x^n$ en $[0, 1]$}
        ]
        \addplot[domain=0:1, samples=100, thick, blue] {x^1};
        \addlegendentry{$f_1$}

        \addplot[domain=0:1, samples=100, thick, red] {x^2};
        \addlegendentry{$f_2$}

        \addplot[domain=0:1, samples=100, thick, green!60!black] {x^3};
        \addlegendentry{$f_3$}

        \addplot[domain=0:1, samples=100, thick, orange] {x^5};
        \addlegendentry{$f_5$}

        \addplot[domain=0:1, samples=100, thick, purple] {x^10};
        \addlegendentry{$f_{10}$}

        \addplot[domain=0:1, samples=100, thick, brown] {x^50};
        \addlegendentry{$f_{50}$}
      \end{axis}
    \end{tikzpicture}
  \end{center}
\end{eg}

\begin{eg}
  Las funciones $f_n\colon [0, 2\pi] \to \R$, $f_n(x) = \cos(nx)$, no convergen puntualmente: para $x = \pi$, $f_n(\pi) = {(-1)}^n$, que no tiene límite.
\end{eg}

Si $f_n\colon X \to \R$ converge puntualmente a $f$, para cada $\e > 0$ y cada $x \in X$ existe $n_0$ tal que $|f_n(x) - f(x)| < \e$ para todo $n > n_0$, pero $n_0$ depende de $\e$ y de $x$. Para $\e$ fijo puede ocurrir que ningún $n_0$ sirva para todo $x \in X$ a la vez.

\begin{eg}
  Sea $f_n(x) = x^n$ en $[0, 1]$, que converge puntualmente a la $f$ del ejemplo anterior. Para $\e = \frac{1}{2}$ y cada $n \in \N$, el punto $x = 2^{-1/n} \in [0, 1)$ cumple $|f_n(x) - f(x)| = x^n = \frac{1}{2} \geq \e$. Luego ningún $n_0$ sirve para todos los $x \in [0, 1)$ simultáneamente.
\end{eg}

\begin{definition}[Convergencia uniforme]
  La sucesión ${(f_n)}_{n \in \N}$ \emph{converge uniformemente} a $f\colon X \to \R$ si
  \[
    (\forall \e > 0)(\exists n_0 \in \N) \colon (\forall n > n_0)(\forall x \in X) \quad |f_n(x) - f(x)| < \e.
  \]
  Notamos $f_n \rightrightarrows f$.
\end{definition}

De la definición se deduce que $f_n \rightrightarrows f$ si y sólo si
\[
  \lim_{n \to +\infty} \sup_{x \in X} |f_n(x) - f(x)| = 0.
\]

\begin{eg}
  Por el ejemplo anterior, $f_n(x) = x^n$ no converge uniformemente en $[0, 1]$. Sea ahora $\delta \in (0, 1)$ y consideremos $f_n\colon [0, 1 - \delta] \to \R$; veamos que la convergencia a $f \equiv 0$ es uniforme. Para $x \in [0, 1 - \delta]$, $0 \leq x^n \leq {(1 - \delta)}^n$. Como $\lim_{n \to \infty} {(1 - \delta)}^n = 0$, dado $\e > 0$ existe $n_0$ con ${(1 - \delta)}^n < \e$ para $n > n_0$, y entonces $|x^n - 0| < \e$ para todo $n > n_0$ y todo $x \in [0, 1 - \delta]$.
\end{eg}

\begin{definition}[Sucesión de Cauchy de funciones]
  La sucesión ${(f_n)}_{n \in \N}$, con $f_n\colon X \to \R$, es \emph{de Cauchy} (uniformemente) si
  \[
    (\forall \e > 0)(\exists n_0 \in \N) \colon (\forall n, m > n_0)(\forall x \in X) \quad |f_n(x) - f_m(x)| < \e.
  \]
\end{definition}

\begin{theorem}
  Una sucesión de funciones $f_n\colon X \to \R$ converge uniformemente si y sólo si es de Cauchy.
\end{theorem}

\begin{proof}
  ($\Rightarrow$) Sea $f_n \rightrightarrows f$ y $\e > 0$. Existe $n_0$ tal que $|f_n(x) - f(x)| < \frac{\e}{2}$ para $n > n_0$ y todo $x \in X$. Si $n, m > n_0$,
  \[
    |f_n(x) - f_m(x)| \leq |f_n(x) - f(x)| + |f(x) - f_m(x)| < \e \qquad (\forall x \in X).
  \]

  ($\Leftarrow$) Para cada $x \in X$, ${(f_n(x))}_{n \in \N}$ es de Cauchy en $\R$, y por lo tanto converge; llamamos $f(x)$ a su límite. Dado $\e > 0$, existe $n_0$ tal que $|f_n(x) - f_m(x)| < \e$ para $n, m > n_0$ y todo $x \in X$. Fijando $x$ y $n > n_0$ y haciendo $m \to +\infty$ se obtiene $|f_n(x) - f(x)| \leq \e$ para todo $n > n_0$ y todo $x \in X$. Luego $f_n \rightrightarrows f$.
\end{proof}

Una serie de funciones $\sum f_n$ es un caso particular: sea $S_n = f_1 + \cdots + f_n$ la suma parcial. Decimos que $\sum f_n$ converge puntual o uniformemente a $f$ si $S_n$ lo hace, y en ese caso escribimos $\sum f_n \rightrightarrows f$. Por el teorema anterior, $\sum f_n$ converge uniformemente si y sólo si
\[
  (\forall \e > 0)(\exists n_0 \in \N) \colon (\forall n > n_0)(\forall p \in \N)(\forall x \in X) \quad |S_{n+p}(x) - S_n(x)| < \e.
\]

\begin{definition}[Convergencia normal]
  Decimos que una serie $\sum f_n$, con $f_n\colon X \to \R$, \emph{converge normalmente} si existe una sucesión ${(a_n)}_{n \in \N} \subset [0, +\infty)$ con $\sum a_n < +\infty$ y $|f_n(x)| \leq a_n$ para todo $n \in \N$ y todo $x \in X$.
\end{definition}

\begin{eg}
  La serie $\sum_{n \geq 1} \frac{\sin(nx)}{n^2}$ converge normalmente en $\R$, pues $\left| \frac{\sin(nx)}{n^2} \right| \leq \frac{1}{n^2}$ para todo $n$ y $\sum \frac{1}{n^2}$ converge.
\end{eg}

\begin{theorem}[Criterio de Weierstrass]
  Si $\sum f_n$ converge normalmente, entonces $\sum f_n$ y $\sum |f_n|$ convergen uniformemente.
\end{theorem}

\begin{proof}
  Sea $\e > 0$. Como $\sum a_n$ converge, por el criterio de Cauchy existe $n_0$ tal que $a_{n+1} + \cdots + a_{n+p} < \e$ para todo $n > n_0$ y todo $p \in \N$. Entonces, para todo $x \in X$,
  \begin{multline*}
    |S_{n+p}(x) - S_n(x)| = |f_{n+1}(x) + \cdots + f_{n+p}(x)| \\
    \leq |f_{n+1}(x)| + \cdots + |f_{n+p}(x)| \leq a_{n+1} + \cdots + a_{n+p} < \e.
  \end{multline*}
  Luego $\sum f_n$ converge uniformemente; el argumento para $\sum |f_n|$ es el mismo.
\end{proof}

La convergencia normal implica la convergencia uniforme, que a su vez implica la convergencia puntual, pero las recíprocas no valen.

\begin{eg}
  Sea $f_n\colon [1, +\infty) \to \R$ dada por
  \[
    f_n(x) = \begin{cases}
      \dfrac{1}{x} & \text{si } x \in [n, n + 1), \\[1ex]
      0            & \text{en caso contrario,}
    \end{cases}
  \]
  y sea $f(x) = \frac{1}{x}$. Afirmamos que $\sum_{n \geq 1} f_n \rightrightarrows f$. En efecto, $S_n(x) = \frac{1}{x}$ si $x \in [1, n + 1)$ y $S_n(x) = 0$ si $x \geq n + 1$. Luego $f(x) - S_n(x) = 0$ si $x < n + 1$, y $f(x) - S_n(x) = \frac{1}{x} \leq \frac{1}{n + 1} < \frac{1}{n}$ si $x \geq n + 1$; así, $0 \leq f - S_n < \frac{1}{n}$ en $[1, +\infty)$.

  Pero la serie no converge normalmente: si $|f_n(x)| \leq a_n$ para todo $x$, tomando $x = n$ resulta $a_n \geq \frac{1}{n}$, y $\sum a_n = +\infty$.
\end{eg}

\begin{note}
  Dada una sucesión ${(f_n)}_{n \in \N}$, nos preguntamos si los límites en $x$ y en $n$ pueden intercambiarse. En general, la respuesta es negativa.
\end{note}

\begin{eg}
  Sean $f_n(x) = x^n$ en $[0, 1]$ y $f$ su límite puntual. Entonces
  \begin{align*}
    \lim_{x \to 1} \left( \lim_{n \to \infty} f_n(x) \right) & = \lim_{x \to 1} f(x) = 0, \\
    \lim_{n \to \infty} \left( \lim_{x \to 1} f_n(x) \right) & = \lim_{n \to \infty} 1 = 1.
  \end{align*}
\end{eg}

\begin{theorem}
  Sea $a$ un punto de acumulación de $X \subset \R$ (o $a = \pm\infty$). Si $f_n \rightrightarrows f$ en $X$ y para cada $n \in \N$ existe $L_n = \lim_{x \to a} f_n(x)$, entonces existe $L = \lim_{n \to \infty} L_n$ y $\lim_{x \to a} f(x) = L$. Es decir,
  \[
    \lim_{n \to \infty} \left( \lim_{x \to a} f_n(x) \right) = \lim_{x \to a} \left( \lim_{n \to \infty} f_n(x) \right).
  \]
\end{theorem}

\begin{proof}
  Veamos primero que ${(L_n)}$ es de Cauchy. Dado $\e > 0$, existe $n_0$ tal que $|f_n(x) - f_m(x)| < \frac{\e}{3}$ para $n, m > n_0$ y todo $x \in X$. Fijemos $n, m > n_0$. Como $f_n(x) \to L_n$ y $f_m(x) \to L_m$ cuando $x \to a$, y $a$ es de acumulación, existe $x \in X$ tal que $|f_n(x) - L_n| < \frac{\e}{3}$ y $|f_m(x) - L_m| < \frac{\e}{3}$. Entonces
  \[
    |L_n - L_m| \leq |L_n - f_n(x)| + |f_n(x) - f_m(x)| + |f_m(x) - L_m| < \e.
  \]
  Luego ${(L_n)}$ converge a un $L \in \R$.

  Veamos ahora que $\lim_{x \to a} f(x) = L$. Dado $\e > 0$, existen $n_1$ con $|f_n(x) - f(x)| < \frac{\e}{3}$ para $n > n_1$ y todo $x \in X$, y $n_2$ con $|L_n - L| < \frac{\e}{3}$ para $n > n_2$. Sea $n > \max\{n_1, n_2\}$. Como $f_n(x) \to L_n$ cuando $x \to a$, existe $\delta > 0$ tal que $0 < |x - a| < \delta$ implica $|f_n(x) - L_n| < \frac{\e}{3}$. Para tales $x$,
  \[
    |f(x) - L| \leq |f(x) - f_n(x)| + |f_n(x) - L_n| + |L_n - L| < \e.
  \]
  Si $a = \pm\infty$, la demostración es la misma reemplazando la condición $0 < |x - a| < \delta$ por $x > R$ (o $x < -R$).
\end{proof}

\begin{corollary}
  Sea $a$ un punto de acumulación de $X$. Si $\sum f_n \rightrightarrows f$ en $X$ y para cada $n$ existe $L_n = \lim_{x \to a} f_n(x)$, entonces $\sum L_n$ converge y $\sum L_n = \lim_{x \to a} f(x)$. Es decir,
  \[
    \lim_{x \to a} \left( \sum f_n(x) \right) = \sum \left( \lim_{x \to a} f_n(x) \right).
  \]
\end{corollary}

\begin{proof}
  Se aplica el teorema a las sumas parciales $S_n$, que convergen uniformemente a $f$ y tienen $\lim_{x \to a} S_n(x) = L_1 + \cdots + L_n$.
\end{proof}

\begin{theorem}
  Sean $f_n\colon X \to \R$ funciones continuas en $a \in X$ tales que $f_n \rightrightarrows f$. Entonces $f$ es continua en $a$.
\end{theorem}

\begin{proof}
  Si $a$ es aislado en $X$, $f$ es continua en $a$ trivialmente. Si $a \in X'$, por continuidad $\lim_{x \to a} f_n(x) = f_n(a)$ para cada $n$, y por el teorema anterior,
  \[
    \lim_{x \to a} f(x) = \lim_{n \to \infty} f_n(a) = f(a). \qedhere
  \]
\end{proof}

En particular, si las $f_n$ son continuas pero el límite puntual no lo es, la convergencia no es uniforme. Por otro lado, que el límite sea continuo no alcanza para que la convergencia sea uniforme: $f_n(x) = x^n$ converge puntualmente a $0$ en $[0, 1)$, pero no uniformemente.

\begin{definition}[Convergencia monótona]
  Decimos que ${(f_n)}_{n \in \N}$ \emph{converge monótonamente} a $f$ si para todo $x \in X$ la sucesión ${(f_n(x))}_{n \in \N}$ es monótona y $f_n(x) \to f(x)$.
\end{definition}

\begin{theorem}[Dini]
  Sea $X \subset \R$ compacto. Si una sucesión de funciones continuas $f_n\colon X \to \R$ converge monótonamente a una función continua $f\colon X \to \R$, entonces la convergencia es uniforme.
\end{theorem}

\begin{proof}
  Dado $\e > 0$ definimos $K_n = \{ x \in X : |f_n(x) - f(x)| \geq \e \}$. Como $|f_n - f|$ es continua, $K_n$ es cerrado relativo a $X$, y por lo tanto compacto. Para cada $x$, como ${(f_n(x))}$ es monótona y converge a $f(x)$, la sucesión ${\bigl( |f_n(x) - f(x)| \bigr)}_n$ es decreciente, de modo que $K_1 \supset K_2 \supset K_3 \supset \cdots$. Además $\bigcap_{n} K_n = \varnothing$, pues si $x$ pertenece a todos los $K_n$, entonces $|f_n(x) - f(x)| \geq \e$ para todo $n$, lo que contradice $f_n(x) \to f(x)$.

  Supongamos que $K_n \neq \varnothing$ para todo $n$. Elegimos $x_n \in K_n$. Como $X$ es compacto, existe una subsucesión $x_{h(k)} \to x \in X$. Para cada $m$ fijo, $x_{h(k)} \in K_{h(k)} \subset K_m$ para $k$ grande, y como $K_m$ es cerrado relativo a $X$, $x \in K_m$. Entonces $x \in \bigcap_m K_m = \varnothing$, absurdo.

  Luego existe $n_0$ con $K_{n_0} = \varnothing$, y entonces $K_n = \varnothing$ para $n \geq n_0$, es decir, $|f_n(x) - f(x)| < \e$ para $n \geq n_0$ y todo $x \in X$.
\end{proof}

\begin{corollary}
  Sea $X \subset \R$ compacto y $f_n\colon X \to \R$ funciones continuas no negativas tales que $\sum f_n(x)$ converge para todo $x \in X$. Entonces $\sum f_n$ converge uniformemente si y sólo si su suma es continua.
\end{corollary}

\begin{proof}
  ($\Rightarrow$) Es el teorema de continuidad del límite uniforme. ($\Leftarrow$) Las sumas parciales $S_n = f_1 + \cdots + f_n$ son continuas y forman una sucesión creciente, pues las $f_n$ son no negativas, que converge puntualmente a una función continua. Por el teorema de Dini, la convergencia es uniforme.
\end{proof}

\section{Integración de sucesiones de funciones}

\begin{theorem}
  Si $f_n\colon [a, b] \to \R$ son integrables y $f_n \rightrightarrows f$, entonces $f$ es integrable y
  \[
    \int_a^b f(x) \,\dd x = \lim_{n \to +\infty} \int_a^b f_n(x) \,\dd x.
  \]
\end{theorem}

\begin{proof}
  Como las $f_n$ son acotadas y $f_n \rightrightarrows f$, $f$ es acotada. Sea $\e > 0$ y sea $\eta = \frac{\e}{4(b - a)}$. Existe $n$ con $|f_n(x) - f(x)| < \eta$ para todo $x \in [a, b]$. En cada subintervalo $[t_{i-1}, t_i]$ de una partición $P$ se tiene entonces $\omega_i(f) \leq \omega_i(f_n) + 2\eta$, de modo que
  \[
    S(f, P) - s(f, P) \leq S(f_n, P) - s(f_n, P) + 2\eta (b - a).
  \]
  Como $f_n$ es integrable, existe $P$ con $S(f_n, P) - s(f_n, P) < \frac{\e}{2}$, y entonces $S(f, P) - s(f, P) < \frac{\e}{2} + \frac{\e}{2} = \e$. Luego $f$ es integrable.

  Ahora, dado $\e > 0$ existe $n_0$ tal que $|f_n(x) - f(x)| < \frac{\e}{b - a}$ para $n > n_0$ y todo $x \in [a, b]$. Para esos $n$,
  \[
    \left| \int_a^b f \,\dd x - \int_a^b f_n \,\dd x \right| \leq \int_a^b |f - f_n| \,\dd x \leq \frac{\e}{b - a} (b - a) = \e. \qedhere
  \]
\end{proof}

\begin{corollary}
  Si $\sum f_n \rightrightarrows f$ en $[a, b]$ y las $f_n$ son integrables, entonces $f$ es integrable y $\int_a^b \sum f_n = \sum \int_a^b f_n$. Es decir, la serie puede integrarse término a término.
\end{corollary}

\begin{note}
  Si $f_n\colon [a, b] \to \R$ son integrables y convergen puntualmente a $f$, puede ocurrir que $f$ no sea integrable. Por ejemplo, si $\Q \cap [a, b] = \{ r_1, r_2, \dots \}$ y
  \[
    f_n(x) = \begin{cases}
      1 & \text{si } x \in \{ r_1, \dots, r_n \}, \\
      0 & \text{en caso contrario,}
    \end{cases}
  \]
  entonces $f_n \to f$ puntualmente, con $f(x) = 1$ si $x \in \Q \cap [a, b]$ y $f(x) = 0$ en caso contrario. Cada $f_n$ es integrable (tiene finitas discontinuidades), pero $f$ no lo es.
\end{note}

\begin{note}
  Aun cuando $f_n$ y $f$ sean integrables y $f_n \to f$ puntualmente, puede ser $\lim_{n \to \infty} \int_a^b f_n \neq \int_a^b f$. Por ejemplo, sea $f_n\colon [0, 1] \to \R$ dada por $f_n(x) = (n + 1) x^n$ si $x \in [0, 1)$ y $f_n(1) = 0$. Entonces $f_n \to 0$ puntualmente y $\int_0^1 f_n(x) \,\dd x = 1$ para todo $n$, pero $\int_0^1 0 \,\dd x = 0$.
\end{note}


\clase{15/11/2024}{Series de potencias y equicontinuidad}
\section{Diferenciación de sucesiones de funciones}

\begin{theorem}
  Sea ${(f_n)}_{n \in \N}$ una sucesión de funciones derivables en $[a, b]$. Si existe $c \in [a, b]$ tal que ${(f_n(c))}_{n \in \N}$ converge y las derivadas $f_n'$ convergen uniformemente en $[a, b]$ a una función $g$, entonces $f_n \rightrightarrows f$ en $[a, b]$ para alguna función $f$, y $f' = g$. Es decir,
  \[
    {\left( \lim_{n \to +\infty} f_n \right)}' = \lim_{n \to +\infty} f_n'
  \]
  siempre que las derivadas converjan uniformemente (y las $f_n$ converjan en algún punto).
\end{theorem}

\begin{proof}
  \emph{Paso 1: ${(f_n)}$ converge uniformemente.} Aplicando el teorema del valor medio a $f_m - f_n$ en el intervalo de extremos $x$ y $c$, existe $d$ entre $x$ y $c$ tal que
  \[
    \bigl( f_m(x) - f_n(x) \bigr) - \bigl( f_m(c) - f_n(c) \bigr) = (x - c) \bigl( f_m'(d) - f_n'(d) \bigr).
  \]
  Luego, para todo $x \in [a, b]$,
  \[
    |f_m(x) - f_n(x)| \leq |f_m(c) - f_n(c)| + (b - a) \sup_{[a, b]} |f_m' - f_n'|.
  \]
  Como ${(f_n(c))}$ converge y ${(f_n')}$ es uniformemente de Cauchy, dado $\e > 0$ ambos términos son menores que $\frac{\e}{2}$ si $n, m$ son grandes, y entonces ${(f_n)}$ es uniformemente de Cauchy. Por lo tanto $f_n \rightrightarrows f$ para alguna $f$.

  \emph{Paso 2: $f' = g$.} Fijemos $x_0 \in [a, b]$ y definamos, para $x \in [a, b] \setminus \{x_0\}$,
  \[
    q_n(x) = \frac{f_n(x) - f_n(x_0)}{x - x_0}, \qquad q(x) = \frac{f(x) - f(x_0)}{x - x_0}.
  \]
  Por el teorema del valor medio aplicado a $f_m - f_n$ entre $x_0$ y $x$, existe $d$ entre $x_0$ y $x$ con $q_m(x) - q_n(x) = f_m'(d) - f_n'(d)$, de modo que $|q_m(x) - q_n(x)| \leq \sup_{[a, b]} |f_m' - f_n'|$ para todo $x$. Entonces ${(q_n)}$ es uniformemente de Cauchy en $[a, b] \setminus \{x_0\}$, y como $q_n \to q$ puntualmente, $q_n \rightrightarrows q$. Además $\lim_{x \to x_0} q_n(x) = f_n'(x_0)$. Por el teorema de intercambio de límites,
  \[
    \lim_{x \to x_0} q(x) = \lim_{n \to \infty} f_n'(x_0) = g(x_0),
  \]
  es decir, $f'(x_0) = g(x_0)$.
\end{proof}

\begin{corollary}
  Sean $f_n$ derivables en $[a, b]$. Si $\sum f_n(c)$ converge para algún $c \in [a, b]$ y $\sum f_n' \rightrightarrows g$, entonces $\sum f_n \rightrightarrows f$ en $[a, b]$, con $f$ derivable y $f' = g$.
\end{corollary}

\begin{proof}
  Se aplica el teorema a las sumas parciales $S_n = f_1 + \cdots + f_n$.
\end{proof}

\section{Series de potencias}

Una serie de la forma $\sum_{n \geq 0} a_n {(x - x_0)}^n$ se llama \emph{serie de potencias}. Por simplicidad consideramos $\sum_{n \geq 0} a_n x^n$; el caso general se reduce a éste mediante el cambio de variable $x \mapsto x - x_0$.

\begin{theorem}[Cauchy--Hadamard]
  Sea $\sum_{n \geq 0} a_n x^n$ una serie de potencias y sea $r \in [0, +\infty]$ definido por $\frac{1}{r} = \limsup_{n \to \infty} \sqrt[n]{|a_n|}$ (con $\frac{1}{0} = +\infty$ y $\frac{1}{+\infty} = 0$). Entonces la serie converge absolutamente si $|x| < r$ y diverge si $|x| > r$. El número $r$ se llama \emph{radio de convergencia}.
\end{theorem}

\begin{proof}
  Sea $\rho = \limsup \sqrt[n]{|a_n|} = \frac{1}{r}$. Para $x \neq 0$, $\limsup \sqrt[n]{|a_n x^n|} = |x| \rho$.

  Si $|x| \rho < 1$, es decir $|x| < r$ (esto incluye $\rho = 0$, donde converge para todo $x$), por el criterio de la raíz la serie converge absolutamente. Si $|x| \rho > 1$, es decir $|x| > r$ (esto incluye $\rho = +\infty$ y $x \neq 0$), existen infinitos $n$ con $\sqrt[n]{|a_n x^n|} > 1$, de modo que $|a_n x^n| > 1$ para infinitos $n$, el término general no tiende a $0$ y la serie diverge.
\end{proof}

\begin{theorem}
  Una serie de potencias converge uniformemente en todo intervalo compacto contenido en su intervalo de convergencia.
\end{theorem}

\begin{proof}
  Sea $r > 0$ el radio de convergencia de $\sum a_n x^n$ y sea $0 < s < r$. Para $x \in [-s, s]$ vale $|a_n x^n| \leq |a_n| s^n$, y como $\sum |a_n| s^n$ converge (la serie converge absolutamente en $s$), por el criterio de Weierstrass $\sum a_n x^n$ converge uniformemente en $[-s, s]$. Todo intervalo compacto contenido en $(-r, r)$ está contenido en algún $[-s, s]$.
\end{proof}

\begin{corollary}
  La función $f\colon (-r, r) \to \R$, $f(x) = \sum a_n x^n$, es continua.
\end{corollary}

\begin{proof}
  Dado $x_0 \in (-r, r)$, existe $s$ con $|x_0| < s < r$. En $[-s, s]$ la serie es un límite uniforme de polinomios, que son continuos, y por lo tanto $f$ es continua en $x_0$.
\end{proof}

Una serie de potencias puede no converger uniformemente en todo su intervalo de convergencia $(-r, r)$.

\begin{eg}
  La serie $\sum x^n$ no converge uniformemente en $(-1, 1)$. En efecto, si $\sum f_n \rightrightarrows f$ entonces $f_{n+1} = S_{n+1} - S_n \rightrightarrows 0$, pero $\sup_{x \in (-1, 1)} |x^n| = 1$ para todo $n$.
\end{eg}

\begin{lemma}
  Sea $\sum_p \alpha_p$ una serie con sumas parciales acotadas: si $S_p = \alpha_1 + \cdots + \alpha_p$, existe $k > 0$ con $|S_p| \leq k$ para todo $p \in \N$. Si $b_1 \geq b_2 \geq \cdots \geq b_p \geq 0$, entonces
  \[
    |\alpha_1 b_1 + \cdots + \alpha_p b_p| \leq k b_1.
  \]
\end{lemma}

\begin{proof}
  Con $S_0 = 0$ y $\alpha_i = S_i - S_{i-1}$,
  \begin{align*}
    |\alpha_1 b_1 + \cdots + \alpha_p b_p| & = |S_1 b_1 + (S_2 - S_1) b_2 + \cdots + (S_p - S_{p-1}) b_p|              \\
                                           & = |S_1 (b_1 - b_2) + \cdots + S_{p-1} (b_{p-1} - b_p) + S_p b_p|          \\
                                           & \leq k \bigl( (b_1 - b_2) + \cdots + (b_{p-1} - b_p) + b_p \bigr) = k b_1. \qedhere
  \end{align*}
\end{proof}

\begin{theorem}[Abel]
  Sea $\sum a_n x^n$ con radio de convergencia $0 < r < +\infty$. Si $\sum a_n r^n$ converge, entonces $\sum a_n x^n$ converge uniformemente en $[0, r]$. En particular, $\lim_{x \to r^{-}} \sum a_n x^n = \sum a_n r^n$.
\end{theorem}

\begin{proof}
  Sea $\e > 0$. Como $\sum a_n r^n$ converge, existe $n_0$ tal que $|a_{n+1} r^{n+1} + \cdots + a_{n+p} r^{n+p}| < \e$ para todo $p \in \N$ y todo $n > n_0$. Fijemos $n > n_0$ y $x \in [0, r]$, y sean $\alpha_p = a_{n+p} r^{n+p}$, cuyas sumas parciales están acotadas en módulo por $k = \e$, y $b_p = {\left( \frac{x}{r} \right)}^{n+p}$, que cumplen $b_1 \geq b_2 \geq \cdots \geq 0$ pues $\frac{x}{r} \in [0, 1]$. Como $a_{n+p} x^{n+p} = \alpha_p b_p$, el lema da
  \[
    |a_{n+1} x^{n+1} + \cdots + a_{n+p} x^{n+p}| \leq \e \, b_1 = \e {\left( \frac{x}{r} \right)}^{n+1} \leq \e
  \]
  para todo $p$ y todo $x \in [0, r]$. Luego la serie converge uniformemente en $[0, r]$. Entonces su suma $f$ es continua en $[0, r]$, por ser límite uniforme de polinomios, y $\sum a_n r^n = f(r) = \lim_{x \to r^{-}} f(x)$.
\end{proof}

\begin{note}
  \begin{enumerate}[label=(\arabic*)]
    \item El mismo resultado vale en $-r$: si $\sum a_n {(-r)}^n$ converge, aplicando el teorema a la serie $\sum {(-1)}^n a_n x^n$ se obtiene que $\sum a_n x^n$ converge uniformemente en $[-r, 0]$.
    \item $\sum a_n x^n$ converge uniformemente en $(-r, r)$ si y sólo si converge en $r$ y en $-r$. En efecto, si converge en ambos extremos, por el teorema de Abel converge uniformemente en $[0, r]$ y en $[-r, 0]$, y por lo tanto en $[-r, r]$. Recíprocamente, si converge uniformemente en $(-r, r)$, para $\e > 0$ existe $n_0$ con $|a_{n+1} x^{n+1} + \cdots + a_{n+p} x^{n+p}| < \e$ para todo $x \in (-r, r)$, $n > n_0$ y todo $p$, y haciendo $x \to \pm r$ resulta que la serie es de Cauchy en $\pm r$.
  \end{enumerate}
\end{note}

\begin{theorem}
  Sea $[\alpha, \beta] \subset [-r, r]$ tal que $\sum a_n x^n$ converge en $\alpha$ y en $\beta$. Entonces
  \[
    \int_{\alpha}^{\beta} \sum_{n \geq 0} a_n x^n \,\dd x = \sum_{n \geq 0} \frac{a_n}{n + 1} \left( \beta^{n+1} - \alpha^{n+1} \right).
  \]
\end{theorem}

\begin{proof}
  Por el criterio de Weierstrass y el teorema de Abel (según si $\alpha$ o $\beta$ son $\pm r$ o están en $(-r, r)$), la serie converge uniformemente en $[\alpha, \beta]$. Por el teorema de integración de series de funciones, puede integrarse término a término.
\end{proof}

Si $\sum a_n x^n$ no converge en un extremo $r$ de su intervalo de convergencia, igualmente puede integrarse término a término la integral impropia $\int_0^r \sum a_n x^n \,\dd x$, siempre que $\sum \frac{a_n r^{n+1}}{n+1}$ converja. En efecto, la serie $\sum \frac{a_n}{n+1} t^{n+1}$ tiene el mismo radio $r$ y converge en $r$, de modo que por el teorema de Abel,
\[
  \int_0^r \sum a_n x^n \,\dd x = \lim_{t \to r^{-}} \int_0^t \sum a_n x^n \,\dd x = \lim_{t \to r^{-}} \sum \frac{a_n}{n + 1} t^{n+1} = \sum \frac{a_n r^{n+1}}{n + 1}.
\]

\begin{theorem}
  La función $f(x) = \sum a_n x^n$ es derivable en todo el intervalo de convergencia $(-r, r)$, y $f'(x) = \sum n a_n x^{n-1}$, serie que tiene el mismo radio de convergencia.
\end{theorem}

\begin{proof}
  Como $\sqrt[n]{n} \to 1$, se tiene $\limsup \sqrt[n]{n |a_n|} = \limsup \sqrt[n]{|a_n|}$. Luego la serie $\sum n a_n x^n$, y por lo tanto $\sum n a_n x^{n-1}$ (que para $x \neq 0$ es $\frac{1}{x} \sum n a_n x^n$), tiene el mismo radio de convergencia $r$.

  Sea $0 < s < r$. La serie de derivadas $\sum n a_n x^{n-1}$ converge uniformemente en $[-s, s]$ por el teorema de convergencia uniforme de series de potencias, y $\sum a_n x^n$ converge en $c = 0$. Por el corolario del teorema de diferenciación de series de funciones, $f$ es derivable en $[-s, s]$ con $f'(x) = \sum n a_n x^{n-1}$. Como $s < r$ es arbitrario, vale en todo $(-r, r)$.
\end{proof}

\begin{corollary}
  Si $f(x) = \sum a_n x^n$ tiene radio de convergencia $r > 0$, entonces $f$ tiene derivadas de todos los órdenes en $(-r, r)$, que se calculan derivando término a término:
  \[
    f^{(k)}(x) = \sum_{n \geq k} n (n - 1) \cdots (n - k + 1)\, a_n x^{n-k}.
  \]
  En particular, $a_k = \dfrac{f^{(k)}(0)}{k!}$, y la serie de potencias que converge a $f$ en $(-r, r)$ es su serie de Taylor alrededor de $0$.
\end{corollary}

\section{Familias equicontinuas}

Dada una familia $E$ de funciones continuas con dominio común, queremos saber cuándo podemos garantizar que toda sucesión ${(f_n)}_{n \in \N} \subset E$ tiene una subsucesión uniformemente convergente.

A diferencia de lo que ocurre con sucesiones numéricas, no alcanza con que la familia esté uniformemente acotada. Por ejemplo, sea $f_n\colon [0, 1] \to \R$, $f_n(x) = x^n (1 - x^n)$. Entonces $0 \leq f_n \leq \frac{1}{4}$ y $f_n \to 0$ puntualmente, pero $f_n \bigl( 2^{-1/n} \bigr) = \frac{1}{2} \cdot \frac{1}{2} = \frac{1}{4}$ para todo $n$. Luego ninguna subsucesión puede converger uniformemente a $0$ y, como el límite uniforme debería coincidir con el puntual, ninguna subsucesión converge uniformemente.

\begin{definition}[Equicontinuidad]
  Sean $X \subset \R$ y $E$ una familia de funciones $f\colon X \to \R$. Dado $x_0 \in X$, decimos que $E$ es \emph{equicontinua en $x_0$} si para todo $\e > 0$ existe $\delta > 0$ tal que $|x - x_0| < \delta \Rightarrow |f(x) - f(x_0)| < \e$ para toda $f \in E$ (el mismo $\delta$ para todas las funciones). Decimos que $E$ es \emph{equicontinua} si lo es en todo punto de $X$.
\end{definition}

\begin{theorem}
  Sean $f_n\colon X \to \R$ continuas y $f_n \rightrightarrows f$. Entonces la sucesión ${(f_n)}_{n \in \N}$ es equicontinua.
\end{theorem}

\begin{proof}
  La función $f$ es continua por ser límite uniforme de continuas. Sean $x_0 \in X$ y $\e > 0$. Existe $n_0$ tal que $|f_n(x) - f(x)| < \frac{\e}{3}$ para todo $n > n_0$ y todo $x \in X$. Por la continuidad de $f$ existe $\delta_0 > 0$ tal que $|x - x_0| < \delta_0 \Rightarrow |f(x) - f(x_0)| < \frac{\e}{3}$, y por la de cada $f_i$ existe $\delta_i > 0$ tal que $|x - x_0| < \delta_i \Rightarrow |f_i(x) - f_i(x_0)| < \e$, para $i \in \{1, \dots, n_0\}$. Sea $\delta = \min\{ \delta_0, \delta_1, \dots, \delta_{n_0} \}$. Si $|x - x_0| < \delta$ y $n \leq n_0$, entonces $|f_n(x) - f_n(x_0)| < \e$. Si $n > n_0$,
  \[
    |f_n(x) - f_n(x_0)| \leq |f_n(x) - f(x)| + |f(x) - f(x_0)| + |f(x_0) - f_n(x_0)| < \e. \qedhere
  \]
\end{proof}

\begin{eg}
  La sucesión $f_n(x) = \dfrac{\sin(nx)}{n}$ es equicontinua en $\R$, pues $|f_n(x)| \leq \frac{1}{n}$ implica $f_n \rightrightarrows 0$.
\end{eg}

\begin{eg}
  La sucesión $f_n(x) = nx$ no es equicontinua en ningún $x_0 \in \R$. Sea $\e = \frac{1}{2}$ y $\delta > 0$. Existe $n \in \N$ con $\frac{1}{n} < \delta$, y el punto $x = x_0 + \frac{1}{n}$ cumple $|x - x_0| = \frac{1}{n} < \delta$, pero $|f_n(x) - f_n(x_0)| = n \cdot \frac{1}{n} = 1 \geq \e$.
\end{eg}


\clase{22/11/2024}{Teorema de Arzelà--Ascoli}
\section{Equicontinuidad uniforme}

\begin{theorem}
  Sea \(E\) una familia de funciones derivables en un intervalo \(I\). Si existe \(c > 0\) tal que \(|f'(x)| \leq c\) para toda \(f \in E\) y todo \(x \in I\), entonces \(E\) es equicontinua en \(I\).
\end{theorem}

\begin{proof}
  Sean \(x_0 \in I\), \(\e > 0\) y \(\delta = \frac{\e}{c}\). Si \(x \in I\) y \(|x - x_0| < \delta\), por el teorema del valor medio, para toda \(f \in E\),
  \[
    |f(x) - f(x_0)| \leq c \, |x - x_0| < c \cdot \frac{\e}{c} = \e. \qedhere
  \]
\end{proof}

\begin{note}
  Con el mismo argumento, si para todo \(x \in I\) existen \(c_x > 0\) y un intervalo \(I_x \subseteq I\) que contiene a \(x\) en su interior relativo a \(I\), tales que \(|f'(y)| \leq c_x\) para toda \(f \in E\) y todo \(y \in I_x\), entonces \(E\) es equicontinua en \(I\). Basta tomar \(\delta = \min\{ \frac{\e}{c_x}, \eta \}\), con \(\eta > 0\) tal que \(I \cap (x - \eta, x + \eta) \subseteq I_x\).
\end{note}

\begin{definition}[Equicontinuidad uniforme]
  Una familia \(E\) de funciones \(f\colon X \to \R\) se dice \emph{uniformemente equicontinua} si para todo \(\e > 0\) existe \(\delta > 0\) tal que, si \(x, y \in X\) y \(|x - y| < \delta\), entonces \(|f(x) - f(y)| < \e\) para toda \(f \in E\).
\end{definition}

\begin{eg}
  Si \(E\) es una familia de funciones derivables en un intervalo \(I\) con \(|f'(x)| \leq c\) para toda \(f \in E\) y todo \(x \in I\), entonces \(E\) es uniformemente equicontinua: el mismo argumento con \(\delta = \frac{\e}{c}\) vale para todo par \(x, y \in I\).
\end{eg}

\begin{eg}
  Una familia formada por una única función continua que no es uniformemente continua, por ejemplo \(f(x) = x^2\) en \(\R\), es equicontinua pero no uniformemente equicontinua.
\end{eg}

\begin{theorem}
  Sea \(K \subset \R\) compacto. Toda familia equicontinua de funciones \(f\colon K \to \R\) es uniformemente equicontinua.
\end{theorem}

\begin{proof}
  Supongamos que una familia equicontinua \(E\) no es uniformemente equicontinua. Entonces existe \(\e > 0\) tal que para todo \(n \in \N\) hay \(x_n, y_n \in K\) y \(f_n \in E\) con \(|x_n - y_n| < \frac{1}{n}\) y \(|f_n(x_n) - f_n(y_n)| \geq \e\).

  Como \(K\) es compacto, pasando a una subsucesión podemos suponer que \(x_n \to x \in K\), y como \(|y_n - x_n| < \frac{1}{n}\), también \(y_n \to x\). Por la equicontinuidad de \(E\) en \(x\) existe \(\delta > 0\) tal que \(z \in K\) y \(|z - x| < \delta\) implican \(|f(z) - f(x)| < \frac{\e}{2}\) para toda \(f \in E\). Para \(n\) suficientemente grande, \(|x_n - x| < \delta\) y \(|y_n - x| < \delta\), de modo que
  \[
    |f_n(x_n) - f_n(y_n)| \leq |f_n(x_n) - f_n(x)| + |f_n(x) - f_n(y_n)| < \e,
  \]
  lo que es absurdo.
\end{proof}

\begin{theorem}
  Sea \({(f_n)}\) una sucesión equicontinua de funciones \(f_n\colon X \to \R\), con \(X \subseteq \R\). Si \({(f_n)}\) converge puntualmente en un subconjunto denso \(D \subseteq X\), entonces converge uniformemente sobre todo compacto \(K \subseteq X\).
\end{theorem}

\begin{proof}
  Sean \(K \subseteq X\) compacto y \(\e > 0\). Veamos que existe \(n_0\) tal que \(|f_m(x) - f_n(x)| < \e\) para todo \(m, n > n_0\) y todo \(x \in K\).

  Para cada \(y \in K\), por la equicontinuidad en \(y\) existe un intervalo abierto \(I_y\) centrado en \(y\) tal que \(|f_n(z) - f_n(y)| < \frac{\e}{6}\) para todo \(z \in X \cap I_y\) y todo \(n \in \N\); luego \(|f_n(x) - f_n(z)| < \frac{\e}{3}\) para todo \(x, z \in X \cap I_y\) y todo \(n\). Como \(K\) es compacto, del cubrimiento \(K \subseteq \bigcup_{y \in K} I_y\) se puede extraer un subcubrimiento finito \(K \subseteq I_1 \cup \cdots \cup I_p\), donde cada \(I_i = I_{y_i}\) con \(y_i \in K \subseteq X\).

  Como \(D\) es denso en \(X\) y \(y_i \in X \cap I_i\), que es abierto relativo a \(X\), podemos elegir \(d_i \in D \cap I_i\) para cada \(i\). Como \(f_n(d_i)\) converge, existe \(n_i\) tal que \(|f_m(d_i) - f_n(d_i)| < \frac{\e}{3}\) para todo \(m, n > n_i\). Sea \(n_0 = \max\{n_1, \dots, n_p\}\).

  Si \(m, n > n_0\) y \(x \in K\), existe \(i\) con \(x \in I_i\), y entonces
  \[
    |f_m(x) - f_n(x)| \leq |f_m(x) - f_m(d_i)| + |f_m(d_i) - f_n(d_i)| + |f_n(d_i) - f_n(x)| < \e.
  \]
  Luego \({(f_n)}\) es uniformemente de Cauchy en \(K\) y, por lo tanto, converge uniformemente en \(K\).
\end{proof}

\begin{definition}[Acotación puntual y uniforme]
  Una familia \(E\) de funciones \(f\colon X \to \R\) se dice \emph{puntualmente acotada} si para todo \(x \in X\) existe \(c_x > 0\) tal que \(|f(x)| \leq c_x\) para toda \(f \in E\). Se dice \emph{uniformemente acotada} si existe \(c > 0\) tal que \(|f(x)| \leq c\) para toda \(f \in E\) y todo \(x \in X\).
\end{definition}

\begin{theorem}[Cantor, Tijonov]
  Sea \(X \subseteq \R\) numerable. Toda sucesión puntualmente acotada de funciones \(f_n\colon X \to \R\) posee una subsucesión puntualmente convergente.
\end{theorem}

\begin{proof}
  Si \(X\) es finito el resultado se sigue de que toda sucesión acotada en \(\R\) tiene una subsucesión convergente, aplicada sucesivamente a cada punto. Sea \(X = \{x_1, x_2, \dots\}\) infinito. Como \({(f_n(x_1))}_{n \in \N}\) es acotada, tiene una subsucesión convergente: existe \(\N_1 \subseteq \N\) infinito tal que existe \(a_1 = \lim_{n \in \N_1} f_n(x_1)\). Como \({(f_n(x_2))}_{n \in \N_1}\) es acotada, existe \(\N_2 \subseteq \N_1\) infinito tal que existe \(a_2 = \lim_{n \in \N_2} f_n(x_2)\). Siguiendo así, para cada \(i \in \N\) obtenemos \(\N_1 \supseteq \N_2 \supseteq \cdots \supseteq \N_i \supseteq \cdots\), infinitos, con \(\lim_{n \in \N_i} f_n(x_i) = a_i\).

  Definimos \(\N^* \subseteq \N\) tomando como \(i\)-ésimo elemento de \(\N^*\) al \(i\)-ésimo elemento de \(\N_i\) (en orden creciente). Como \(\N_{i+1} \subseteq \N_i\), el \((i+1)\)-ésimo elemento de \(\N_{i+1}\) es mayor o igual que el \((i+1)\)-ésimo de \(\N_i\), que es mayor que el \(i\)-ésimo de \(\N_i\); luego \(\N^*\) es infinito y estrictamente creciente. Para cada \(i\), a partir del \(i\)-ésimo elemento, \({(f_n(x_i))}_{n \in \N^*}\) es una subsucesión de \({(f_n(x_i))}_{n \in \N_i}\), que converge. Por lo tanto \({(f_n)}_{n \in \N^*}\) converge en todo \(x_i \in X\).
\end{proof}

\begin{theorem}[Arzelà, Ascoli]
  Sea \(K \subset \R\) compacto. Toda sucesión equicontinua y puntualmente acotada de funciones \(f_n\colon K \to \R\) posee una subsucesión uniformemente convergente.
\end{theorem}

\begin{proof}
  Veamos primero que \(K\) tiene un subconjunto numerable y denso \(X\). Para cada \(m \in \N\), los intervalos abiertos de radio \(\frac{1}{m}\) centrados en puntos de \(K\) cubren \(K\); por compacidad bastan finitos, con centros en un conjunto finito \(X_m \subseteq K\). Entonces \(X = \bigcup_m X_m\) es numerable y todo punto de \(K\) está a distancia menor que \(\frac{1}{m}\) de \(X_m\) para cada \(m\), así que \(X\) es denso en \(K\).

  Por el teorema de Cantor--Tijonov, como \(X\) es numerable y \({(f_n)}\) es puntualmente acotada, existe una subsucesión \({(f_{n_k})}_{k \in \N}\) que converge puntualmente en \(X\). Esta subsucesión sigue siendo equicontinua, y converge puntualmente en el subconjunto denso \(X \subseteq K\). Por el teorema anterior, converge uniformemente sobre el compacto \(K\).
\end{proof}

\section{Ecuaciones diferenciales}

TODO


\capitulosinnumero{Parciales}
\section{Primer parcial, primera fecha, 14/10}

\begin{enumerate}[label=(\arabic*)]
  \item Dada \(f\colon X \to Y\), mostrar que \(f\) es inyectiva si y sólo si \(f(A \setminus B) = f(A) \setminus f(B)\) para todo par de conjuntos \(A, B \subseteq X\).
  \item Considerar \(K\) un cuerpo ordenado y \(f\colon K \to K\) definida por \(f(x) = {(1 + x^2)}^{-1}\). Calcular, en el caso que existan, \(\inf f(\Z)\) y \(\sup f(\Z)\).
  \item Considerar \(B \subseteq A\) conjuntos no vacíos de números reales. Se sabe que \(A\) está acotado superiormente y que para todo \(x \in A\) existe \(y \in B\) tal que \(x \leq y\). Demostrar que \(\sup A = \sup B\).
  \item Demostrar, usando la definición, que \(\dfrac{4^n - 6}{4^{n+1} + 10} \to \dfrac{1}{4}\).
  \item Decidir si las siguientes sucesiones son convergentes. En caso afirmativo, determinar su límite.
        \begin{enumerate}[label=(\alph*)]
          \item \(a_1 = -\dfrac{3}{2}\), \(3 a_{n+1} = 2 + a_n^3\). Dar un valor de \(a_1\) para que la sucesión converja a \(-2\).
          \item \(b_1 = 1\), \(b_{n+1} = 1 + \dfrac{1}{b_n}\).
        \end{enumerate}
  \item Considerar una sucesión \({(a_n)}\) de términos positivos tal que \(a_n \to L\). Demostrar que
        \[
          \dfrac{n}{\dfrac{1}{a_1} + \dfrac{1}{a_2} + \cdots + \dfrac{1}{a_n}} \to L.
        \]
  \item Considerar \({(a_n)}\) y \({(b_n)}\) sucesiones acotadas superiormente. Mostrar que
        \[
          \limsup (a_n + b_n) \leq \limsup a_n + \limsup b_n.
        \]
        Dar un ejemplo donde la desigualdad sea estricta.
  \item Considerar \({(a_n)}\) una sucesión de términos positivos tal que \(\sum a_n\) es convergente. Decidir si la serie \(\sum \dfrac{\sqrt{a_n}}{n}\) es convergente o divergente.
\end{enumerate}

\clearpage

\section{Primer parcial, segunda fecha, 28/10}

\begin{enumerate}[label=(\arabic*)]
  \item Considerar \({(a_n)}\) y \({(b_n)}\) tales que \(a_n \to M\) y \(b_n \to L \neq 0\). Demostrar, usando la definición, que \(\dfrac{a_n}{b_n} \to \dfrac{M}{L}\).
  \item Considerar \(f\colon X \to Y\) y \(B \subseteq Y\). Demostrar que \(f\bigl( f^{-1}(B) \bigr) \subseteq B\). Dar un ejemplo donde la inclusión sea estricta.
  \item Decidir si la sucesión \(b_{n+1} = \sqrt{2 + \sqrt{b_n}}\), con \(b_1 = 1\), es convergente. En caso afirmativo, determinar su límite.
  \item Considerar \({(x_n)}\) una sucesión acotada tal que \(\liminf x_n = \limsup x_n = L\). Demostrar que \(x_n \to L\).
  \item Decidir si \(\sum \dfrac{n! \, e^n}{n^{n+1}}\) es convergente o divergente.
  \item Considerar \({(a_n)}\) una sucesión decreciente tal que \(\sum a_n\) es convergente. Demostrar que \(n a_n \to 0\).
  \item Considerar \(a_1 \geq a_2 \geq \cdots \geq 0\). Demostrar que \(\sum a_n\) converge si y sólo si \(\sum 2^n a_{2^n}\) converge.
\end{enumerate}

\clearpage

\section{Segundo parcial, primera fecha, 20/12}

\begin{enumerate}[label=(\arabic*)]
  \item Sea \(f\colon \R \to \R\) y \(f_n(x) = f(nx)\) para \(n \in \N\). Si \({(f_n)}\) es equicontinua en \(0\), demostrar que \(f\) es constante.
  \item Considerar \(F\) cerrado y \(K\) compacto en \(\R^n\) tales que \(F \cap K = \varnothing\). Demostrar que \(d(F, K) > 0\), donde \(d(F, K) = \inf\{ d(f, k) : f \in F, \, k \in K \}\).

        Dar un ejemplo donde \(F\) y \(G\) sean cerrados disjuntos tales que \(d(F, G) = 0\).
  \item Considerar \(f\colon X \subseteq \R^n \to \R\) una función tal que para todo \(\e > 0\) existe \(g_{\e}\colon X \to \R\) continua tal que \(|f(x) - g_{\e}(x)| < \e\) para todo \(x \in X\). Mostrar que \(f\) es continua.
  \item \begin{enumerate}[label=(\alph*)]
          \item Considerar \(f_n\colon X \to Y\) tales que \(f_n \rightrightarrows f\) y \(g\colon Y \to Z\) uniformemente continua. Demostrar que \(g \circ f_n \rightrightarrows g \circ f\).
          \item Considerar \(g_n\colon Y \to Z\) tales que \(g_n \rightrightarrows g\) y \(f\colon X \to Y\). ¿Es cierto que \(g_n \circ f \rightrightarrows g \circ f\)?
        \end{enumerate}
  \item \begin{enumerate}[label=(\alph*)]
          \item Demostrar que \(x^n (1 - x)\) converge uniformemente a \(0\) en \([0, 1]\).
          \item Estudiar la convergencia puntual y uniforme de las siguientes series en el intervalo \([0, 1]\).
                \begin{enumerate}[label=(\roman*)]
                  \item \(\sum_{n \geq 1} x^n (1 - x)\)
                  \item \(\sum_{n \geq 1} {(-1)}^n x^n (1 - x)\)
                \end{enumerate}
        \end{enumerate}
\end{enumerate}

\clearpage

\section{Segundo parcial, segunda fecha, 10/02/2025}

\begin{enumerate}[label=(\arabic*)]
  \item Sea \(f\colon \R^n \to \R^m\). Mostrar que las siguientes afirmaciones son equivalentes.
        \begin{enumerate}[label=(\alph*)]
          \item \(f\) es continua.
          \item Para todo \(E \subseteq \R^n\) se cumple que \(f(\cl{E}) \subseteq \cl{f(E)}\).
        \end{enumerate}
  \item Considerar \(f_n\colon [0, 1] \to \R\) dadas por \(f_n(x) = \dfrac{n x^2}{1 + n x^2}\). Estudiar, por definición, la convergencia puntual y uniforme de \({(f_n)}\) y \({(f_n')}\).
  \item Considerar \(f_n\colon A \subseteq \R^n \to \R^m\) una sucesión de funciones uniformemente continuas que converge uniformemente a una función \(f\colon A \to \R^m\). Demostrar que \(f\) es uniformemente continua.
  \item Considerar, para \(n \geq 2\), las funciones \(f_n(x) = \dfrac{x^n}{\ln n}\) definidas en \([0, 1]\).
        \begin{enumerate}[label=(\alph*)]
          \item Mostrar que \({(f_n)}\) converge uniformemente.
          \item Mostrar que \({(f_n)}\) es equicontinua en \([0, 1]\).
        \end{enumerate}
  \item Considerar \(A \subseteq \R^2\) conexo y no acotado. Demostrar que para todo \(a \in A\) y \(r > 0\) el conjunto \(\{ x \in A : d(x, a) = r \}\) es no vacío.
\end{enumerate}

\clearpage

\section{Segundo parcial, tercera fecha, 21/02/2025}

\begin{enumerate}[label=(\arabic*)]
  \item Sean \(f_n\colon E \subseteq \R^n \to \R\) continuas. Si \(f_n \rightrightarrows f\), entonces para toda sucesión \({(x_n)} \subseteq E\) y todo \(x \in E\) tales que \(x_n \to x\) se tiene que \(f_n(x_n) \to f(x)\).
  \item Considerar \(f_n \rightrightarrows f\) en \([a, +\infty)\). Suponer que para cada \(n \in \N\) vale \(\lim_{x \to +\infty} f_n(x) = a_n\) con \(a_n \in \R\). Probar que existe \(\lim_{n \to +\infty} a_n\). ¿Vale lo mismo si se cambia \(f_n \rightrightarrows f\) por \(f_n \to f\) puntualmente?
  \item Considerar \(A\) y \(B\) conjuntos no vacíos en \(\R\). Mostrar que si \(A\) es compacto y \(B\) es cerrado, entonces \(A + B\) es cerrado.
  \item Demostrar que la clausura de un conjunto conexo es conexa.
  \item Decidir si las siguientes sucesiones de funciones son equicontinuas en los dominios indicados.
        \begin{enumerate}[label=(\alph*)]
          \item \(f_n(x) = n\) en \(\R\).
          \item \(g_n(x) = x + n\) en \(\R\).
          \item \(h_n(x) = n x\) en \(\R\).
        \end{enumerate}
  \item Considerar \({(f_n)}\) y \({(g_n)}\) dos sucesiones de funciones equicontinuas en un conjunto \(X\). Mostrar que la sucesión \({(f_n + g_n)}\) es equicontinua en \(X\). ¿Es cierto para la sucesión de los productos \({(f_n g_n)}\)?
  \item Estudiar la convergencia uniforme de las funciones \(f_n(x) = \dfrac{2 n x^2}{1 + n^2 x^4}\) en \([0, 1]\) y en \([1, +\infty)\).
\end{enumerate}


\cleardoublepage
\phantomsection
\addcontentsline{toc}{chapter}{\bibname}
\bibliographystyle{unsrt}
\bibliography{references}
\nocite{*}

\end{document}
